%% Template for the submission to:
%%   The Annals of Statistics [AOS]
%%
%%%%%%%%%%%%%%%%%%%%%%%%%%%%%%%%%%%%%%%%%%%%%%
%% In this template, the places where you   %%
%% need to fill in your information are     %%
%% indicated by '???'.                      %%
%%                                          %%
%% Please do not use \input{...} to include %%
%% other tex files. Submit your LaTeX       %%
%% manuscript as one .tex document.         %%
%%%%%%%%%%%%%%%%%%%%%%%%%%%%%%%%%%%%%%%%%%%%%%

\documentclass[aos]{imsart}

%% Packages
\RequirePackage{amsthm,amsmath,amsfonts,amssymb}
\RequirePackage[numbers]{natbib}
%\RequirePackage[authoryear]{natbib}%% uncomment this for author-year citations
%\RequirePackage[colorlinks,citecolor=blue,urlcolor=blue]{hyperref}%% uncomment this for coloring bibliography citations and linked URLs
%\RequirePackage{graphicx}%% uncomment this for including figures

\startlocaldefs
%%%%%%%%%%%%%%%%%%%%%%%%%%%%%%%%%%%%%%%%%%%%%%
%%                                          %%
%% Uncomment next line to change            %%
%% the type of equation numbering           %%
%%                                          %%
%%%%%%%%%%%%%%%%%%%%%%%%%%%%%%%%%%%%%%%%%%%%%%
%\numberwithin{equation}{section}
%%%%%%%%%%%%%%%%%%%%%%%%%%%%%%%%%%%%%%%%%%%%%%
%%                                          %%
%% For Axiom, Claim, Corollary, Hypothesis, %%
%% Lemma, Theorem, Proposition              %%
%% use \theoremstyle{plain}                 %%
%%                                          %%
%%%%%%%%%%%%%%%%%%%%%%%%%%%%%%%%%%%%%%%%%%%%%%
\theoremstyle{plain}
%\newtheorem{???}{???}
%\newtheorem*{???}{???}
%\newtheorem{???}{???}[???]
%\newtheorem{???}[???]{???}
\newtheorem{thm}{\textsc{Theorem}}[section]
\newtheorem{prop}{\textsc{Proposition}}[section]
\newtheorem{cor}{\textsc{Corollary}}[section]
\newtheorem{lem}{\textsc{Lemma}}[section]
%%%%%%%%%%%%%%%%%%%%%%%%%%%%%%%%%%%%%%%%%%%%%%
%%                                          %%
%% For Assumption, Definition, Example,     %%
%% Notation, Property, Remark, Fact         %%
%% use \theoremstyle{remark}                %%
%%                                          %%
%%%%%%%%%%%%%%%%%%%%%%%%%%%%%%%%%%%%%%%%%%%%%%
%\theoremstyle{remark}
%\newtheorem{???}{???}
%\newtheorem*{???}{???}
%\newtheorem{???}{???}[???]
%\newtheorem{???}[???]{???}
%%%%%%%%%%%%%%%%%%%%%%%%%%%%%%%%%%%%%%%%%%%%%%
%% Please put your definitions here:        %%
%%%%%%%%%%%%%%%%%%%%%%%%%%%%%%%%%%%%%%%%%%%%%%
\newcommand{\david}[1]{{\small{\color{blue} \sf $\heartsuit$ David: #1}}}
\newcommand{\piotr}[1]{{\small{\color{red} \sf $\clubsuit$ Piotr: #1}}}
\newcommand{\paul}[1]{{\small{\color{green} \sf $\spadesuit$ Paul: #1}}}
\newcommand{\E}{\mathbb{E}}
\newcommand{\N}{\mathbb{N}}
\newcommand{\R}{\mathbb{R}}
\newcommand{\I}{\mathbb{I}}
\newcommand{\M}{\mathcal{M}}
\newcommand{\T}{\mathcal{T}}
\newcommand{\bs}{\boldsymbol}
\endlocaldefs

\begin{document}

\begin{frontmatter}
%%%%%%%%%%%%%%%%%%%%%%%%%%%%%%%%%%%%%%%%%%%%%%
%%                                          %%
%% Enter the title of your article here     %%
%%                                          %%
%%%%%%%%%%%%%%%%%%%%%%%%%%%%%%%%%%%%%%%%%%%%%%
\title{???}
%\title{A sample article title with some additional note\thanksref{T1}}
\runtitle{???}
%\thankstext{T1}{A sample of additional note to the title.}

\begin{aug}
%%%%%%%%%%%%%%%%%%%%%%%%%%%%%%%%%%%%%%%%%%%%%%%
%% Only one address is permitted per author. %%
%% Only division, organization and e-mail is %%
%% included in the address.                  %%
%% Additional information can be included in %%
%% the Acknowledgments section if necessary. %%
%% ORCID can be inserted by command:         %%
%% \orcid{0000-0000-0000-0000}               %%
%%%%%%%%%%%%%%%%%%%%%%%%%%%%%%%%%%%%%%%%%%%%%%%
\author[A]{\fnms{???}~\snm{???}\ead[label=e1]{???@???}},
\author[B]{\fnms{???}~\snm{???}\ead[label=e2]{???@???}}
\and
\author[B]{\fnms{???}~\snm{???}\ead[label=e3]{???@???}}
%%%%%%%%%%%%%%%%%%%%%%%%%%%%%%%%%%%%%%%%%%%%%%
%% Addresses                                %%
%%%%%%%%%%%%%%%%%%%%%%%%%%%%%%%%%%%%%%%%%%%%%%
\address[A]{???\printead[presep={,\ }]{e1}}

\address[B]{???\printead[presep={,\ }]{e2,e3}}
\end{aug}

\begin{abstract}
Lorem ipsum dolor sit amet, consectetur adipiscing elit. Sed auctor enim metus, vel commodo neque laoreet ut. Praesent ac maximus mauris. Aliquam at lorem vitae nisi porta viverra. Integer dictum at ante at lobortis. Nulla placerat tortor eros, non condimentum elit molestie a. Aliquam orci est, congue vel hendrerit sed, scelerisque in ante. Duis mattis id urna nec tempus. Pellentesque efficitur auctor ex a egestas. Cras euismod tincidunt nisi eget egestas. Vivamus purus augue, semper vitae iaculis eu, egestas sit amet odio. Interdum et malesuada fames ac ante ipsum primis in faucibus.

Pellentesque nisi nunc, commodo id orci ac, imperdiet ullamcorper ante. Donec iaculis consequat interdum. Vivamus at efficitur est. Phasellus vitae nisi ac ipsum pellentesque faucibus in quis odio. Nunc vehicula nulla ac egestas posuere. Morbi sit amet neque feugiat, semper neque quis, finibus quam. Sed eget hendrerit mi. Nam fermentum tempus vulputate. Pellentesque convallis risus tellus, eget viverra elit egestas et. Ut felis turpis, pharetra ac convallis quis, facilisis quis urna. Quisque quis ipsum sit amet odio imperdiet blandit ut non ex.

Vestibulum rhoncus pretium ex, egestas facilisis neque efficitur quis. Morbi malesuada volutpat erat, sed aliquet lectus dictum ac. Aenean ligula massa, consequat a blandit id, placerat.
\end{abstract}

\begin{keyword}[class=MSC]
\kwd[Primary ]{???}
\kwd{???}
\kwd[; secondary ]{???}
\end{keyword}

\begin{keyword}
\kwd{???}
\kwd{???}
\end{keyword}

\end{frontmatter}


%%%%%%%%%%%%%%%%%%%%%%%%%%%%%%%%%%%%%%%%%%%%%%
%% Please use \tableofcontents for articles %%
%% with 50 pages and more                   %%
%%%%%%%%%%%%%%%%%%%%%%%%%%%%%%%%%%%%%%%%%%%%%%
%\tableofcontents

%%%%%%%%%%%%%%%%%%%%%%%%%%%%%%%%%%%%%%%%%%%%%%
%%%% Main text entry area:



\section{Introduction} 

\david{I re-wrote the opening paragraph a bit. We referred to estimation when we don't obtain any estimation results, etc.} \paul{I feel like the first few lines flow less naturally now. I don't see the link between the first sentence and the next one.} \paul{"Estimation" does not seem wrong to me. Aren't we estimating the support of $\beta^*$. Maybe "inference" instead of "estimation"?}

High-dimensional techniques designed for settings where there are many parameters relative to the sample size $n$ frequently rely on strict assumptions on sparsity and signal strength.\\
\\
%In many modern applications, the sample size is insufficient to reliably estimate the parameters in a model. This limitation has motivated the development of high-dimensional techniques that frequently rely on strict assumptions on sparsity and signal strength. %Yet, often, external information is available and can be leveraged to enhance parameter estimation. 
\paul{My proposal as opening paragraph, basically the previous one but with inference instead of estimation:}
In many modern applications, the sample size is insufficient to reliably do inference. This limitation has motivated the development of high-dimensional techniques that frequently rely on strict assumptions on sparsity and signal strength. Yet, often, external information is available and can be leveraged to enhance inference.
%Our motivation stems from the observation that, often, 
In particular, external information can guide our decisions regarding which parameters should be included in a given model. In regression, this information may originate from the inherent nature of each variable (e.g., clinical history versus genomic markers, sociodemographics versus job history) or external data that distinguishes relevant variables from irrelevant ones (e.g., functional annotations on genes). 
Our paper investigates this concept in the the Gaussian Sequence Model (GSM) and linear regression, where external information helps partition variables into groups with potentially distinct characteristics.

Incorporating external information is often advocated within the framework of data integration, as joint learning can lead to more accurate inference compared to analyzing datasets separately. Indeed, numerous applied works have employed external information to guide inference. For instance, \cite{stingo:2011, quintana, cassese:2014} proposed Bayesian variable selection methods for gene expression, where prior probabilities for non-zero coefficients depend on biological knowledge and meta-covariates. Additionally, \cite{chen_tinghuei:2021} predicted disease outcomes by allowing LASSO penalties to depend on functional annotation categories. Beyond regression, node and edge covariates have been incorporated in \cite{peterson} and \cite{jewson2023graphical} to drive edge inclusion in Gaussian graphical models, while \cite{schiavon2022generalized} used meta-covariates to determine non-zero loadings in factor models. In causal analysis, covariates of interest (referred to as treatments) guide the inclusion of other covariates based on their degree of association with the treatments \cite{belloni2014inference,antonelli2021bayesian}. Collectively, empirical evidence consistently demonstrates improved structural learning when external data is suitably integrated. However, despite this empirical success, a theoretical framework explaining precisely why this occurs remains elusive.

Here we show that incorporating external information allows pushing the theoretical limits under which consistent variable selection is possible, understood as recovering the true model with probability going to 1 as $n$ and the number of variables $p$ grow. Previous literature has extensively explored how high-dimensional variable selection is constrained by inherent characteristics of the dataset, such as $n$, $p$, the signal strength, the correlation between variables and the number of variables truly associated to the outcome \cite{buhlmannvandegeerbook,overviewvarsel,handbookbayesvarsel,wainwrighthighdimbook}. In this paper we analyse how conditions for consistency can be relaxed in the presence of external information. We focus on variable selection but using external information is also expected to benefit parameter estimation and prediction problems. 

To make our ideas concrete, consider the variable selection problem in the linear model with Gaussian noise:
\begin{equation}\label{eq:linearmodel}
	\bs y \;=\; \bs X \bs \beta^{\star} + \bs\epsilon,\qquad \bs \epsilon \sim N(0,\sigma^2 I_n),
\end{equation}
where $\bs X\in \R^{n\times p}$, $\bs y\in \R^n$, and $\sigma >0$. A large class of methods operate by penalizing the size of an estimated $\hat{\bs  \beta}$. For instance, penalized likelihood methods are based on optimizing the log-likelihood plus a penalty term driven by the $\ell_q$ "norm" of $\hat{\bs  \beta}$ for $q\in [0,1]$ or folded concave penalties \citep{lasso,overviewvarsel,bertsimas}. 
In Bayesian settings, variable selection is often based on posterior model probabilities that are directly connected to $\ell_0$ penalties \citep{BIC,EBIC,rossell2022concentration}. 
%In frequentist settings, recent advances in discrete optimization have also rendered $\ell_0$ penalization much more computationally efficient \cite{bertsimas,mazumdergroup}.
% Roughly speaking, the penalty must grow with the number of truly inactive variables ($\beta_j^\star=0$) to avoid selecting an overfitting model, but it must also be small enough to avoid missing small signals \cite{buhlmannvandegeerbook,handbookbayesvarsel,wainwrighthighdimbook}.
% Those requirements results in so-called betamin conditions that involve the smallest non-zero absolute parameter value $\beta_{\min}^{\star}= \min_j |\beta_j^{\star}|$.  

We study the case where the external information partitions the $p$ variables into $b$ blocks denoted $B_j$, $j=1,\ldots,b$ across which key characteristics such as the level of sparsity or signal strength may vary. We introduce block $\ell_0$ penalties that allow modulating the strength of the penalty in each block. Unlike standard $\ell_0$ penalties, like BIC \cite{BIC}, RIC \cite{RIC} or EBIC \cite{EBIC}, block penalties are non-exchangeable in the sense that the complexity penalty for adding a variable $j$ may depend on its block $B_j$.

\paul{Reorganized the intro to have this contribution paragraph.}
Our main contributions are as follows. We show that block-based non-exchangeable $\ell_0$ penalties soften
%the use of non-exchangeable $\ell_0$ penalties driven by external information softens
the theoretical conditions for consistent model recovery in the Gaussian sequence model and in linear regression. %high-dimensional linear models. 
Specifically, %We further demonstrate 
that it is possible for an oracle to set block-based penalties so that variable selection consistency is either attained where otherwise not possible or at a faster rate. Another important contribution is that we introduce data-based variable selection procedures, %with informed block penalties 
based on information criteria and empirical Bayes, that realize the theoretical benefits without requiring an oracle. Other more minor contributions include a tight necessary condition on signal strength in high-dimensional orthogonal settings and consistency results for the EBIC in a wider class of asymptotic regimes than proved in \cite{EBIC} and \cite{ebicdiverging}.

Our work has connections with ideas from multiple hypothesis testing. In this line of research, \cite{genovese2006} proposed a false discovery rate (FDR) %adjustment 
procedure in which p-values are weighted based on prior information. They show that if the weights are positively associated to the null hypotheses being false, their procedure improves power. Subsequent works discussed oracle choices of weights either based on external information or derived from the data %with weighting based on external information and weights derived from the data 
\cite{roederwasserman,basu}. 
Adaptive FDR methods \cite{blanchard} estimate the %Another stream of FDR control methods are called adaptive \cite{blanchard} and use estimate of the
proportion of truly null hypotheses to improve statistical power.
Recently, \cite{ramdas2019prior} combined external information %weighting 
and adaptivity in a group-based FDR control. 
Relative to this work, we study %focus on studying 
variable selection consistency in high-dimensional regression %a high-dimensional regression setting 
as well as the associated conditions on sparsity and signal strength.

\david{I added a paragraph on Scarlett et al.}
\color{blue} Also related to our work, \cite{scarlett2012compressed} studied compressed sensing when prior information allows splitting parameters into blocks, where one knows the true proportion of non-zero parameters in each block. They show that sparse signal recovery requires smaller $n$ than with exchangeable penalties. Key differences with our work are their assuming independence across covariates, which renders the results inapplicable to regression. Also, they do not consider the sequence model, nor that the proportions of non-zero parameters are unknown in practice.   \color{black} \paul{I would also add that the design is assumed random, and consistency is taken with respect to this source of randomness.}

The paper is organized as follows. Section~\ref{sec:intro_blockpen} introduces block $\ell_0$ penalization. Section~\ref{sec:properties_orth} presents its %the theoretical 
variable selection properties %of block penalization 
in the Gaussian sequence model, %, a high-dimensional linear and orthogonal model that generalizes orthogonal linear regression. 
in an oracle setting where the sparsity and betamin conditions are known for all groups.
%We also discuss the benefits of an oracle block penalization in relation to the discriminative power of the external information. 
Those results apply to essentially any penalized likelihood or Bayesian method, including $\ell_0$ and $\ell_1$ penalties, because in that setting selection operates by thresholding \cite{davidomiros,rigollethighdimensional}. Section~\ref{sec:properties_legression} studies block $\ell_0$ penalties in linear regression, and shows analogous benefits to those in Section~\ref{sec:properties_orth}.
%analyzes the properties of block $\ell_0$ penalties in the more general linear regression setting where variables are correlated, and shows that the benefits of the penalties are analogous to those discussed in Section~\ref{sec:properties_orth}.
%We remark that results in this section
These results can be extended to a wide class of Bayesian variable selection methods. Section~\ref{sec:informed_l0} presents data-based %variable selection 
procedures %with block $\ell_0$ penalties that leverage external information to 
that achieve faster convergence without requiring an oracle. Section~\ref{sec:simulation} shows empirical examples and Section~\ref{sec:discussion} concludes. Proofs are gathered at the end of the paper.
\medskip

\noindent \textbf{Notation:}  
We denote by $\bs \beta\in \R^p$ the parameters of interest and by $\bs \beta^*$ their true values. Let  $V=\{1,\ldots,p\}$. 
%In a linear model parameterized by $\bs \beta\in \R^p$ with true value $\bs \beta^*$, we denote by $V=\{1,\ldots,p\}$ the set of indices of $\bs \beta$. 
For any $A\subseteq V$ and any vector $\bs x\in \R^p$, $\bs x_A$ denotes the subvector of $\bs x$ with entries corresponding to indices in $A$. For any matrix $\bs X \in \R^{n\times p}$, $\bs X_A$ denotes the submatrix of $\bs X$ obtained by selecting the \emph{columns} with indices in $A$. We denote by
$S=\left\{i\in V:\,\,  \beta^{\star}_i \neq 0 \right\}$ the true support of $\bs \beta^{\star}$, its size is $s=|S|$, and hence $p-s$ is the number of truly inactive parameters.
Denote $\beta_{\min}^{\star}=\min_{i \in S}|\beta^{\star}_i|$ the smallest true signal.
We assume that $V$ is partitioned into a fixed number $b$ of disjoint blocks $B_j \subseteq V$ for $j=1,\ldots,b$.
We denote by $S_j=\left\{i\in B_j:\,\,  \beta^{\star}_i \neq 0\right\}$ the set of active parameters in block $B_j$,  $s_j=|S_j|$ its size, $p_j-s_j$ the number of inactive parameters in the block,
and $\beta_{\min,j}^{\star}=\min_{i \in S_j}|\beta^{\star}_i|$.
Given sequences $f(n) > 0$ and $g(n) >0$, $f(n)=O(g(n))$ means that there exists a constant $c<\infty$ such that $f(n) \leq c g(n)$ for all $n \geq n_0$ and some fixed $n_0$.
$f(n)=o(g(n))$ means that $f(n) / g(n) \to 0$, and $f(n)=\Theta(g(n))$ means that $f(n) = O(g(n))$ and $g(n) = O(f(n))$. 
\medskip

\noindent\textbf{Our asymptotic regime:} Although not denoted explicitly, $s_j$, $p_j$ and $\beta_{\min,j}^{\star}$ are  functions of $n$, and so are $s$, $p$ and $\beta_{\min}^{\star}$. We study asymptotic regimes where
%in this paper high-dimensional asymptotic results and assume:
\begin{enumerate}[start=1,label={(A\arabic*)}]
    \item $n \,\to \, \infty$, and the number of blocks $b$ is constant.
    \item For all $j$, $p_j-s_j \,\to \, \infty$.
\end{enumerate}
Assumption A2 implies that $p-s$ goes to infinity but we do not make any assumption on the asymptotic regime linking $n$ and $p$. We do not impose any restriction either on the total numbers of active parameters, $s$, nor across blocks, the $s_j$'s. 
\david{But surely we will end up needing some conditions involving the $s_j$'s, at least in regression. And on how $p$ can grow with $n$. To be cautious, I wouldn't write "we do not impose any restriction...", it may raise eyebrows and possibly disappoint later if we end up setting conditions (even if indirectly) on $(p,n)$ or $(s,s_j)$. For example, you could just say "Assumption A2 is a mild assumption requiring that one has a growing number of inactive covariates in each block."}\paul{add $s>1$ or all $s_j>1$}
%The number of blocks $b$ is considered constant. \david{Maybe we should list as an a}
% \begin{equation*}
%     p-s \,\to \, \infty, \;\;\;\;s\to \, \infty.
% \end{equation*}

% Note that the family of regimes considered includes all regimes such that $s_j \to \infty$ and $p_j/s_j \to +\infty$, and allows some regimes in which $p_j/s_j \to 1$.

\section{Variable selection via informed block penalization}\label{sec:intro_blockpen}

Consider a set of candidates models $\M$ given by subsets $M\subseteq V$ and their corresponding coordinate subspaces
$$
\mathcal{L}_M\;:=\;\{\beta\in \R^p: \beta_i=0\mbox{ if }i\notin M\}.
$$
By definition the true model $S$ is the smallest subset such that $\bs\beta^\star\in L_S$, which for simplicity we assume to be unique \piotr{Two points: (i) S was defined in the notation section as the support of $\beta^*$ and its easiest to just to repeat this definition, (ii) the way it is defined, it is unique so we do not need to assume this.} \paul{I think $S$ is unique as defined too. But David explained in an email on Tue, 1 Oct, 14:27 that S can be non-unique. Based on David's explanation, what can be non-unique is "S" defined as "smallest subset of covariates minimizing ssr" but this is not how we defined it here.}\piotr{Then we need to sort out, which definition we use. Because if $S$ is defined as the support of $\beta^*$ (as here) then this is uniquely defined. Perhaps you means that $\widehat S$ may not be unique?}.
Consider a standard $\ell_0$ selection procedure with penalty $\eta(M)$ that depends on $M$ linearly 
%Consider the class of variable selection procedures based on standard linear $\ell_0$-penalties $\eta$, that is, penalties that depend on the model $M$ linearly 
through its cardinality, $\eta(M)=\kappa |M|$ for some $\kappa>0$. The selected model is
%The selected model is the smallest maximizer, over all $M\in \M$, of the penalized likelihood: 
\begin{equation}\label{eq:stdl0}
\hat{S}= \arg\max_{M\in \M} \left\{\max_{\beta\in \mathcal{L}_M}\ell(\bs y; \bs \beta)\;-\;\kappa |M|\right\},
\end{equation}
where $\ell(\bs y; \bs \beta)$ is the log-likelihood function. 
%More generally, consider a fixed partition $B_1/\cdots/B_b$ of $V$.
%, that is, a collection of disjoint subsets $B_j\subseteq V$ whose union is equal to $V$. We fix this partition throughout and $b$ refers to the number of its blocks 
%We consider a procedure that selects a model $M$ for $\bs\beta^\star$ using different penalties for different blocks $B_1,\ldots,B_b$, defining a partition of $V$. 
We study a block-informed procedure that penalizes differently each block $B_1,\ldots,B_b$,
%Fixing $(\kappa_1,\ldots,\kappa_b)$, the selected model $\hat S^b$ is the smallest maximizer, over all $M\in \M$, of the  likelihood penalized with \textbf{block informed} linear $\ell_0$-penalty:
\begin{equation}\label{eq:Mhat}
  \hat S^b\;\in\;\arg\max_{M\in \M} \left\{\max_{\beta\in \mathcal{L}_M}\ell(\bs y; \bs \beta)-\sum_{j=1}^b\kappa_j |M_j| \right\},
\end{equation}
where $\kappa_1>0,\ldots,\kappa_b>0$, we denote a model by $M=M_1\cup\cdots\cup M_b$, and $M_j\subseteq B_j$ represents the  selected parameters in block $j$.
Note that $\hat{S}^1$ is the standard $\ell_0$ selector in \eqref{eq:stdl0}.
\piotr{Section 2 is half-page. Section 3 takes 10 pages. Is there a way to reorganize sectioning so that things are lighter and more balanced? (this is not critical)}
\section{Gaussian sequence model}\label{sec:properties_orth}

We now introduce the Gaussian sequence model, its connection to orthogonal linear regression and discuss that $\hat S^b$ performs block-wise hard thresholding. We also provide its variable selection properties, discuss its benefits in terms of lessening the conditions for consistency variable selection and convergence rates, and illustrate said advantages under different asymptotic regimes.
%In this section, we start by introducing the Gaussian sequence model, its connections to orthogonal linear regression and show that in that setting $\hat S^b$ is a particular instance of a more general class of procedures that perform blockwise hard thresholding. 
%We proceed with a derivation of the variable selection properties of block thresholding in the Gaussian sequence model. We then discuss the benefits they bring in terms of conditions for selection consistency and rates of convergence. %We particularly highlight the impact of the quality of external information and the resulting partition on those benefits. 
%We end with examples exhibiting the benefits under different asymptotic regimes and scenarios of external information.

In this section, in addition to A1-A2, we assume that
%The results and discussion of this section focus on asymptotic regimes in which, in addition to A1 and A2, the following assumption on the numbers of active parameters holds:
\begin{enumerate}[start=3,label={(A\arabic*)}]
    \item $\text{for all } j, \;\; s_j \,\to \, \infty$
\end{enumerate}
Assumption A3 is made to simplify the exposition, and because we derive a new necessary condition on signal strength under A3, but we obtained analogous results when $s_j$'s can be finite (Supplement XXXX).
%Assumption A3 implies that the total number of active parameters $s$ goes to infinity. 
%Similar results hold in the case where for all $j$ $s_j$ is finite. We assume A3 holds to simplify the exposition and because we derive a new necessary condition on signal strength in that setting. Results in the finite case are discussed in the Appendix~\ref{sec:appendixGSMsfinite}.

\subsection{Gaussian sequence model and thresholding}\label{sec:GSMgen}

The Gaussian sequence model assumes:
\begin{equation}\label{eq:GSM}
	\bs y \;=\; \sqrt{n} \bs \beta^{\star} + \bs\epsilon,\qquad \bs \epsilon \sim N(0,\sigma^2 I_p),
\end{equation}
where $\bs y \in \mathbb{R}^p$ and without loss of generality we set $\sigma^2=1$ to simplify notation. It is a popular simplified model for high-dimensional inference, and is also central to non-parametric statistics. See Chapter 3 in \cite{johnstone} for a complete introduction. 

Let $\tilde{\bs \beta}=\frac{1}{\sqrt{n}}\bs y \sim N(\bs \beta^*,\tfrac{1}{n}I_p)$ be the MLE under \eqref{eq:GSM}, and $\tilde{\bs \beta}_M$ the MLE under $M$ (with $j^{th}$ entry equal to $\tilde{\beta}_j$ if $j \in M$ and 0 otherwise).

\david{Lemma \ref{lem:useful-reformulation-GSM} can go into the supplement without harm, to save space and ease the paper's flow.}\piotr{I agree}
\begin{lem}\label{lem:useful-reformulation-GSM}
In the Gaussian sequence model \eqref{eq:GSM}, $\hat{S}^b$ defined in \eqref{eq:Mhat} satisfies:
    \begin{equation}\label{eq:MhatGSM}
  \hat S^b\;\in\;\arg\max_{M\in \M} \bigg\{\tfrac{1}{2}\| \sqrt{n}\tilde{\bs \beta}_M\|^2-\sum_{j=1}^b\kappa_j |M_j|\bigg\}, \qquad M=M_1\cup \ldots\cup M_b.	
\end{equation}
\end{lem}

The sequence model generalizes orthogonal linear regression with normalized columns, i.e. when $\bs X$ in \eqref{eq:linearmodel} satisfies $\bs X^\top \bs X=n I_p$.
%variable selection in the Gaussian sequence model generalizes selection in orthogonal linear regression with normalized columns, that is when $\bs X$ in \eqref{eq:linearmodel} satisfies $\bs X^\top \bs X=n I_p$. 
Then, $\tilde{\bs \beta}^{\prime}=\frac{1}{n}\bs X^{\top}\bs y \sim N(\bs \beta^*,\tfrac{1}{n}I_p)$, and selection with $\hat{S}^b$ operates as in the sequence model.\piotr{What is $\tilde{\bs\beta}'$? Also, for the record, I still cannot understand why you say that the GSM generalizes orthogonal regression not the other way around.} \paul{Some of the explanation has been commented out, maybe we need to add some back. $\tilde{\bs\beta}'$ is the MLE under the "full model" $V$ in orth linear regression. $\tilde{\bs \beta}$ is the MLE under the "full model" $V$ in GSM.} \paul{I think GSM generalizes orthogonal linear regression only from the the point of view of selection/selection with penalized methods? I think we said it before but GSM allows $n<p$ which is not possible with orth. linear regression. That's not enough for you?}\piotr{In GSM $n=p$ really, you just introduced a parameter, which you call $n$. But this is not a big point for me. David said that everybody in the literature agrees that GSM is more general and I see no point in fightting over this :) Please add back everything needed to understand the notation. }
%In that case, the MLE under the full model $V$ is $\tilde{\bs \beta}^{\prime}=\frac{1}{n}\bs X^{\top}\bs y$, it is distributed $N(\bs \beta^*,\tfrac{1}{n}I_p)$ and for any $M\in\M$, the MLE $\tilde{\bs \beta}^{(M)}$ is $\tilde{\bs\beta}_M$. 
%It follows that $\hat{S}^b$ satisfies too the characterisation in \eqref{eq:MhatGSM} replacing $\tilde{\bs \beta}$ by $\tilde{\bs \beta}^{\prime}$. 
%As a result, selection with $\hat{S}^b$ operates in the same way in the two settings and the discussions and results presented in this Section apply to both. 
%Orthogonal linear regression is however a more restrictive model since, unlike the Gaussian sequence model, it requires $n\geq p$.

% We first consider the case where the columns of $\bs X$ are orthogonal and normalized: $\bs X^\top \bs X=n I_p$. This case covers orthogonal designs in the linear regression \eqref{eq:linearmodel} and necessarily $p\leq n$. It also covers 
% In that model, the design matrix can be taken to be $\bs X=\sqrt{n}I_p$ and the number of samples is $p$. The quantity $n$ controls the shrinking rate of the MLE variance and does not constrain $p$. The Gaussian sequence model \piotr{Citing a book always requires a more precise citation (section number etc)}\paul{It is an unpublished book, but I cited it as a book and added the chapter number}. 

% Denote $\tilde{\bs \beta}=\frac{1}{n}\bs X^{\top}\bs y$. 

The next proposition shows that the block-wise $\ell_0$ penalty reduces to hard-thresholding $\tilde{\bs \beta}$, with block-dependent thresholds.
%From the characterization of Lemma~\ref{lem:useful-reformulation-GSM}, we can show that linear $\ell_0$-penalty approaches are exactly equivalent to hard thresholding $\tilde{\bs\beta}$.
\begin{prop}\label{prop:isthresh}
	In the Gaussian sequence model \eqref{eq:GSM}, $\hat{S}^b$ defined in \eqref{eq:Mhat} satisfies $\hat S^b=\hat S^b_1\cup \cdots \cup \hat S^b_b$, where, for each $j=1,\ldots,b$,
	$$
	\hat S_j^b\;:=\;\left\{i\in B_j:\; |\tilde{\beta}_i|> \sqrt{\frac{2\kappa_j}{n}}\right\}.
	$$ 
	% Each $\hat S^b_j$ is defined uniquely unless $|\tilde{\beta}_i|= \sqrt{\frac{2\kappa_j}{n}}$ for some $j=1,\ldots,b$ and $i\in B_j$.
 % \paul{I removed the discussion on whether $\hat S^b_j$ is defined uniquely. Because I define it as the smallest minimizer now. Also now $\hat S_j^b\;:=\;\left\{i\in B_j:\; |\tilde{\beta}_i|> \sqrt{\frac{2\kappa_j}{n}}\right\}$ and not $\hat S_j^b\;:=\;\left\{i\in B_j:\; |\tilde{\beta}_i|\geq \sqrt{\frac{2\kappa_j}{n}}\right\}$ anymore. It makes it more consistent with the $\ell_1$ selectors that have $>$ instead of $\geq$ in Lemma~\ref{lem:l1isthreshold} }
\end{prop}

%From Proposition~\ref{prop:isthresh} it follows that with standard linear $\ell_0$-penalties there is a single threshold while block informed linear $\ell_0$-penalty allow for different thresholds across blocks. 
Thus, equivalently, we study generic hard thresholding model selectors of the form
\begin{equation}\label{eq:thresholdS}
\hat S^{T,b}\;:=\;\{i\in B_j: |\tilde{\beta}_i|>\tau_j\}, \qquad\bs \tau=(\tau_1,\ldots,\tau_b)\in \R^b_{>0}.	
\end{equation}
\paul{I feel like we should highlight more this definition. Because we say in the introduction our object of study is $\hat S^b$ and then have an entire section on $\hat S^{T,b}$. Someone not reading very carefully could easily get lost. Maybe as we have special formats for Lemmas, Prop and theorem. We could have one for Definition to be used for selectors $\hat S^b$ and $\hat S^{T,b}$.}
%$\hat S^{T,b}$ generalizes standard hard thresholding. Taking $b=1$ gives the standard hard threshold selector. 
\paul{I think we could drop the 1 when we refer to the standard hard threshold selector $\hat S^{T,1}$ to make the notation lighter and more consistent with other notations in the discussion of the benefits.}
\david{Agreed. I added $\hat{S}$ to the lhs in \eqref{eq:stdl0} so that it's introduced upfront. Moreover, I would just use $\hat{S}^b$ instead of $\hat{S}^{T,b}$ to simplify notation, since $\hat{S}^{b}=\hat{S}^{T,b}$ in the sequence model. And then just $\hat{S}$ for the standard selector.}

Most popular variable selection procedures also take this block-based thresholding form, including for the LASSO and adaptive LASSO (Lemma \ref{lem:l1isthreshold}), Bayesian procedures under most standard priors \citep{davidomiros} where one sets a different prior inclusion probability in each block.

\david{I would end the section here, and move Lemma \ref{lem:l1isthreshold} to the supplement.}\piotr{I agree. Also, I do not know what happened to this lemma. Now the link to adaptive LASSO is not given explicitly and it is hard to decipher what is the meaning of the second part of this result.} \paul{Something that we lost in the trimming of text (a while ago) is the following argument that seemed important to me. When we say we do non-exchangeable penalties Adaptive LASSO comes to mind naturally. Initially we explained that adalasso with MLE / LASSO weights does not help with selection because it is eventually equivalent to having a single threshold.}\piotr{Perhaps, as David suggested, move this lemma to the appendix and add back all the necessary discussion there. }

%Results analogous to Proposition~\ref{prop:isthresh} hold for many other variable selection procedures. We now show it for $\ell_1$-penalized procedures. 
\begin{lem}\label{lem:l1isthreshold}
	In the sequence model \eqref{eq:GSM}, selecting the non-zero entries  of the solution to 
\begin{equation}\label{eq:genl1}
	\underset{\bs \beta \in \R^p}{\rm minimize} \;\;\tfrac12\|\bs y-\sqrt{n}\bs\beta\|^2+\sum_{j=1}^b \lambda_j \sum_{i \in B_j}|\beta_{i}|
\end{equation}
for non-negative %penalty parameters 
$\lambda_1,\ldots,\lambda_b$  is equivalent to taking the selector $\hat S^{T,b}$ with $\bs \tau=(\lambda_1/n,\ldots,\lambda_b/n)$. Let $\bs\beta^\circ$ be either the MLE or the LASSO estimator with penalty $\lambda^\circ$. Selecting the non-zero entries  of the solution to the problem
\begin{equation}\label{eq:genl1ada}
	\underset{\bs \beta \in \R^p}{\rm minimize} \;\;\tfrac12\|\bs y-\sqrt{n}\bs\beta\|^2+\sum_{j=1}^b \lambda_j \sum_{i \in B_j}\tfrac{|\beta_{i}|}{|\beta_i^\circ|}.
\end{equation}
is equivalent to taking the selector $\hat S^{T,b}$ with $\bs \tau=(\sqrt{\lambda_1/n},\ldots,\sqrt{\lambda_b/n})$ if $\bs\beta^\circ$ is the MLE and with $\bs \tau=(\lambda^\circ/2n+\sqrt{(\lambda^\circ/2n)^2+\lambda_1/n},\ldots,\lambda^\circ/2n+\sqrt{(\lambda^\circ/2n)^2+\lambda_b/n})$ if $\bs\beta^\circ$ is a LASSO estimate.
\end{lem}

%Lemma~\ref{lem:l1isthreshold} introduces block informed versions of the $\ell_1$ penalties in LASSO \cite{lasso} and adaptive LASSO \cite{adalasso} that recover the original penalties when $b=1$. The lemma shows that in the sequence model the statistical properties of $\hat{S}^{T,b}$ extend to $\ell_1$ penalties too. 
%This result holds also in orthogonal linear regression replacing $ \tilde{\bs\beta}$ by $\tilde{\bs \beta}^{\prime}$.

%Although we do not discuss it in detail here, the block-thresholded selector $\hat{S}^{T,b}$ also includes Bayesian procedures. Under most standard priors Bayesian variable selection operates by hard thresholding $\tilde{\bs\beta}$ in the orthogonal case \citep{davidomiros}. One may set different prior inclusion probabilities in each block, and the highest posterior probability model results in block thresholding. 


\subsection{Selection based on block thresholds}
\label{ssec:block_thresholds}


%\paul{Added this outline of subsection } \david{great, thanks}
We now study the statistical performance of the block-based thresholding operator in \eqref{eq:thresholdS}.
% the model selector in \eqref{eq:Mhat}. By Proposition~\ref{prop:isthresh}, we equivalently study 
%model selector $\hat S^{T,b}$. By Proposition~\ref{prop:isthresh}, statistical properties of $\hat S^{T,b}$ apply to our selector of interest $\hat S^{b}$ in the Gaussian sequence model. 
Block thresholding was previously studied in the context of parameter estimation for wavelet-based models with equally sized blocks; see Chapters 7 to 9 in \cite{johnstone}. % for a complete discussion. 
We focus here on its variable selection properties. We successively analyze properties relative to false positives, false negatives, exact recovery and rates of convergence.  
%The relevant quantities are $p_j=|B_j|$ and $s_j=|S_j|$ for $j=1,\ldots,b$. 
\bigskip

\noindent\textbf{False positives.} 
%\noindent\textbf{False positives (guarantees for $\hat S^{T,b}\subseteq S$).} 
A false positive %in a block threshold selector 
occurs when there exists $j=1,\ldots,b$ and $i\in B_j$ such that $\beta^\star_i=0$ %(equiv. $i\in B_j\setminus S_j$) 
but $|\tilde \beta_i|>\tau_j$. In other words, $\hat S^{T,b}\not\subseteq S$. A sufficient condition to guarantee that asymptotically $\hat S^{T,b}\subseteq S$ is:
\begin{enumerate}[start=4,label={(A\arabic*)}]
    \item For all $j=1,\ldots,b$ and all sufficiently large $n$
\begin{equation*}\label{cond:suffIallblockthresh} 
\tau_j \;\geq\;  \sqrt{\frac{{2\ln(p_j-s_j)}}{{n}}},	
\end{equation*}
\end{enumerate}
The next proposition shows that %in addition to being sufficient 
A4 is also essentially necessary.
\piotr{edited the proposition}
\begin{prop}\label{prop:suffnecIblockthresh}In the %Gaussian
sequence model \eqref{eq:GSM}, assume A1 and A2.
\begin{enumerate}
    \item [a)] If A4 holds then % $\hat{S}^{T,b}$ satisfies:
$P(\hat{S}^{T,b}\subseteq S) \to 1$.
\item [b)] If for some $j\in\{1,\ldots,b\}$,
\begin{equation}\label{cond:necIallblockthresh0}
\lim_{n\to \infty} \frac{\tau_j}{\sqrt{\frac{{2\ln(p_j-s_j)}}{{n}}}} \;<\;1\; ,	\end{equation}
then  $P(\hat{S}^{T,b}\subseteq S ) \not \to 1$.
\end{enumerate}
\end{prop}
Proposition~\ref{prop:suffnecIblockthresh} shows that %assumption~A4 is  sufficient and that 
a necessary condition to have, asymptotically, no false positives %and model recovery 
is: 
\begin{equation}\label{cond:necIallblockthresh}
\lim_{n\to \infty} \frac{\tau_j}{\sqrt{\tfrac{2\ln(p_j-s_j)}{n}}}\;\geq\;1\qquad\mbox{for all }j=1,\ldots,b.	
\end{equation}
Proposition~\ref{prop:suffnecIblockthresh} extends to $\hat{S}^{T, b}$ results known for $\hat{S}^{T, 1}$ in orthogonal linear regression \cite{ slope, buhlmannvandegeerbook, wainwrighthighdimbook} and in the Gaussian sequence model in \cite{ariascastro} and Chapter 3 of \cite{johnstone}.
\bigskip

\noindent\textbf{False negatives.} 
%\noindent\textbf{False negatives (guarantees for $S\subseteq \hat S^{T,b}$).} 
False negatives occur when $S\not\subseteq \hat S^{T,b}$, i.e. for some $j \in \{1,\ldots,b\}$ and $i\in B_j$,  we have that $\beta^\star_i\neq 0$ and $|\tilde\beta_i|\leq \tau_j$. %This means that $\tau_j \geq \min_{i \in S_j}|y_i/\sqrt{n}|$ for some $j=1,\ldots, b$. 
Control of false negatives in a block $B_j$ requires bounding the minimum of $|S_j|=s_j$ independent Gaussians with non-zero means $\beta^{\star}_i$, $i\in S_j$. This can be achieved by setting a betamin condition %, that is a condition on the smallest non-zero parameter 
on $\beta_{\min,j}^{\star}\;:=\;\min\{|\beta^\star_i|: i\in S_j\}$,
and a condition on the thresholds $\tau_j$ for $j=1,\ldots,b$.  
Intuitively, to have no false positives (also known as screening property), each $\tau_j$ needs to be smaller than $\beta^*_{\min,j}$.
%It is clear that, to have that $S\subseteq \hat S^{T,b}$, also known as screening property, each $\tau_j$ cannot be too large; in particular $\tau_j$ needs to be smaller than $\beta^*_{\min,j}$. 
A closely-related sufficient condition is: 
\begin{enumerate}[start=5,label={(A\arabic*)}]
    \item For all $j=1,\ldots,b$ and sufficiently large $n$
\begin{equation*}\label{cond:suffIIallblockthresh}
	 \beta_{\min,j}^{\star} - \tau_j \;\geq\; \sqrt{\frac{2\ln(s_j)}{n}},
%	 \beta_{\min,j}^{\star}-\tau_j \geq \sqrt{\frac{2\ln(s_j)}{n}},
\end{equation*}
\end{enumerate}

\begin{prop}\label{prop:suffnecIIblockthresh}
 In the sequence model \eqref{eq:GSM}, assume A1-A3 and A5. Then %$\hat{S}^{T,b}$ satisfies:
%\begin{equation}\label{cond:suffIIallblockthresh}
%	\lim_{n\to \infty}\frac{\beta_{\min,j}^{\star}-\tau_j}{\sqrt{\frac{2\ln(s_j)}{n}}} \;>\;1 ,
%\end{equation}
$P(S\subseteq \hat{S}^{T,b} ) \to 1$.
\end{prop}
Condition A5 for standard thresholding was proved to be sufficient for selection consistency in orthogonal linear regression in \cite{wainwrighthighdimbook}. Proposition~\ref{prop:suffnecIIblockthresh} generalizes it to block thresholds. The study of {\it necessary} conditions %for $S\subseteq \hat S^{T,b}$ asymptotically
is more complicated. We provide below a first general condition valid for any vector $\bs \beta^*$ and a second tight condition for a worst case where $\beta_i^*=\beta_{\min}^*$ for all $i\in S$.
\piotr{edited the proposition. also in the whole paper I replaced $min$ with $\min$ for consistency}
\begin{prop}\label{prop:necIIblockthresh}In the %Gaussian
sequence model \eqref{eq:GSM}, assume A1 and A2.
\begin{enumerate}
    \item [a)] Suppose that $\tau_j<\beta^\star_{\min,j}$ satisfies \eqref{cond:necIallblockthresh} for some $j \in \{1,\ldots,b\}$. If \begin{equation}\label{cond:necIIallblockthresh2}
    \lim_{n\to \infty} \sqrt{n}(\beta_{\min,j}^{\star}-\tau_j)\; <\; \infty
\end{equation}
then $P(S\subseteq \hat{S}^{T,b} ) \not \to 1$.
\item [b)] Suppose there exists $j\in \{1,\ldots,b\}$ such that  $\beta^{\star}_i=\beta_{\min,j}^{\star}$ for all $i \in S_j$, $\tau_j<\beta^\star_{\min,j}$ satisfies \eqref{cond:necIallblockthresh}, and  $s_j/p_j<1$. If A3 holds and, in addiiton, 
\begin{equation}\label{cond:necIIallblockthresh}
\lim_{n\to \infty} 	\frac{\beta_{\min,j}^{\star}-\tau_j}{\sqrt{\frac{\pi}{2}\frac{\ln(s_j)}{n}}} \;\leq\;  1\end{equation}
then $P(S\subseteq \hat{S}^{T,b} ) \not \to 1$.
\end{enumerate}\end{prop}
\cite{abraham2023sharp} showed the necessity of \eqref{cond:necIIallblockthresh2} for consistency in the single-block case %in the sequence model 
in more restrictive asymptotic regimes satisfying $p/s \to \infty$. %Proposition~\ref{prop:necIIblockthresh}a) shows it generalizes to $b>1$ and in a larger class of regimes. 
Proposition~\ref{prop:necIIblockthresh}b) gives a new precise necessary  rate of growth of $\sqrt{n}\left(\beta_{\min}^{\star}-\tau\right)$ for a worst case in which all signals are equal. It shows that in that case sufficient assumption~A5 is essentially necessary up to a constant factor, and that A5 cannot be relaxed in general.
\bigskip

\noindent\textbf{Exact model recovery guarantees.} 
Combining Conditions A4 (Proposition~\ref{prop:suffnecIblockthresh}) and A5 (Proposition~\ref{prop:suffnecIIblockthresh}),
we asymptotically recover the correct support if for all $j=1,\ldots,b$,
%To guarantee model recovery %with high probability the block thresholds $\tau_j$ must be $\geq \sqrt{2\ln(p_j-s_j)/n}$ %greater or equal to  $\sqrt{2\ln(p_j-s_j)/n}$ 
%(assumption~A4, Proposition~\ref{prop:suffnecIblockthresh}), and $\leq {\beta^\star_{\min,2}}+\sqrt{2\ln(s_j)/n}$ (assumption~A5, Proposition~\ref{prop:suffnecIIblockthresh}). In particular, we get asymptotic guarantees for correct support recovery as long as for all $j=1,\ldots,b$, $\tau_j$ and $\beta_{\min,j}^{\star}$ satisfy both A4 and A5 that is, for all $n$ sufficiently large:
\begin{equation}\label{eq:thetamin}
  \sqrt{\frac{2\ln(p_j-s_j)}{n}}\;\;\leq\;\;\tau_j\;\;\leq\;\;\beta_{\min,j}^{\star}-\sqrt{\frac{2\ln(s_j)}{n}}
\end{equation}
for sufficiently large $n$. This implies that %A4 and A5 then imply 
that for all $j=1,\ldots,b$, %for all $n$ sufficiently large:
\begin{equation}\label{eq:betamin}
    {\beta_{\min}^{\star}}_j \geq \sqrt{\frac{2\ln(p_j-s_j)}{n}} + \sqrt{\frac{2\ln(s_j)}{n}}
\end{equation}
By Propositions~\ref{prop:suffnecIblockthresh} and~\ref{prop:suffnecIIblockthresh} signals as small as the right-hand side in \eqref{eq:betamin} can be recovered in block $B_j$. As we analyze the benefits of block-based thresholding, %based on external information, 
we will discuss signal recoverable with $\hat{S}^{T,b}$ and $\hat{S}^{T,1}$ and will refer to betamin condition \eqref{eq:betamin} as the \textbf{minimal betamin condition}. \paul{David does not like this expression? Maybe "smallest signal recoverable"? Although we show it is sharp only up to constant factor and in a particular worst case} \david{Maybe refer to it as "blockwise betamin condition? "As we analyze the benefits of block-based thresholding, we will refer to \eqref{eq:betamin} as the block-wise betamin condition".}\piotr{i like ``blockwise betamin condition''}

%Conversely for betamin smaller than the right-hand side in \eqref{eq:betamin}, there is no choice of thresholds $(\tau_1,\ldots,\tau_b)$ for which selection consistency is guaranteed. 
Propositions~\ref{prop:suffnecIblockthresh} b) and \ref{prop:necIIblockthresh} show that Conditions A4-A5 are essentially necessary. It follows that betamin condition \eqref{eq:betamin} is near-necessary for selection consistency in the following sense.
\piotr{Can we have the same numbering for everybody? Now lemmas have different numbering than propositions, which makes harder to look for results.}
\begin{lem}\label{lem:necbetaminblockthresh}
    If there exists $j\in\{1,\ldots,b\}$ such that $\beta^{\star}_i=\beta_{\min,j}^{\star}$ for all $i \in S_j$, $s_j/p_j<1$ and
    \begin{equation}\label{eq:necbetaminblockthresh}
      \lim_{n\to \infty}\frac{\beta_{\min,j}^*}{\sqrt{\tfrac{2\ln(p_j-s_j)}{n}}+\sqrt{\tfrac{\pi}{2}\tfrac{\ln(s_j)}{n}}}\;<\; 1  
    \end{equation}
then under Conditions A1-A3,\;  $P(\hat{S}^{T,b} = S)\not \to 1$
\end{lem}
Lemma~\ref{lem:necbetaminblockthresh} shows the minimal betamin condition \eqref{eq:betamin} cannot be relaxed much in general. Under the conditions of Lemma~\ref{lem:necbetaminblockthresh}, if $\beta_{\min,j}^*$ remains strictly smaller than $\sqrt{2\ln(p_j-s_j)/n}+\sqrt{\pi\ln(s_j)/2n}$, selection consistency is not possible. \paul{David suggested to delete this last sentence, I don't undertand why. It seems correct and important to me} \david{It's important but it seems repetitive, we already stated this just before the lemma.}\piotr{Better to remove the part before the lemma and move it here. Placing the conclusion in the end strenghtens the message. (but I am fine either way)}


\bigskip

\noindent\textbf{Rate of convergence.} From the %proofs of the 
preceding results, we can bound the rate of convergence of $P(\hat S^{T, b}\neq S)$.\piotr{Here we do not assume neither A1 nor A2? also I slightly edited this result so check if you agree} %We also give oracle threshold values $\tau^\star_j$ that approximately minimize the bound, and the resulting oracle rate.
\begin{thm}\label{theo:fullrecovGS}
Assume A4 and A5. Then, for sufficiently large $n$,
%	Assume that the thresholds $\tau_j$ and the $\beta^\star_{\min,j}$ satisfy A4 and A5 then for all sufficiently large $n$
\begin{equation} \label{eq:rateconvgsm}
P(\hat{S}^{T,b} \neq S)\;\leq\;\sum_{j=1}^b\,\,\frac{e^{-\tfrac{n}{2}\big[\tau_j^2-\tfrac{2\ln(p_j-s_j)}{n}\big]}}{\sqrt{\pi\ln(p_j-s_j)}}+\frac{e^{-\tfrac{n}{2}\big[(\beta^\star_{\min,j}-\tau_j)^2-\tfrac{2\ln(s_j)}{n}\big]}}{\sqrt{\pi\ln(s_j)}}.
\end{equation}
On the other hand, if $\beta^\star_{\min,j}$ satisfy \eqref{eq:betamin} for every block $j$, and the thresholds $\tau_j$ take oracle values
	\begin{equation} \label{eq:optimtauj}
\tau^\star_j \;=\; \frac{1}{2}\beta^\star_{\min,j}+\frac{1}{n\beta^\star_{\min,j}}(\ln({p_j-s_j})-\ln({s_j})),
\end{equation}
then A4 and A5 hold and, for sufficiently large $n$, %there exists $n_0$ such that for all $n\geq n_0$:
\begin{equation}\label{eq:optimrateconvgsm}
    P(\hat{S}^{T,b} \neq S) \; \leq \; 2 \sum_{j=1}^be^{-[\frac{n}{2}{\beta^\star_{\min,j}}^2-\max\{\ln(p_j-s_j),\ln(s_j)\}]}
\end{equation}
\paul{what about the $\sqrt{\pi\ln(p-s)}$ and $\sqrt{\pi\ln(s)}$}
\end{thm}
In \eqref{eq:rateconvgsm} the first component of each term in the sum decreases exponentially in $n\tau^2/2$, while the second component decreases exponentially in $n(\beta_{\min,j}^*-\tau)^2/2$. The choice $\tau^*_j$ ensures that both terms are equal and hence approximates the values minimizing \eqref{eq:rateconvgsm}.  \piotr{Good moment to explain very briefly why we call them ``oracle''.} 
%For such oracle threshold, we obtain an oracle rate in \eqref{eq:optimrateconvgsm} that is driven by the distances between ${\beta^\star_{\min,j}}\!\!\!\!\!\!\!^{2}\;\;\;$ and $\max\{\ln(p_j-s_j),\ln(s_j)\}$. 
The resulting bound in \eqref{eq:optimrateconvgsm} is tightest in the worst-case scenario where $\beta_i^*=\beta_{\min}^*$ for all $i \in S$, and approximates the fastest rate achievable in that worst case.
%That bound is valid for all vectors $\bs\beta^*$ and is thus the tightest for the worst case in which $\beta_i^*=\beta_{\min}^*$ for all $i \in S$. The oracle rate in \eqref{eq:optimrateconvgsm} approximates the fastest rate achievable in that worst case.


\subsection{Gains and losses of block thresholding}\label{sec:prosconsblock}

%\paul{added this outline}
We now examine the benefits brought by block penalization. %thresholding based on externally defined blocks. 
We start with a simple example to build intuition and then discuss %go on with discussing separately 
two types of benefit: those relative to conditions for selection consistency and those relative to rates of convergence.
\bigskip

Intuitively, block penalties allow adapting 
%Broadly speaking, block penalization provides the possibility to adapt 
to varying sparsity $s_j$  or signal strengths $\beta^\star_{\min,j}$ across blocks. 
When $p_j-s_j$ is small one may use smaller thresholds $\tau_j$ and still control false positives, and when $\beta^\star_{\min,j}$ is large one may use larger $\tau_j$ and still control false negatives.
As the blocks become more
%As the external information driving the blocks becomes more 
discriminative, that is, more distinct $p_j-s_j$ and $\beta^\star_{\min,j}$ across blocks, the larger the potential benefits.
Consider a somewhat extreme example where $n = 1$, $p = 100$, $s = 10$ and $\beta_i^\star=4$ for all $i\in S$. There are $b=2$ blocks: a large sparse block with $s_1=0$, $p_1-s_1=80$ and a smaller non-sparse block with $s_2=10$, $p_2-s_2=10$. In the left panel of Figure~\ref{fig:starter}, a standard threshold selector is used with $\tau=\sqrt{2 \ln(p-s)/ n}$, the minimal threshold such that the FWER vanishes (dashed line). The large number of truly zero parameters $p-s$ forces a high threshold, which leads to a false negative. In the right panel, block thresholds are set to $\tau_j=\sqrt{2 \ln(p_j-s_j)/ n}$.
Adjusting $\tau_2$ to the lower number of truly zero means in block 2 ($p_2-s_2$) allows recovering the true model. 
\begin{figure}[t]
\includegraphics[scale=.6]{images/starter.example.pdf}
\centering
\caption{In both panels, the 90 white dots are generated from $N(0, 1)$ and the 10 black dots are generated from $N(4, 1)$. Left panel: the dashed line is $\sqrt{2 \ln (p-s)/n}$, the minimal threshold such that FWER vanishes. Right panel: observations are grouped in two blocks. The dashed lines are the minimal block thresholds for the vanishing of the FWER} %with $\hat{S}^{T,b}$.
\david{In the figure x-axis, I would delete "observation" and "observations grouped". Then may add "global thresholding" and "block-based thresholding" on top of the figure. Also, journals usually require that each panel is uploaded as a separate pdf. Then use tabular to create the two panels, as you would with a table.}
\label{fig:starter}
\end{figure}
\bigskip

\david{I moved this paragraph that you introduced later (in the examples) to here. To me, this captures the essence of why block thresholding improves upon global thresholding. I also placed the corollary immediately afterwards, to connect it with the general principle.}

%\noindent\textbf{Thresholds for $\hat{S}^{T,1}$ and $\hat{S}^{T,b}$.}
More precisely, Conditions A4-A5 give ranges of threshold values that are essentially necessary and sufficient for asymptotic support recovery. For the standard threshold selector $\hat{S}^{T,1}$, the range for the single threshold $\tau$ is
\begin{equation}\label{eq:rangethreshex}
  \sqrt{\frac{2\ln(p-s)}{n}}\;\;\leq\;\;\tau\;\;\leq\;\;\beta^{\star}_{\min}-\sqrt{\frac{2\ln(s)}{n}}.	
\end{equation}
For a block threshold selector $\hat{S}^{T,b}$, in the $b=2$ block case the ranges for $\tau_1$ and $\tau_2$ are
\begin{equation}
    \begin{aligned}\label{eq:rangeblockthreshex}
  \sqrt{\frac{2\ln(p_1-s_1)}{n}}\;\;\leq\;\;\tau_1\;\;\leq\;\;\beta^{\star}_{\min,1}-\sqrt{\frac{2\ln(s_1)}{n}}\\
  \sqrt{\frac{2\ln(p_2-s_2)}{n}}\;\;\leq\;\;\tau_2\;\;\leq\;\;\beta^{\star}_{\min}-\sqrt{\frac{2\ln(s_2)}{n}}        
    \end{aligned}
\end{equation}

\color{blue}
If the range for $\tau_1$ does not overlap with that for $\tau_2$, 
which, for example, occurs if $\beta_{\min,1}^{\star} - \sqrt{(2/n) \ln(s_1)} < \sqrt{(2/n) \ln(p_2-s_2)}$,
then a global threshold $\tau$ cannot possibly satisfy both conditions in 
\eqref{eq:rangeblockthreshex} and consistent model recovery is not possible.
We next give a precise instance where this occurs.
\color{black}\piotr{nice!}
%The betamin condition \eqref{eq:necbetaminblockthresh} in Lemma~\ref{lem:necbetaminblockthresh} allows to describe settings where consistent model recovery is possible for block thresholding, but not for standard thresholding.
%The necessary betamin condition of Lemma~\ref{lem:necbetaminblockthresh} also allows to precisely describe settings where block selection makes possible model recovery. 
\david{Corollary 3.1 is similar to my added discussion in blue above, but if features $(s,p)$ instead of $(s_j,p_j)$. Is it a refined version? How does it connect to the previous discussion? I just wrote "a precise instance" for lack of a better way how to set up the result.}

\begin{cor}\label{cor:selconstimpossible}
Assume A1-A3, that $s/p<1$, and $\beta^{\star}_i=\beta_{\min}^{\star}$ for all $i \in S$. If
\begin{align*}
    &\lim_{n\to \infty}\frac{\beta_{\min}^*}{\sqrt{\tfrac{2\ln(p-s)}{n}}+\sqrt{\tfrac{\pi}{2}\tfrac{\ln(s)}{n}}}\;<\; 1 
\end{align*}
and for $j=1,\ldots,b$, $(\tau_j,\beta^\star_{\min,j})$ satisfy A4-A5 \piotr{Cannot we just assume A1-A5 holds in the beginning?},
then 
$P(\hat{S}^{T,1} = S)\not \to 1$ and $P(\hat{S}^{T,b} = S)\to 1$.
\end{cor}
\paul{Do I need to have proof or can we say it is immediate?} \david{I normally err on the cautious side and add a brief proof where I say "The result follows immediately from..."}\piotr{If you call it ``Corollary'' perhaps fine to skip the proof. But you can write a sentence below just to guide readers.}

\noindent\textbf{Minimal betamin conditions.} The possibility to lower thresholds in the less sparse blocks as in the example above make betamin conditions for selection consistency milder with block thresholding. To simplify the exposition and without loss of generality, assume the smallest $|\beta_i^*|$ for all $i\in S$, $\beta_{\min}^{\star}$, is located in block $b$. 
From \eqref{eq:betamin}, we obtain the minimal betamin condition for the standard threshold selector $\hat{S}^{1}$:
\david{Maybe just say betamin, for simplicity?}
\begin{equation}\label{eq:betamincon.suffthresh}
    \beta_{\min}^{\star} \;>\; \sqrt{\frac{2\ln(p-s)}{n}}+\sqrt{\frac{2\ln(s)}{n}}
\end{equation} 
For the block threshold selector $\hat{S}^{T,b}$, since $\beta_{\min}^{\star}$ is located in block $b$, condition \eqref{eq:betamin} requires:
\david{Why is it enough to look only at block $b$ ? I mean, even if $\beta_{\min}^\star$ is in block $b$, $(p_j,s_j)$ from other blocks could be such that \eqref{eq:betamincon.suffblockthresh} holds in block $b$, but not for some other block $j$ no?}\piotr{Yes, this should be clarified.}
\begin{equation}\label{eq:betamincon.suffblockthresh}
    \beta_{\min}^{\star} \;>\; \sqrt{\frac{2\ln(p_b-s_b)}{n}}+\sqrt{\frac{2\ln(s_b)}{n}}
\end{equation} 
Denote the right hand-side in inequalities above: 
\begin{equation}\label{eq:betaminbardef}
    \underline{\beta}_{\min}^{*} := \sqrt{\frac{2\ln(p-s)}{n}} + \sqrt{\frac{2\ln(s)}{n}},\qquad
    \underline{\beta}^{*}_{\min,b} := \sqrt{\frac{2\ln(p_b-s_b)}{n}} + \sqrt{\frac{2\ln(s_b)}{n}}
\end{equation} Since $p_b-s_b \leq p - s$ and $s_b \leq s$, $\underline{\beta}^{*}_{\min,b} \leq \underline{\beta}_{\min}^{*}$ and condition \eqref{eq:betamincon.suffblockthresh} is always milder than \eqref{eq:betamincon.suffthresh}. Selector $\hat{S}^{T,b}$ is then able to recover smaller signals than $\hat{S}^{T,1}$. The magnitude of the difference depends on how $p-s$ and $s$ compare to $p_b-s_b$ and $s_b$. The gain thus depends on whether the block structure induced by the external information is indicative of varying sparsity. \david{I didn't find sub-section "Minimal betamin conditions" super-informative (given what we already said earlier). This is however used later for Figure Figure~\ref{fig:ratiobetamincond}, so maybe we can introduce $\underline{\beta}_{\min}^{*}$ later, when presenting that example.}\piotr{agreed}
%The analysis extends to the blocks $B_j\neq B_b$  and in general in every block, the strictness of the requirements on the smallest signal depends on how $p-s$ and $s$ compare to $p_j-s_j$ and $s_j$.  To illustrate this, we compare the growth rates required in \eqref{eq:betamincon.suffblockthresh} and in \eqref{eq:betamincon.suffthresh} in a few examples. 

\bigskip

\noindent\textbf{Rates of convergence.}   We just saw that block thresholding can attain consistent variable selection where standard thresholding cannot. We now discuss differences in the probability of correct variable selection, when consistent selection is possible for both procedures. We assume the block thresholds %in $\hat{S}^{T,b}$ 
are set to their oracle values $\tau_j^\star$ %in \eqref{eq:optimtauj} 
and similarly the single threshold %in $\hat{S}^{T,1}$  
takes its oracle value $\tau^\star$.
Let ${OR}^b_{orth}$ be the oracle convergence rate for $\tau_j^\star$ %$\hat{S}^{T,b}$ 
in \eqref{eq:optimrateconvgsm} and ${OR}_{orth}$ that for $\tau^\star$. %$\hat{S}^{T,1}$.
Then
\begin{equation}\label{eq:ratioCR}
    \frac{{OR}^b_{orth}}{{OR}_{orth}}\;=\;\sum_{j=1}^be^{-\frac{n}{2}\left({\beta^\star_{\min,j}}^2-{\beta_{\min}^\star}^2\right)+\max\{\ln(p-s),\ln(s)\}-\max\{\ln(p_j-s_j),\ln(s_j)\}}
\end{equation}
Since $\beta_{\min,j}^\star \geq \beta_{\min}^\star$, Equation \eqref{eq:ratioCR} shows that for $n$ large enough ${OR}^b_{orth} \leq {OR}_{orth}$ for any partition in blocks, i.e. the convergence rate for $\tau_j^\star$ is never worse than for $\tau^\star$. \david{As currently written, when $\beta_{\min,j}^\star = \beta_{\min}^\star$ for all $j$, then \eqref{eq:ratioCR} doesn't vanish. Also, ${\beta^\star_{\min,j}}^2-{\beta_{\min}^\star}^2$ must be larger than the difference between the max's no? In \eqref{eq:ratioCR} should it be $-\max\{\ln(p-s),\ln(s)\}+\max\{\ln(p_j-s_j)$? I ask because the display and discussion below suggest that.}
The magnitude of the gain depends on the discriminative power of the blocks.
Consider a mildly sparse setting where $s<p-s$ and $s_j<p_j-s_j$ for every $j$. Then, assuming without loss of generality that $\beta_{\min,b}^\star=\beta_{\min}^\star$, we have
\begin{equation*}
    \frac{{OR}^b_{orth}}{{OR}_{orth}}\;=\; \frac{p_b-s_b}{p-s} \,+\,\sum_{j<b}\frac{p_j-s_j}{p-s}e^{-\frac{n}{2}\left({\beta^\star_{\min,j}}^2-{\beta_{\min}^\star}^2\right)}
\end{equation*}

There are then two sources of improvement in convergence rates: in any block $j$, a smaller number of truly zero means, that is $p_j-s_j\leq p-s$, or a larger smallest signal, $\beta^\star_{\min,j} > \beta_{\min}^\star$. Depending on the sparsity of the block, one of the two effects usually prevails. In blocks where $\beta^\star_{\min,j} \approx \beta_{\min}^\star$, improvement comes mainly from having a smaller number of truly zero means. In those blocks
one sets $\tau^\star_j < \tau^\star$, so that one may detect smaller signals. 
If %In blocks where 
$\beta^\star_{\min,j} \gg \beta_{\min}^\star$,
one sets $\tau^\star_j > \tau^\star$
and the probability of false positives is reduced.


On the contrary, there is no improvement 
(${OR}^b_{orth}={OR}_{orth}$), 
when for example 
we have a non-sparse setting where $s_j>p_j-s_j$, with
$p_j-s_j=(p-s)/b$ and $\beta^\star_{\min,j}=\beta_{\min}^\star$ for all $j$. In such case, the external information fails to identify blocks that are substantially differently sparse. 

\subsection{Examples with two blocks}\label{sec:examples}

We illustrate our results with concrete examples contemplating different sparsity regimes and informativeness of the blocks. 
We focus on a mildly sparse setting where $\ln(s)=o(\ln(p-s))$, with two blocks ($b=2$). We assume that $\beta_{\min}^{\star}$, the smallest non-zero $|\beta_i^*|$ %for all $i\in S$ 
is located in block $2$. We start by plotting the range of valid thresholds for $\hat{S}^{T,1}$ and $\hat{S}^{T,b}$ as given in conditions A4 and A5. We then go on with comparing the minimal betamin conditions and the %end with a comparison of 
oracle convergence rates.

In all scenarios we set $s=e^{\sqrt{\ln(n)}}$ and, for simplicity, $s_1=s_2=s/2$. The difference in sparsity between blocks %block 2 and $V$ 
across scenarios varies as follows:
\begin{itemize}
    \item Scenario 1: $p-s=n$ and $p_2-s_2=n/2$. Here $\bs\beta^*$ is mildly sparse ($s/(p-s)= e^{\sqrt{\ln n} - \ln(n)}$) and the blocks are non-discriminative (both have the same number of truly 0 means $p_1-s_1=p_2-s_2$). %splitting $V$ in two blocks with the same number of truly 0 means.
    \item Scenario 2: $p-s=n^2$ and $p_2-s_2=n$. Here $\bs\beta^*$ is moderately sparse ($s/(p-s)= e^{\sqrt{\ln n} - 2\ln(n)}$) and the blocks are %external information is 
    discriminative ($(p_1-s_1)/(p_2-s_2)$ is a power of $n$). %with the ratio of number of truly 0 means between the two blocks being a power of $n$.
    \item Scenario 3: $p-s=e^{n/10}$ and $p_2-s_2=n^2$. Here $\bs\beta^*$ is highly sparse ($s/(p-s)= e^{\sqrt{\ln n} - n/10}$) and the blocks are %external information is 
    discriminative ($(p_1-s_1)/(p_2-s_2)$ is a power of $n$). %with the ratio of number of truly 0 means between the two blocks being exponential in $n$.
    \item Scenario 4: $p-s=n$ and $p_2-s_2=\ln(n)$. Here $\bs\beta^*$ is mildly sparse ($s/(p-s)= e^{\sqrt{\ln n} - \ln(n)}$) and the blocks are %external information is 
    highly discriminative (block 2 is non-sparse, since $\ln(p_2-s_2)=o(\ln(s_2))$).
    %with $\bs \beta^*_{B_2}$ being non sparse in the sense that $\ln(p_2-s_2)=o(\ln(s_2))$.
\end{itemize}

\david{Serving as a guinea pig, I found it hard to follow the scenarios, so I re-wrote them a bit to make them easier to follow. Now you become the new guine pigs, if unclear then we modify further! :-) Btw, can you also indicate $p_1-s_1$ in each scenario, so that the reader doesn't have to figure it out?}


Figure \ref{fig:boundthresh} shows the upper and lower bounds for $\tau$ as in \eqref{eq:rangethreshex} and $\tau_2$ as in \eqref{eq:rangeblockthreshex} in the four scenarios assuming $\beta^*_{\min}=1/10$. Observe that, for $\tau$, the sparser $\bs \beta^*$ is, the larger the $n$ needed to make the upper bound to exceed the lower bound, i.e. for which support recovery is possible. 
The block threshold $\tau_2$ requires smaller $n$, 
%Upper bounds for $\tau_2$ exceed the lower bounds for comparatively smaller values of $n$, 
with the gap with respect to $\tau$ increasing as the difference in sparsity between blocks %$V$ and block 2 
increases. The gap is close to null in Scenario 1 when external information is non-discriminative. In Scenario 3 on the contrary, recovery with $\hat{S}^{T,1}$ is not possible in the range of values of $n$ considered while it is with $\hat{S}^{T,b}$. Scenario 4 shows that the gain with $\hat{S}^{T,b}$ can be large even in setting where $V$ is only somewhat sparse.\paul{I am not so comfortable with this discussion, because conditions A4 and A5 are asymptotic bounds. For any fixed $n$, model recovery is never "impossible" it just has very low probability.}\piotr{Can we simply early in the paper write explicitly that when we say recovery is possible/impossible, this is an asymptotic statement and you should read it more like ``nearly impossible/possible''.}\\
\begin{figure}[t]
\includegraphics[scale=1]{images/boundthresh.pdf}
\centering
\caption{Bounds on $\tau$ and $\tau_2$ as per conditions A4 and A5 in Scenarios 1 to 4}
 \label{fig:boundthresh}
\end{figure}

\david{Why is $\tau_1$ not shown in the figure? Can you please label scenarios "Scenario 1", "Scenario 2", instead of "1,2,..."? (again, this is best done having a separate pdf per panel, using tabular, and adding "Scenario 1" etc as entries in the table). More importantly, I got confused with the shaded areas. I suggest to only shade the interval when it's non-empty, i.e. when lower bound $<$ upper bound. Finally, is Scenario 3 a setting where recovery with a global threshold is theoretically not possible as $n \rightarrow \infty$? It would be good to have 1-2 such scenarios, and perhaps briefly refer to them earlier, in the discussion after (19)-(20)} \paul{Recovery is not possible asymptotically in Scenario 3, since $\sqrt{\frac{2\ln(p-s)}{n}}+\sqrt{\frac{\ln(s)}{n}} \to \sqrt{2/10} \approx 0.44$ and $\beta_{\min}=1/10$. We can clarify that. Should we have another such scenario?}


\bigskip
\noindent\textbf{Minimal betamin conditions.} To illustrate how minimal betamin conditions change between $S^{T,1}$ and $S^{T,b}$, we analyze the ratio of $\underline{\beta}^{*}_{\min,b}$ and $\underline{\beta}^{*}_{\min}$ as defined in \eqref{eq:betaminbardef}. %Given the sparsity of $\bs\beta^*_V$ and $\bs\beta^*_{B_2}$ assumed in Scenarios 1 to 3 we have:
Under Scenarios 1 to 3 we have
\begin{equation}
    \frac{\underline{\beta}^{*}_{\min,b}}{\underline{\beta}_{\min}^{*} }\;\; = \;\;\Theta\Bigg(\sqrt{\frac{\ln(p_b-s_b)}{\ln(p-s)}}\Bigg).
\end{equation}
In Scenario 4 where block 2 is non sparse, we have:
\begin{equation}
    \frac{\underline{\beta}^{*}_{\min,b}}{\underline{\beta}_{\min}^{*} }\;\; = \;\;\Theta\Bigg(\sqrt{\frac{\ln(s_b)}{\ln(p-s)}}\Bigg).
\end{equation}
Figure~\ref{fig:ratiobetamincond} plots the ratio above in our four scenarios. In Scenario 1, %Observe that in Scenario 1, 
$\underline{\beta}^{*}_{\min,b}/\underline{\beta}_{\min}^{*} $ converges to 1. With non-discriminative blocks, %external information, 
the benefits %of block penalization 
in terms of %minimal
betamin condition fade away as $n \to +\infty$. In Scenario 2, with discriminative blocks, the %external information, the minimal 
betamin condition with $\hat{S}^{T,b}$ is asymptotically twice smaller than with the standard selector $\hat{S}^{T,1}$. 
In Scenario 3 where blocks are even more discriminative, $\underline{\beta}^{*}_{\min,b}/\underline{\beta}^{*}_{\min}$ converges to 0.
%With even more discriminative external information, in Scenario 3,
%$\underline{\beta}^{*}_{\min,b}$ is asymptotically infinitely smaller than $\underline{\beta}^{*}_{\min}$. 
Scenario 4 illustrates how highly discriminative blocks %external information 
can also bring significant benefits in terms of minimal betamin condition in a regime that is very mildly sparse.
\begin{figure}[t]
\includegraphics[scale=0.7]{images/betamin.bound.ratio.pdf}
\centering
\caption{Ratio of minimal betamin bounds $\underline{\beta}^{*}_{\min,b}$ and $ \underline{\beta}_{\min}^{*} $ for Scenario 1 to 4}
\label{fig:ratiobetamincond}
\end{figure}

\bigskip
\noindent\textbf{Oracle convergence rates.} We illustrate the ratio of oracle rates ${OR}^b_{orth}/{OR}_{orth}$ in \eqref{eq:ratioCR} in Scenarios 1 to 3 assuming $\beta_{\min }^{\star}=\sqrt{\frac{8 \ln (p-s)}{n} \frac{2 \ln (\ln (n))}{3}}$. \david{Is there a reason for choosing this betamin? (before we used a fixed betamin)} \paul{If betamin is fixed the ratio drops to 0 almost immediately for sparser scenarios. It is a rate of increase that give visually good results for the 3 scenarios in a same plot.} Since the oracle rate for $\hat{S}^{T,b}$ depends not only on the $p_j-s_j$'s but also the $\beta_{\min,j}^*$'s, we consider two cases within each scenario: %of the Scenarios 1 to 3: 
one case in which signals $\beta^*_{\min,1}\approx\beta_{\min,2}^*$ \david{$\approx$ is a bit vague, can we say $\beta^*_{\min,1}=\beta_{\min,2}^*$ for simplicity?} and one case in which $\beta^*_{\min,1}=1.1 \, \beta_{\min,2}^*$. Figure~\ref{fig:ratioconvrate} shows that, when the blocks are discriminative (Scenarios 2-3), the oracle rate for $\hat S^{T,b}$ is much
%, at the oracle, the convergence rate of $\hat S^{T,b}$ is asymptotically infinitely faster 
than for $\hat S^{T,1}$. Particularly so when blocks are %The difference in rates is even larger when external information is 
discriminative not only in in terms of the number of truly 0 parameters, but also in terms of signal strength. In the non-discriminative Scenario 1 the improvement stagnates as $n$ grows. % even with larger signals in the sparser block.
\david{In the left panel, the ratio doesn't seem to stagnate. Also, can we use the same line types in Fig 4 as we did in Fig 3? (e.g. the solid line is Scenario 1 in Fig 1, Scenario 2 in Fig 2.). Btw, in scenario one would we get a ratio=1 in Fig 4? Even if obvious, maybe worth adding that line. Finally, little things: in the figure top legend I'd just write $\beta_{\min}^1=1.1\beta_{\min,2}$. And can we show normal labels for n? (100 instead of 1e+02, etc)}


\begin{figure}[t]
\includegraphics[scale=0.9]{images/comp.bound.prob.rec.optim.pdf}
\centering
\caption{Ratio of bounds on convergence rates ${OR}^b$ and ${OR}^1$ in Scenarios 1 to 3. In the left panels  smallest signals are equal across blocks while they differ in the right panel}
\label{fig:ratioconvrate}
\end{figure}

\david{Devoting 3 figures to these examples seemed a little too much. How about merging Figs 3-4 into a single figure, with 3 panels?}


\section{High-dimensional linear regression}
\label{sec:properties_legression}

In Section~\ref{sec:properties_orth} we showed that, in the sequence model, using block penalties allows consistent model recovery in settings where it is otherwise not possible (e.g. smaller signals), and improves oracle consistency rates. We now extend the results to linear regression. 
We present sufficient and necessary conditions, 
betamin conditions and rates for the probability of correct selection for the block informed $\ell_0$-penalized selector $\hat S^b$ in \eqref{eq:Mhat}. We also compare the properties of $\hat S^b$ to those of the standard $\ell_0$ selector $\hat S^1$. 
Whereas the framework is more involved than in the sequence model and the required proof techniques are different, the results remain conceptually the same. 

Here $V$ is the set of all variables and $p_j-s_j$ and $s_j$ are respectively the numbers of inactive and active variables in block $B_j$. We require Assumptions A1-A2, but the number of active variables $s_j$ may either be fixed or go to infinity. 
That is, we do not assume A3 anymore. For simplicity, we assume that $\hat{S}^b$ is unique and that one constrains attention to the set of models $\M$,
such that
$\bs X_M$ has full column rank for any $M \subset \M$, and that the data-generating true support $S$ lies in $\M$. Non-uniqueness of $\hat{S}^b$ and non-full rank models can be accommodated though, at the expense of a slightly more involved treatment.

Our analysis relies on a classical connection
between $\ell_0$ penalties and Bayesian variable selection. 
Section \ref{ssec:regression_proof_strategy} reviews this connection and the proof strategy for our results.
The technical nature of the proof precludes a detailed exposition, we instead only present the most important ideas.
In Section \ref{ssec:regression_mainresult} we state our main theorem
giving variable selection consistency of $\hat{S}^{b}$, we also give necessary conditions, betamin conditions and oracle convergence rates as we did for the sequence model. Section \ref{ssec:regression_gains} discusses the gains of block penalization in further detail. Section \ref{ssec:nobetaminhyp} gives a general convergence result for $\hat{S}^{b}$ under no betamin assumption which is used in Section \ref{sec:informed_l0}.


\subsection{Proof strategy}
\label{ssec:regression_proof_strategy}


We state a well-known
reformulation of $\hat S^b$ in linear regression \eqref{eq:linearmodel}. 
\begin{lem}\label{lem:useful-reformulation} Let $\tilde{\bs \beta}^{(M)}= (\bs X_{\!\!M}^\top \bs X_{\!\!M})^{-1}\bs X_{\!\!M}^\top\bs y \in \R^{|M|}$ be the MLE under model $M\in \M$. Then $\hat S^b\;\in\;\arg\max_{M\in \M} {\rm C}(M)$, where $M=M_1\cup \ldots\cup M_b$ and
\begin{equation}\label{eq:CM}
    {\rm C}(M)\;: =\;\tfrac{1}{2}\|\bs X_M\tilde{\bs\beta}^{(M)}\|^2-\sum_{j=1}^b\kappa_j |M_j|,
\end{equation}
\end{lem}
By Lemma~\ref{lem:useful-reformulation}, $\hat S^b$ selects %proceeds by selecting the model 
$M\in \M$ with the largest score $C(M)$, or equivalently the largest normalized score
%Equivalently, $\hat S^b$ retains the model with the largest normalized score
\begin{equation}\label{eq:normalizedscore}
    NC(M)\;:=\;\frac{e^{{\rm C}(M)}}{\underset{M' \in \M}{\sum} e^{{\rm C}(M')}}.
\end{equation}
This normalized score can be understood as a pseudo-posterior probability for model $M$ in a Bayesian variable selection framework \citep{BIC}. \paul{We used "citep" instead "cite" here. is that for a specific reason?}
Proposition 1 in \cite{rossell2022concentration}, reproduced below as Lemma \ref{lem:L0toL1convergence}, proves that the expectation of such posterior probabilities bounds the probability of an incorrect variable selection $(\hat S^b\neq S)$. 

\begin{lem}\label{lem:L0toL1convergence}
\begin{longlist}
\item 
$P(\hat{S}^{b} \neq S) \,\,\leq \,\, 2 \sum_{M \in \M\setminus \{S\}} \E\left(NC(M)\right).$
\item  For any $M,M'\in\M$, such that $M\neq M'$, 
$NC(M)\;\leq \;\big(1+e^{C(M')-C(M)}\big)^{-1}$.
\end{longlist}
\end{lem}
By Lemma~\ref{lem:L0toL1convergence}(i) the vanishing, in expectation, of the sum of the normalized scores 
guarantees selection consistency. Such vanishing is similar to the \textit{strong selection consistency} property \citep{narisettyhe2014}: the concentration of pseudo-posterior model probabilities on the truth $S$ . We prove however $L_1$ convergence instead of convergence in probability.
To bound $\E(NC(M))$, by Lemma~\ref{lem:L0toL1convergence} (ii) with $M'=S$, we may instead bound a simple function of $C(S)-C(M)$: the pairwise comparison of criteria values of $S$ and $M$. 
This is achieved by noting that $C(S)-C(M)=\tfrac12 L_{SM}+\Delta_{MS}$, where for any two models $M,T\subseteq V$, we denote

\begin{equation}\label{eq:DeltaM}
\Delta_{MT}\;:=\;\sum_{j=1}^b \kappa_j(|M_j|-|T_j|). 
\end{equation}
and for any two models $Q,M\subseteq V$, we denote 
$$L_{QM}\;:=\;\|\bs X_{Q}\tilde{\bs\beta}^{(Q)}\|^2-\|\bs X_{M}\tilde{\bs\beta}^{(M)}\|^2.$$
To lower bound $C(S)-C(M)$ we use that $L_{SM}$ can be expressed in terms of chi-squared variables.
For any over-fitted model $M \supset S$, we have $\Delta_{MS}>0$ and it is well-known that $-L_{SM} \sim \chi_{|M \setminus S|}^2$. Hence, if one sets large enough $\kappa_j$ (and thus $\Delta_{MS}$), then $NC(M)$ vanishes. 
%vanishing of $\E\left( NC(M)\right)$ can be guaranteed by requiring the $\kappa_j$'s to be large enough.
For non-overfitted  models, the idea is to take the union model $Q_S=S\cup M$, and to note that $L_{SM}=L_{Q_SM}-L_{Q_SS}$. Since $M \subset Q_S$, we may use the next lemma to bound $L_{Q_SM}$ and hence also $L_{SM}$.

\begin{lem}\label{lemma:s7}
    Let $M,Q$ be any two nested models such that $M \subseteq Q$. Then
$$
L_{QM}\;=\;\|\bs X_{Q}\tilde{\bs\beta}^{(Q)}\|^2-\|\bs X_{M}\tilde{\bs\beta}^{(M)}\|^2\;\sim\; \chi_{|Q\setminus M|}^2\left(\mu_{Q M}\right), 
$$
where $\chi_{k}^2(\mu)$ denotes the noncentral chi-squared distribution with $k$ degrees of freedom and noncentrality parameter $\mu$. The noncentrality parameter $\mu_{QM}$ is given by \begin{equation}\label{eq:muQM}
    \mu_{QM}\;\;:=\;\;\|\big(I_n-P_M\big) \bs X_{Q\setminus M} \bs \beta^{\star}_{Q\setminus M}\|^2
\end{equation}
where $P_M=\bs X_{M}\left(\bs X_{M}^{\top} \bs X_{M}\right)^{-1} \bs X_{M}^{\top}$.
\end{lem}

For example, by Lemma~\ref{lemma:s7}, for a under-overfitted $M\subset S$ we have $L_{SM}\sim \chi^2_{|S|-|M|}(\mu_{SM})$, which is centered around $|S|-|M|+\mu_{SM}$.
If the noncentrality parameter $\mu_{SM}$ is large enough, then $L_{SM}$ and $C(S)-C(M)$ are also large, and $NC(M)$ vanishes. Lemma \ref{lem:boundrho} shows that large $\mu_{Q_SM}$ can be achieved by setting a betamin condition on signal strength and an eigenvalue condition on $\bs X$ involving the following quantity.  
 \begin{equation}\label{eq:rhoX}
        \rho(\bs X)\,=\,\min_{M \in \M:\, S \not \subseteq M} \lambda_{\min}\big(\tfrac{1}{n} \boldsymbol{X}_{S \setminus M}^{\top}\left(I_n- P_M\right) \boldsymbol{X}_{S \setminus M}\big) \geq 0.
    \end{equation}
\piotr{General math style comment. Avoid defining things in lemmas, propositions, etc. Here $\mu_{Q_SM}$, $P_M$, and $\rho$ are all defined in lemmas. Ideally results are just clean and short mathematical statements. (this is just a general guideline)} \paul{I have moved it outisde the lemma}
The quantity $\rho(\bs X)$ is nonnegative and relates to how distinguishable the other models $M\in \M$ are from $S$ \citep{infotheowainwright}.  More specifically, $\tfrac{1}{n} \boldsymbol{X}_{S \setminus M}^{\top}\left(I_n- P_M\right) \boldsymbol{X}_{S \setminus M}$ is the sample covariance matrix of the residuals when regressing $\boldsymbol{X}_{S \setminus M}$ on $\boldsymbol{X}_M$.
That is, $\rho(\bs X)$ measures how well one can approximate the 
$\boldsymbol{X}_{S \setminus M}$ with $\boldsymbol{X}_M$.
In an orthonormal case where $\bs X^{\top}\bs X=n\I_p$, then $\rho(\bs X)=1$\piotr{Not $1/n$?}\paul{no, it's 1}.  

\begin{lem}\label{lem:boundrho}
 For any $M\in\M$, let $Q_S=M \cup S$. The non-centrality parameter defined in \eqref{eq:muQM} satisfies: 
    \begin{equation}\label{eq:musmbound}
    \mu_{Q_SM}\; \geq \;n\,\rho(\bs X)\,\sum_{j=1}^b|S_j\setminus M_j|\,{\beta_{\min,\,j}^{\star}}^2
\end{equation}
\end{lem}



\subsection{Selection based on block $\ell_0$ penalty}
\label{ssec:regression_mainresult}

We state two sufficient conditions for asymptotic support recovery. %with selector $\hat{S}^{b}$. 
Building on our previous discussion, the first requires the block penalties $\kappa_j$ to be large enough,  
and the second is a betamin condition: 
\begin{enumerate}[start=6,label={(A\arabic*)}]
    \item For each block $j$, there exists $b_{j}\to \infty$ (as $n\to \infty$) such that for $n$ sufficiently large, %, and for all $j=1,\ldots,b$
    \david{Since we use $j=1,\ldots,b$, using $b_j$ may lead to confusion, maybe we can change for some other letter ($r_j$ or whatever)?}\piotr{agreed!}
\begin{equation}\label{cond:suffIreg}
    \kappa_j \;=\;\ln(p_j-s_j) + b_{j}
\end{equation}\item For each block $j$, there exists $c_{j}\to +\infty$ such that for $n$ sufficiently large, %for all $j=1,\ldots,b$
    \begin{equation}\label{cond:suffIIreg}
        \sqrt{\frac{(1-\gamma) n\rho(\bs X)}{2}}{{\beta_{\min}^{\star}}_j} \,- \,\sqrt{\kappa_j} \;=\; \sqrt{\ln(s_j)} + c_{j} .
    \end{equation}
where $\gamma:=\frac12(1+\max_j \alpha_j)\in (\tfrac12,1)$, and
$\alpha_j:=\frac{\ln(p_j-s_j)}{\kappa_j}$.
\end{enumerate}

These conditions are similar to A4-A5 formulated for the sequence model.
 By Proposition~\ref{prop:isthresh}, in the orthonormal %orthogonal and normalized 
 case $\bs X^{\top}\bs X=n\I_p$, setting block penalties $\kappa_j$ %$j=1, \ldots,b$ 
 is equivalent to hard thresholding with block thresholds $\tau_j=\sqrt{2 \kappa_j/n}$.
 Conditions A4-A5 can then be directly translated into conditions on $\kappa_j$, taking $\rho(\bs X)=1$. In that case, A6 is of the same order as A4, up to a term $b_j$ that can grow at an arbitrarily slow rate. The betamin assumption~A7 is the same as A5 up to, again, a term $c_j$ increasing at arbitrarily slow rate, and the factor $(1-\gamma)^{1/2}/2$. This factor depends on how much larger the $\kappa_j$'s are than $\ln(p_j-s_j)$. It is smaller than 1 and bounded away from 0 unless $b_j= \kappa_j - \ln(p_j-s_j)= o(\ln(p_j-s_j))$ for some $j$.
\david{Earlier we said that $b_j$ can grow at an arbitrarily slow rate, but we now see that $b_j$ has to be at least of order $p_j-s_j$. How about saying: This factor depends on the choice of $b_j$ in A6. For example, by default we take $b_j=...$, in which case $\gamma$ is bounded below by ... Perhaps also relate the choice of $(b_j,c_j)$ to the oracle $\kappa_j^\star$'s discussed later in the theorem, like "an ideal choice of $(b_j,c_j)$ leads to oracle penalties $\kappa_j^\star$ given in Theorem 4.2 below" (and, if it helps, briefly explain how the oracle $(b_j,c_j)$ depend on $s_j,p_j,\beta_{\min}^\star$)} \paul{$b_j$ can grow at any rate.  For any growth rate you will discard asymptotically the overfitting models. When $b_j= o(\ln(p_j-s_j))$ for some $j$, the betamin requirement is stronger than $\Theta(\sqrt{\kappa_j} +\; \sqrt{\ln(s_j)} + c_{j})$. To see it, the betamin requirement is essentially: $\sqrt{\frac{ n\rho(\bs X)}{2}}{{\beta_{\min}^{\star}}_j}  \;=2(\max_j\frac{\ln(p_j-s_j)}{b_j}+1)(\sqrt{\kappa_j} +\; \sqrt{\ln(s_j)}) + c_{j}$. It is not very intuitive to have such change in betamin condition i know. It is linked to a technicality of the proof. I have tried hard to get rid of it but could not. \\}
%As in the Gaussian sequence model, 
\david{The rest of the paragraph appeared after Theorem 4.1, but felt more natural here, as it discussed A6-A7. I'm not sure that it's strictly needed though, it basically repeats the notion that A6-A7 are very similar to the sequence model no?} In fact, from A6-A7 we can infer the smallest signals that can be recovered with $\hat{S}^{b}$. %as per conditions A6 and A7. 
We get the minimal betamin condition:
\begin{equation}\label{betamincond:minimalreg} \mbox{for all }j=1,\ldots,b \qquad
    \lim_{n\to \infty}\frac{\sqrt{(1-\gamma) n\rho(\bs X)}{\beta_{\min}^{\star}}_j}{\sqrt{2\ln(p_j-s_j)} + \sqrt{2\ln(s_j)}} >1
\end{equation}
Note that, up to the factor $\sqrt{(1-\gamma)\rho(\bs X)}$, Expression \eqref{betamincond:minimalreg} is very similar to betamin conditions for the sequence model (Corollary \ref{cor:selconstimpossible}, Expression \eqref{eq:betamincon.suffthresh}).


Finally, we remark that in our proof we discuss that one could take a betamin condition that is slightly less strict than A7, but we present A7 here to compare it more directly to A5 in the sequence model.  

We can now state our main theorem on the strong variable selection consistency of $\hat{S}^{b}$:
\begin{thm}\label{theo:suffcondlinearmodel}

Under assumptions A1, A2, A6 and, A7, we have
$$\sum_{M \in \M\setminus \{S\}}\E\left(NC(M)\right)\to 0 \qquad \mbox{and} \qquad P(\hat{S}^b = S) \to 1$$
\end{thm}


Theorem~\ref{theo:suffcondlinearmodel} considers an additive penalty across blocks
$\sum_{j=1}^b \kappa_j|M_j|$. This additivity assumption can be easily relaxed in the proof. 
For example, the EBIC penalty is non-additive \citep{EBIC}.
In Supplement XXXX we give a consistency results for non-additive penalties, which in particular improves on known betamin conditions for the EBIC.
%The consequent result shown in Appendix \ref{sec:appendixnonlinear} includes selection with the EBIC \citep{EBIC} and improves on known betamin conditions for the criterion. 
It also shows that EBIC is consistent in a class of asymptotic regimes larger than that originally considered in \cite{EBIC}, and which includes ultra-high dimensions.

From the proof of Theorem~\ref{theo:suffcondlinearmodel}, we can bound the convergence rate of $P(\hat S^{b}\neq S)$, given in Theorem \ref{prop:fullrecovreg}. We also give oracle block penalties that approximately optimize the bound, and provide the resulting oracle rate of convergence. 
In the statement, $\delta<1$ and $r_b>1$ should be understood as being arbitrarily close to 1 \paul{check 'arbitrarily close to 1'}.
\begin{thm}\label{prop:fullrecovreg}
Assume A6-A7. Then, for all sufficiently large $n$ and any $\delta \in (0,1)$,
\begin{equation}\label{eq:rateconvreg}
 P(\hat{S}^{b} \neq S)\;\leq\; 6(2^{2b}-2b)r_b\,\sum_{j=1}^b\,\, e^{-\tfrac{\delta}{2}\big[\kappa_j-\ln(p_j-s_j)\big]} + 
    e^{-\tfrac{\delta}{2}\big[(\sqrt{\frac{(1-\gamma) n\rho(\bs X)}{2}}{{\beta_{\min}^{\star}}_j} - \sqrt{\kappa_j})^2 -\ln(s_j) \big]}
\end{equation}
where $r_b\to 1$ as $n\to +\infty$.

Moreover, suppose that the $\beta^\star_{\min,j}$ satisfy \eqref{betamincond:minimalreg} and the $\kappa_j$ are set at the oracle values:
	\begin{equation} \label{eq:oraclepenreg}
\sqrt{\kappa^\star_j\,} \;=\; \frac14\sqrt{\frac{(1-\gamma) n\rho(\bs X)}{2}}{{\beta_{\min}^{\star}}_j}+\sqrt{\frac{2}{(1-\gamma) n\rho(\bs X)}}\frac{1}{{\beta_{\min}^{\star}}_j}\big(\ln({p_j-s_j})-\ln({s_j})\big).
\end{equation}
Then for all sufficiently large $n$, A6 and A7 hold and for any $\delta \in (0,1)$: 
\begin{equation}\label{eq:oraclerateconvreg}
    P(\hat{S}^{b} \neq S) \; \leq \; 12(2^{2b}-2b)r_b \sum_{j=1}^b e^{-\tfrac{\delta}{2}\big[\sqrt{(1-\gamma) n\rho(\bs X)}{32}{{\beta_{\min}^{\star}}_j}^2-\max\{\ln(p_j-s_j),\ln(s_j)\}\big]},
\end{equation}
where $r_b\to 1$ as $n\to+\infty$.
\end{thm}


We now derive conditions on the $\kappa_j$ and $\beta_{min,j}$ that are necessary for consistent variable selection. Recall that the strength of the $\kappa_j$ drives the vanishing of the pseudo posterior probabilities of over-fitted models. The $\kappa_j$ need then to be large enough for the pseudo-posterior probabilities of even the least penalized overfitting models to go to 0. For every $j=1,\ldots,b$, denote the set of models that over-fit by only one variable from block $B_j$:
$$O_j \:=\,\{M\in \M\;|\; M=S\cup\{i\} \text{ where } i\in B_j\setminus S_j\}$$
The asymptotic recovery of $S$ implies that, $\max_{M\in O_j}NC(M)/NC(S)<1$ with probability going to 1 as $n$ grows, for every $j=1,\ldots,b$, . For every $M\in O_j$, the ratio $NC(M)/NC(S)$ grows with $L_{MS}$ which, by Lemma~\ref{lemma:s7}, is $\chi^2_1$ distributed. There exists then, for every $M\in O_j$, $Z_M \sim N(0,1)$ such that $L_{MS}=Z_M^2$. A necessary condition on $\kappa_j$ follows from the properties of the maximum of the correlated $Z_M$. Denote the smallest eigenvalue of their correlation matrix:
\begin{equation}
    \underline{\lambda}_j\;:=\;\lambda_{min}(R^j)\;\;\text{where} \;\;R^j_{k,l}=\operatorname{corr}\big(Z_{M_k},Z_{M_l}\big),\;\;\forall M_k,M_l \in O_j
\end{equation}
\david{Above it should be $\mbox{cor}(Z_m,Z_n)$ for $m,n \in O_j$ right? Also, best to change the $n$ subindex to avoid confusion with the sample size $n$}\paul{Yes. I am struggling with good notation, because we use upper case for models/sets but matrix indexes are usually lower case.}
\paul{note it is not the correlation matrix of the $\sqrt{L_{MS}}=|Z_M|$. This result is based on \cite{hartigan} that gives a lower bound for maximum of correlated Gaussians and not maximum of absolute value of correlated Gaussians. The lower bound is probably the same but proving it would require additional lemmas.}
\paul{the notation for the correlation matrix is a bit sloppy but I think rather clear.}
A lower bound on the maximum of correlated Gaussians in \cite{hartigan} yields the following result.
\begin{prop}\label{prop:nectypeIreg}
     If for some $j=1,\ldots,b$,  $\lim_{n\to\infty}\frac{\kappa_j}{\underline{\lambda}_j^2\ln(p_j-s_j)} <1$, then under assumptions A1 and A2, $$P\Big(\max_{M\in O_j}\frac{NC(M)}{NC(S)}<1\Big)\to 0 \quad\text{and}\quad P(\hat{S}^b=S) \to 0.$$
\end{prop}

\paul{Note in passing: \cite{hartigan} nicely complements the result Piotr found on independence being the "worst case" for the maximum of Gaussians.\\
Interestingly it suggests that we could set lower penalty when inactive variables are highly correlated between each other. I remember very well that someone in Toronto remarked that it was somewhat strange that sufficient conditions for penalty strength did not depend on the level of correlation between variables. I could not answer well back then!}

Proposition~\ref{prop:nectypeIreg} shows that 
a necessary condition to recover S asymptotically %and model recovery 
is: 
\begin{equation}\label{cond:necIallblockpen}
\lim_{n\to\infty}\frac{\kappa_j}{\underline{\lambda}_j^2\ln(p_j-s_j)}\;\geq\;1\qquad\mbox{for all }j=1,\ldots,b.
\end{equation}
When the inactive variables are orthogonal, that is $\bs X_{V\setminus S}^{\top}\bs X_{V\setminus S}=nI$, for every $j=1,\ldots,b$ and $M\in O_j$, the $Z_M$ are independent and $\underline{\lambda}_j=1$. Under orthogonal design, condition~\eqref{cond:necIallblockpen} then recovers condition~\eqref{cond:necIallblockthresh}.

Consider now an under-fitted model $M\subset S$. The recovery of $S$ implies that with probability going to 1, $NC(M)/NC(S)<1$. By definition of $NC$ the ratio grows with penalties $\kappa_j$ and shrinks with $L_{SM}$. which is distributed $\chi^2_{|S\setminus M|}(\mu_{SM})$ by Lemma~\ref{lemma:s7}. Small values of the ratio are attained only if $\mu_{SM}$ grows fast enough compared to the $\kappa_j$. Lemma~\ref{lem:upperboundnoncentral} shows that the growth $\mu_{SM}$ is bounded by the growth of the largest active signals in $S$ that is not in $M$.
\begin{lem}\label{lem:upperboundnoncentral}
    For any $T\subseteq S$, $$\mu_{ST} \leq n\,\bar{\lambda}\,\sum_{j=1}^b|S_j\setminus T_j|\,\max_{i\in S_j\setminus T_j}{\beta_{i}^{\star}}^2 \quad \text{where}\quad \bar{\lambda}\,:=\,\lambda_{\max}\Big(\frac1n \bs X_{S}^\top\bs X_{S}\Big)$$
\end{lem}

For a fixed set of penalties $\kappa_j$, we denote the set of active signals that are small with respect to those penalties: 
\begin{equation}\label{eq:defss}
    S_j^{S}\;:=\;\Bigg\{\beta_i^{\star} \in S_j \,\Big|\,\sqrt{n\bar{\lambda}}|\beta_{i}^{\star}|=o\big(\sqrt{\kappa_j}\big) \Bigg\}
\end{equation}
Take $M=S\setminus S^S$ where $S^S=\cup_{j=1}^b S_j^S$. If $S^S\neq \emptyset$, the missing active signals in $M$ and thus $\mu_{SM}$ are small compared to the $\kappa_j$. Using the properties of the noncental $\chi^2$ distribution, we then establish the following convergence of the probabilities of $(NC(M)/NC(S)<1)$ and thus of $(\hat{S}^{b}=S)$.

\begin{prop}\label{prop:nectypeIIreg}
    If $S^S\neq\emptyset$ and for all $j=1,\ldots,b$, $\kappa_j\to \infty$, then, under assumptions A1,
    $$P\Big(\frac{NC(S\setminus S^S)}{NC(S)}<1\Big) \to 0 \quad\text{and}\quad P(\hat{S}^b=S)\to 0$$
\end{prop}
Property~\ref{prop:nectypeIIreg} shows that for any increasing penalty, when $S$ contains signals that are too small with respect to that penalty, the smaller model $S\setminus S^S$ is favored at the expense of $S$ and recovery fails. It follows that a necessary condition for variable selection consistency is:
\begin{equation}\label{cond:necIIallblockpen}
\lim_{n\to\infty}\frac{\sqrt{n\bar{\lambda}}\beta_{\min,j}^{\star}}{\sqrt{\kappa_j}}\;>\;0\qquad\mbox{for all }j=1,\ldots,b.
\end{equation}
When active variables are orthogonal, $\bs X_S^{\top}\bs X_S = nI$, $\bar{\lambda}=1$. Under orthogonal design, the condition~\eqref{cond:necIIallblockpen} is milder and less tight than the previously derived condition~\eqref{cond:necIallblockthresh} for the sequence model. A corollary of Proposition~\ref{prop:nectypeIreg} and Proposition~\ref{prop:nectypeIIreg} is as follows. 
\begin{cor}\label{cor:necrecovreg}
    If for some $j\in\{1,\ldots,b\}$   $\lim_{n\to\infty}\frac{\sqrt{n\bar{\lambda}}\beta_{\min,j}^{\star}}{\sqrt{\underline{\lambda}_j\ln(p_j-s_j)}} = 0$, then, under assumption A1 and A2,  $P(\hat{S}^b=S)\not\to 1$
\end{cor}
Corollary~\ref{cor:necrecovreg} yields a necessary condition for consistent variables selection under A1 and A2:
\begin{equation}\label{eq:betaminnecreg}
\lim_{n\to\infty}\frac{\sqrt{n\bar{\lambda}}\beta_{\min,j}^{\star}}{\sqrt{\underline{\lambda}_j\ln(p_j-s_j)}} > 0\qquad\mbox{for all }j=1,\ldots,b.
\end{equation}
\david{In the denominator above it should be $\underline{\lambda}_j \sqrt{p_j-s_j}$ no? That's what you have in Prop 4.1.}\paul{No, $\beta_{min}$ has to be compared to $\sqrt{\kappa_j}$.}
Condition \eqref{eq:betaminnecreg} is a general necessary betamin condition that applies in correlated, orthogonal, low and high dimensional setting, for $s$ fixed or diverging. \paul{I am happy about this result because I think there was no such tight necessary condition for variable selection consistency in regression in the literature} In an orthogonal design such that $\bs X^{\top} \bs X=nI$, for all $j=1,\ldots,b$, $\underline{\lambda}_j=1$ and $\bar{\lambda}=1$. In the particular case of orthogonality and diverging $s$, the previously derived necessary condition \eqref{eq:betamin} is stricter and tighter that Condition \eqref{eq:betaminnecreg}.


\subsection{Gains and losses of block penalties}
%\subsection{Gains and losses of penalizing base on externally defined blocks}
\label{ssec:regression_gains}

%\paul{added this outline }
%We now discuss the benefits of penalizing based on externally defined blocks. 
We discuss separately the benefits in the conditions for variable selection consistency and those in the convergence rate. The main message %takeaway 
is that benefits discussed in the sequence model in Section~\ref{sec:properties_orth} extend to linear regression.
\bigskip

\david{Throughout, the notion of minimal betamin condition wasn't clear to me. Probably just me being obtuse, but this is a non-standard notion that could confuse readers. People are used to hearing about betamin conditions, and about whether they are either necessary or sufficient. Would it be possible to frame them using this standard terminology? For example, the discussion below would be clearer if one said that "these are necessary betamin conditions for standard and for block penalties", or something along those lines.}

\paul{By minimal betamin condition, I meant to answer the question: what is the smallest signal that we are guaranteed to be able to recover asymptotically. It is not a necessary condition. It a sufficient betamin condition where the penalty is set to the smallest value that guarantees no type I error asymptotically. There is a small gap between the minimal betamin condition and a necessary betamin condition.}

\david{I see, so it's a betamin condition for a particular data analysis method, one that penalizes as little as possible. We should think about an alternative naming / way to refer to this condition, to avoid confusions. We should also explain clearly why this is an important setting, i.e. why do we need to discuss "minimal betamin conditions" at all, instead of just focusing on betamin conditions for a generic penalty $\kappa$. We can chat about this.}

\noindent Assume without loss of generality that $\beta_{\min}^{\star}$ is located in block $b$.  By \eqref{betamincond:minimalreg}, the minimal betamin conditions for the block informed $\ell_0$ penalty $\hat{S}^b$ and standard $\ell_0$ penalty $\hat{S}^1$ are respectively: 
$$\lim_{n\to \infty}\frac{\sqrt{(1-\gamma) n\rho(\bs X)}\,{\beta_{\min}^{\star}}}{\sqrt{2\ln(p_b-s_b)} + \sqrt{2\ln(s_b)}} >1 \qquad \mbox{and} \qquad \lim_{n\to \infty}\frac{\sqrt{(1-\gamma) n\rho(\bs X)}\,{\beta_{\min}^{\star}}}{\sqrt{2\ln(p-s)} + \sqrt{2\ln(s)}} >1$$
The requirements on $\beta_{\min,j}^\star$ and $\beta_{\min}^\star$ again relate to
%The relevant quantities to compare the minimal requirement on the growth rate of betamin are thus again 
$\sqrt{\ln(p_b-s_b)} + \sqrt{\ln(s_b)}$ and $\sqrt{\ln(p-s)} + \sqrt{\ln(s)}$, hence the discussion in Section \ref{sec:prosconsblock} on the benefits in betamin conditions extend directly to linear regression.
%The discussion in Section \ref{sec:prosconsblock} on the benefits in terms of betamin conditions of the block informed $\ell_0$ penalty in the %Gaussian 
%sequence model then extend to linear regression. %the linear regression settings. 
In particular, there are linear regression settings where variable selection consistency is not possible with $\hat{S}^1$, but is with $\hat{S}^b$. For example, in a simplest case where $\bs X^\top \bs X=n I_p$, consistency results for the sequence model apply here
%When one assumes $\bs X$ satisfies $\bs X^\top \bs X=n I_p$, results on selection in the Gaussian sequence model apply 
and Corollary \ref{cor:selconstimpossible} outlines such instances.

Let ${OR}^b_{reg}$ be the oracle convergence rate for $\hat{S}^{b}$ in Theorem~\ref{prop:fullrecovreg}, and ${OR}_{reg}$ that for $\hat{S}^{1}$ . Then
\begin{equation*}
    \frac{{OR}^b_{reg}}{{OR}_{reg}}=(2^{2b-1}-b)\frac{r_b}{r_1}\sum_{j=1}^be^{-\frac{\delta}{2}\left[\frac{(1-\gamma) n\rho(\bs X)}{32}\left({\beta^\star_{\min,j}}^2-{\beta_{\min}^\star}^2\right)+\max\{\ln(p-s),\ln(s)\}-\max\{\ln(p_j-s_j),\ln(s_j)\}\right]}
\end{equation*}
which is essentially the same as in the sequence model %the ratio in the analysis of the Gaussian Sequence model 
in \eqref{eq:ratioCR}, up to certain %mainly constant 
factors. 
Specifically, $r_b/r_1$ and $\delta$ are both close to 1, whereas $(1-\gamma) \rho(\bs X)/32$ slows the convergence rate gains %gain in convergence rate with $\hat{S}^{b}$ 
but does not alter the essence of the ratio %monotonicity and asymptotic behaviour of the ratio, 
as long as $\rho(\bs X)$ remains bounded above 0. The factor $2^{2b-1}-b$ highlights a potential limitation: guarantees of gains with block $\ell_0$ penalties deteriorate when one considers a large number of blocks $b$.
We remark that such deterioration may be a consequence of our proof strategy, rather than an inherent limitation of block penalties. Studying cases with large $b$, including $b \to \infty$, is left as future research.
%and may vanish when $b\to \infty$.

\subsection{Convergence with no betamin assumption}\label{ssec:nobetaminhyp}

We study now the asymptotics of pseudo-posterior probabilities under no assumption on the minimal signal strength. In \eqref{eq:defss}, we defined, for a choice of penalties $\kappa_j$, the set of signals that are small with respect to those penalties. We now distinguish between large and intermediate active signals in $S \setminus S^S$. For any $j=1,\ldots,b$, $\gamma$ and $c_j$ defined as in A7, let: 
\david{I suggest introducing small, intermediate and large signals $S_j^S$, $S_j^I$ and $S_j^L$ all at once earlier on, perhaps at the start of Section 4 (right before Sec 4.1). There briefly mention the role that each of them is gonna play.}
\begin{align}\label{eq:defslsi}
 \begin{split}
    &S_j^{L}\;:=\;\Bigg\{\beta_i^{\star} \in S_j \,\Big|\,\sqrt{\frac{(1-\gamma) n\rho(\bs X)}{2}}|\beta_{i}^{\star}| \,- \,\sqrt{\kappa_j} \;=\; \sqrt{\ln(s_j)} + c_{j}\Bigg\}\\
    &S_j^{I}\;:=\;S_j \setminus \big(S_j^L \cup S_j^S\big)
 \end{split}
\end{align}

The subset $S_j^{L}$ gathers the signals in $S_j$ that large in the sense that they satisfy assumption~A7. The subset $S_j^{I}$ gathers the signals in $S_j$ that are neither large nor small. Consider the set of models that contain $S^L=\cup_{j=1}^b S^L_j$ and some or none of the signals in $S^I=\cup_{j=1}^b S^I_j$.
\begin{equation}
    \T \;:=\;\Big\{M \in \M \,\big|\, M\,=\, S^L\,\cup\,R,\; R\in \mathcal{P}\big(S^I\big)  \Big\}
\end{equation}
where for any set $A$, $ \mathcal{P}(A)$ denotes the power set of $A$.
Following a strategy similar to that of the proof of Theorem~\ref{theo:suffcondlinearmodel}, we demonstrate, assuming only sufficiently large penalties, the vanishing of the sum of pseudo-probabilities of all models but those in $\T$. The main difference with respect to the proof of Theorem~\ref{theo:suffcondlinearmodel} is that, for any $M\not\in \T$, we bound $\E(NC(M))$ by a function of $C(M)-C(T)$ for a well chosen model $T$ in $\T$. 
\begin{thm}\label{thm:convtoT}
    Assume A1, A2, A6 and $|S^S_j|=O(p_j-s_j)$ for every $j=1,\ldots,b$:
    \begin{longlist}
        \item $\sum_{M \in \M\setminus \T} \E\left(NC(M)\right)\to 0$,
        \item if $|S^I|=O(1)$, then $P(\hat{S}^b \in \T)\to 1$.
    \end{longlist}
\end{thm}
Theorem~\ref{thm:convtoT} shows that, asymptotically, with only assuming sufficiently large penalties, all inactive and small active signals are discarded. On the contrary, all large signals are retained. Intermediate signals in $S^I$ may or may not be retained. The assumption $|S^I|=O(1)$ in $(ii)$ was made for simplicity and $|S^I|$ can be allowed to increase moderately. 

\section{Data analysis with informed $\ell_0$ penalization}\label{sec:informed_l0}

\david{I made a number of edits to the section, to clarify the connection between block $\ell_0$ and the Bayesian paradigm, particularly for non-Bayesian readers. Piotr in particular please chime in, if things don't make sense.}

In Sections~\ref{sec:properties_orth} and~\ref{sec:properties_legression}, we derived oracle properties for block penalties where one sets the penalties to oracle values.
We now %In this section we 
propose two data analysis methods anchored in a Bayesian perspective %on the problem to incorporate external information in selection. %and realize benefits discussed in sections ~\ref{sec:properties_orth} and ~\ref{sec:properties_legression}.
that do not require oracle values, and adapt to the unknown sparsity in the data-generating truth.
A Bayesian approach provides a natural estimator of the number of truly active variables in each block and thus, by using empirical Bayes, a straightforward way to adapt penalties to sparsity in each block. % closer to the oracle \eqref{eq:oraclepenreg}.
\\
\\
Consider a Bayesian framework for \eqref{eq:linearmodel} where %$M$ is taken as a random vector over the space of models. More precisely let 
the model $M$ is a $p$-dimensional vector of variable inclusion indicators $m_i=I(\beta_i \neq 0)$. Consider a joint prior on parameters and models 
\begin{equation}\label{eq:bayesianframework}
    p(\boldsymbol{\beta},M \mid \boldsymbol{\theta})=p( \boldsymbol{\beta} \mid  M) p(M \mid \boldsymbol{\theta})
\end{equation}
where $p(M \mid \boldsymbol{\theta})$ is the prior probability of model $M$, which depends on hyperparameter $\bs \theta$,
and $p( \boldsymbol{\beta} \mid  M)$ is a prior on regression coefficients given the model. 

Posterior model probabilities are $p(M \mid \bs y, \bs \theta) \propto p(\bs y \mid M) p(M \mid \bs \theta)$, where $p(\bs y \mid M)$ is the so-called marginal likelihood of model $M$. The BIC approximation \citep{BIC} gives
\begin{align*}
 p( \bs y  \mid M)\approx p(\bs y \mid \tilde{\bs \beta}^{(M)})n^{-|M|/2}.
\end{align*}
\paul{It was $p( M \mid \bs y ) \approx p(\bs y \mid \hat{\bs \beta}^{(M)})n^{-|M|/2}$, I haved changed to $p( \bs y  \mid M) \approx p(\bs y \mid \tilde{\bs \beta}_M)n^{-|M|/2}$. $y$ and $M$ were switched and we are using $\tilde{\bs \beta^{(M)}}$ for the MLE. That's correct no?}
It follows that the posterior model probabilities satisfy
\begin{align*}
\ln p(M \mid\bs y) \approx \ln p(\bs y \mid \tilde{\bs \beta}^{(M)}) - \frac{1}{2}\ln(n)|M| + \ln p(M\mid\bs \theta) + \text{constant}
\end{align*}
\paul{I have added equation above so it is more explicit to a reader non familiar with l0 /Bayesian selection.}
and any block or standard $\ell_0$ penalty selector discussed in earlier sections can be written as $\hat{M}= \arg\max_M p(M \mid \bs y)$. 
In particular, suppose that variable inclusions are independent a priori, and that inclusion probabilities are constant within each block. That is,
\begin{align} \label{eq:block_modelprior}
p( M \mid \boldsymbol{\theta})=
        \prod_{i=1}^p \operatorname{Bern}\left(m_i ; \theta_i\right) I(M \in \M) \mbox{, where } \forall \, i\in B_j, \: \theta_i=\theta^{(j)}
\end{align}
then $\ln p(M \mid \bs \theta)$ defines the additive block penalties studied in Section \ref{sec:properties_legression} with:
\begin{equation}\label{eq:blockpenpriorinclu}
    \kappa_j=\tfrac{1}{2}\ln(n)+\ln\big({\theta^{(j)}}^{-1}-1\big).
\end{equation}

This connection motivates an estimator of the number of truly active variables in block $B_j$, $s_j$, for any block $\ell_0$ procedure:
\begin{equation}\label{eq:shat}
    \widehat{s}_j \,:=\, \sum_{i\in B_j}\,\,\sum_{M\in \M | i \in M}NC(M)
\end{equation}
\paul{I think it is better to define $\widehat{s}_j$ from normalized scores directly. Because we see it as a Bayesian object only for the connection to Empirical Bayes. If we keep it as it was, that is, defining it as sum of posterior inclusion prop with a paragraph explaining that we have the same notation for a sum of normalized criterion values. That paragraph would not very visible so the reader could get lost.}
The normalized score $NC(M)$ can be seen as the posterior probability for model $M$ in Bayesian framework \eqref{eq:bayesianframework}. Replacing normalized scores  by posterior model probabilities in \eqref{eq:shat} gives:
\begin{equation*}
    \sum_{i\in B_j}\,\,\sum_{M\in \M | i \in M}P(M\mid \bs y) = \sum_{i\in B_j}P(\beta_i \neq 0 \mid \bs y) = E\Big(s_j \mid \bs y\Big)
\end{equation*}
From a Bayesian perspective, $\widehat{s}_j$ is then the sum of posterior inclusion probabilities and the posterior mean of $s_j$. The estimator of the proportion of truly active variables $\widehat{s}_j/p_j$ can also be formally justified as an empirical Bayes procedure. The empirical Bayes estimator of $\bs \theta$ (cf \eqref{eq:block_modelprior}) maximizes the marginal likelihood of $\bs y$ given $\bs \theta$ and is defined as:
\begin{align}
\hat{\bs \theta}^{EB}\,:=\, \arg\max_{\bs \theta} p(\bs y \mid \bs \theta)=
\arg\max_{\bs \theta} \int p(\bs y \mid \bs \beta, M) d P(\bs \beta, M \mid \bs \theta).
\nonumber
\end{align}
 In the Bayesian framework \eqref{eq:block_modelprior}, $\widehat{s}_j/p_j$ corresponds to the empirical Bayes estimate of $\bs \theta$.
\begin{lem}\label{lem:empiricalbayesinterpret} For $p(M \mid \bs \theta)$ in \eqref{eq:block_modelprior}, the empirical Bayes estimator satisfies $$\hat{\theta}^{EB}_i = \frac{1}{p_j}\sum_{i\in B_j}P(\beta_i \neq 0 \mid \bs y)$$
\end{lem}

Estimator $\widehat{s}_j/p_j$ has also attractive frequentist properties under mild assumptions. Theorem~\ref{thm:convtoT} yields the following asymptotic bounds on its expectation.
\begin{prop}
    \label{lem:shatconsistence}
Let $S^S_j$ be the set of small signals in block $j$ defined in \eqref{eq:defss}, and $S^L_j$ the larger signals in \eqref{eq:defslsi}.
Assume A1, A2, A6 and $|S^S_j|=O(p_j-s_j)$ for every $j=1,\ldots,b$. We have:$$ \,\, \frac{|S^L_j|}{p_j}\leq \lim_{n\to\infty}\E\bigg(\frac{\widehat{s}_j}{p_j}\bigg)\leq \frac{s_j-|S^S_j|}{p_j}\qquad\text{for all }j=1,\ldots,b$$
\end{prop}
\david{I reminded readers what's $S_j^S$ and $S_j^L$.}
An immediate consequence is that, if A7 also holds, then $\widehat{s}_j/p_j \stackrel{L_1}{\longrightarrow} s_j/p_j$ %is a consistent estimator of $s_j/p_j$ 
because $S_j=S_j^L$ and $S_j^S=\emptyset$ for all $j$. The proposition also shows that, when A7 is not met, $\widehat{s}_j/p_j$ is asymptotically downward biased by at least $|S^S_j|/p_j$ but larger than the proportion of signals satisfying A7. 

The two selection methods that we propose rely on $\widehat{s}_j$ to adapt the $\ell_0$ penalty in each block. The first method builds on the connection in \eqref{eq:blockpenpriorinclu} and follows the empirical Bayes approach of Lemma~\ref{lem:empiricalbayesinterpret} to set block penalties.\\ \\
%\david{On the discussion surrounding Method 1. (a) Saying that we run a BVS method sounds a bit vague. A natural strategy would be to use EBIC pseudo-posterior probabilities in Step 1, as a reasonable default that's coherent with our using EBIC to approximate posterior model probabilities. We can add a mention that one could alternatively consider other BVS methods, if so desired. But then, in Step 2 the coherent thing would be to take $\hat{M}= \arg\max p(M \mid \bs y)$, using the same priors that in Step 1, except that one uses $\hat{\theta}_i^{EB}$. That is, we either use $\ell_0$ in both steps, or a Bayesian framework in both steps. (b) To motivate Method 2, we say that Method 1 requires variable selection consistency and $s=o(n)$. It's not clear how Method 2 avoids this, e.g. in the discussion after Method 2 we say that it asymptotically uses $\kappa_j= \ln(p_j - s_j) + 0.5 \ln(n)$, and that this requires variable selection consistency in Step 1. Perhaps we should motivate Method 2 differently, as an information criterion that's driven by our theoretical analysis, and which we observed to compare positively to the empirical Bayes strategy (Method 1) in some simulation settings.}\\ \\

\david{To avoid repetition and facilitate discussion I mixed Methods 1-2 into a single algorithm. I also edited a bit the subsequent discussion, and suggested names for $\hat{S}^{EB}$ and $\hat{S}^{IB}$. Paul, please edit over anything that you don't like! :-)}

\textbf{Algorithm 1 - Empirical Bayes}
%\textbf{Method 1 - Empirical Bayes:}
\begin{enumerate}
    \item Set %For $\ell_0$ penalty 
    $\kappa_j=\kappa^{\circ}=\ln(p)+\frac{1}{2}\ln(n)$ for $j=1,\ldots,b$. Compute $\widehat{s}_j/p_j$ in \eqref{eq:shat} for $j=1,\ldots,b$.

    \item Obtain $\hat{S}^{EB}$ solving \eqref{eq:Mhat} with $\kappa_j^{EB}= \ln\bigg(p_j/\widehat{s}_j -1\bigg)+ \tfrac{1}{2}\ln(n)$.
    Alternatively, obtain $\hat{S}^{IB}$ solving \eqref{eq:Mhat} with
    $\kappa_j^{IB} = \ln(p_j-\widehat{s}_j)+ \tfrac{1}{2}\ln(n)$.
%    Select the model:
%    \begin{equation*}
%    \hat{S}^{EB}:= \arg\max_{M\in \M} \left\{\max_{\beta\in \mathcal{L}_M}\ell(\bs y; \bs \beta)-\sum_{j=1}^b\kappa_j |M_j| \right\}
%    \text{where} \,\,\forall j\,\,\kappa_j = \ln\bigg(\frac{p_j}{\widehat{s}_j}-1\bigg)+ \tfrac{1}{2}\ln(n)
%    \end{equation*}
\end{enumerate}

We refer to $\hat{S}^{EB}$ as the empirical Bayes block selector, and to $\hat{S}^{IB}$ as the information block selector.
Step 1 approximates posterior model probabilities under prior inclusion probabilities $\theta^{(j)}=(p+1)^{-1}$ in \eqref{eq:block_modelprior} that are equal across blocks, and estimates the proportion of truly active coefficients with $\hat{s}_j/p_j$. Step 2 selects a model using the block penalties $\kappa_j^{EB}$ induced by $\theta^{(j)}=\hat{s}_j/p_j$, %matching the empirical Bayes choice for prior inclusion probabilities in the equivalent Bayesian framework \eqref{eq:block_modelprior}. 
or alternatively setting $\kappa_j^{IB}= \kappa_j^{EB} + \ln(\hat{s}_j)$.
Under mild assumptions both $\kappa_j^{EB}$ and $\kappa_j^{IB}$ lead to variable selection consistency.
As shown below, the main difference is that in non-sparse settings where $s_j$ is fairly large, the penalty $\kappa_j^{EB}$ can be insufficient and there $\kappa_j^{IB}$ might be preferred. %(intuitively, $p_j/\hat{s}_j$ can be too small).

To show the consistency of $\hat{S}^{EB}$ we rely on Theorem~\ref{theo:suffcondlinearmodel}. The main challenge is showing that the $\kappa_j^{EB}$'s set in Step 2 satisfy Assumptions A6-A7. By Proposition~\ref{lem:shatconsistence}, $\widehat{s}_j/p_j$ in Step 1 is asymptotically downward biased and the bias depends on the number of signals that are not large with respect to the initial penalty $\kappa^{\circ}$. For every $j=1,\ldots,b$, denote:
\begin{align*}
    S_j^{S}(\kappa^{\circ})\;&:=\;\Bigg\{\beta_i^{\star} \in S_j \,\Big|\,\sqrt{n\bar{\lambda}}|\beta_{i}^{\star}|=o\big(\sqrt{\kappa^{\circ}}\big) \Bigg\} \\
    S_j^{L}(\kappa^{\circ})\;&:=\;\Bigg\{\beta_i^{\star} \in S_j \,\Big|\,\sqrt{\frac{(1-\gamma)n\rho(\bs X)}{2}}{\beta_{i}^{\star}} \,- \,\sqrt{\kappa^{\circ}} \;=\; \sqrt{\ln(s_j)} + c_{j}\Bigg\}
\end{align*}
\david{Maybe we could denote $S_j^S(\kappa)$ from the start, when we first introduce $S_j^S$ (same for $S^I$, $S^L$). To avoid repetition and save some space.}
where $\gamma:=\frac12(1+\max_j \frac{\ln(p_j-s_j)}{\kappa^{\circ}})$, and
some $c_j\to+\infty$. $S_j^{S}(\kappa^{\circ})$ and $S_j^{L}(\kappa^{\circ})$ are simply $S_j^{S}$ and $S_j^{L}$ as defined in \eqref{eq:defss} and \eqref{eq:defslsi} for the particular choice of penalty $\kappa^{\circ}$ in step 1. Consider assumptions:
\begin{enumerate}[start=8,label={(A\arabic*)}]
    \item $\sqrt{n}/s \to +\infty$  \david{We can just say $s=o(\sqrt{n})$ no? (more familiar to readers)}\\
\item For each block $j$, there exists $c_{j}\to +\infty$ \david{Is this the same $c_j$ as when defining $S_j^L$? If not, denote $c_j'$ to avoid confusions?} such that for $n$ sufficiently large, %for all $j=1,\ldots,b$
    \begin{equation*}\label{cond:suffIIreg2}
        \sqrt{\frac{ (1-\psi)n\rho(\bs X)}{2}}{{\beta_{\min}^{\star}}_j} \,- \,\sqrt{\ln\big(\frac{p_j}{|S^L_j(\kappa^{\circ})|}-1\big)+\frac12\ln(n)} \;=\; \sqrt{\ln(s_j)} + c_{j} .
    \end{equation*}
where $\psi=\frac12\big(1+\max_j \frac{\ln(p_j-s_j)}{\ln(p_j-s_j)+\ln(\sqrt{n}s_j^{-1}\ln^{-1}(\sqrt{n}s^{-1})}\big)$.
\end{enumerate}
\david{You're missing a parenthesis somewhere in the denominator of $\psi$ (not sure where).}

We require A8 so that $\kappa_j^{EB}$ do not shrink too fast when the $s_j$ grow and then can meet assumption~A6. \david{This first sentence doesn't seem particularly relevant for the main text, perhaps just mention in the proof.} Assumption A9 guarantees that $\beta_{min,j}$ and $\kappa_j^{EB}$ satisfy A7. A9 takes the same form as A7 where, in the targeted block penalty $\ln\big(\frac{p_j}{s_j}-1\big)+\tfrac 12\ln(n)$, $s_j$ is replaced by $|S^L_j(\kappa^{\circ})|$. That difference arises from the downward bias in $\widehat{s}_j/p_j$.
\david{I couldn't follow the discussion on A9. All I could see is that in A7 I see $\kappa_j$, which here seems to be replaced by other stuff.}
\david{Since $|S_j^L(\kappa^\circ)| \leq s_j$, is A9 to be interpreted as a slightly more stringent condition than A7? More critically, we should explain how these assumptions improve upon those for the uninformed penalties.} 
\begin{thm}\label{theo:empbayesselconsist}
    Assume A1, A2, A8, A9, $|S^S_j|=O(p_j-s_j)$ and $s_j / (p_j-s_j)<+\infty$ \david{Isn't this just saying that $s_j < p_j$?} for every $j=1,\ldots,b$. Then
$P(\hat{S}^{EB} = S) \to 1$.
\end{thm}


%In non-sparse settings, A8 may not hold and the variable selection consistency of the empirical Bayes method is not guaranteed. 
%To address this, we introduce a second method where Step 2 is inspired by the sufficient conditions of Theorem~\ref{theo:suffcondlinearmodel}.  \\ \\
%\textbf{Method 2 - Informed block $\ell_0$ criterion}
%\begin{enumerate}
%    \item Set %For $\ell_0$ penalty 
%    $\kappa_j=\kappa^{\circ}=\ln(p)+\frac{1}{2}\ln(n)$ for $j=1,\ldots,b$. Compute $\widehat{s}_j/p_j$ in \eqref{eq:shat} for $j=1,\ldots,b$.
%    
%    \item Select the model:
%    \begin{equation*}
%    \hat{S}^{IB\ell_0}:= \arg\max_{M\in \M} \left\{\max_{\beta\in \mathcal{L}_M}\ell(\bs y; \bs \beta)-\sum_{j=1}^b\kappa_j |M_j| \right\}
%    \text{where} \,\,\forall j\,\,\kappa_j = \ln(p_j-\widehat{s}_j)+ \tfrac{1}{2}\ln(n)
%    \end{equation*}  
%\end{enumerate}

For the consistency of $\hat{S}^{IB}$, we no longer require A8, and A9 is replaced by A10 below.
\david{Since $\ln(p_j - |S_j^L|)= \ln(|S_j^L|) + \ln(p_j/|S_j^L| -1)$, it seems that A10 is a bit more stringent than A9? If so, mention that A10 is a slightly stronger betamin condition than A9. Perhaps also mention whether $\psi$ in A9 is slightly larger/smaller than $\psi$ in A10.}
%In that second proposal, the step 2 penalties do not shrink with $\hat{s}_j$ as fast as in the empirical Bayes method and A8 is then not required for variable selection consistency. Consider assumptions:
\begin{enumerate}[start=10,label={(A\arabic*)}]
\item For each block $j$, there exists $c_{j}\to +\infty$ such that for $n$ sufficiently large, %for all $j=1,\ldots,b$
\begin{equation*}
        \sqrt{\frac{ (1-\psi)n\rho(\bs X)}{2}}{\beta_{\min}^{\star}}_j \,- \,\sqrt{\ln\big(p_j-|S^L_j(\kappa^{\circ})|\big)+\frac12\ln(n)} \;=\; \sqrt{\ln(s_j)} + c_{j} .
    \end{equation*}
where  $\psi=\frac12(1+\max_j \frac{\ln(p_j-s_j)}{\ln(p_j-s_j)+\ln(\sqrt{n})})$.
\end{enumerate}


\begin{thm}\label{theo:informedl0consist}
    Assume A1, A2, A10, $|S^S_j|=O(p_j-s_j)$ and $s_j / (p_j-s_j)<+\infty$ \david{isn't the latter condition just that $s_j < p_j$? If so, maybe write that for simplicity.}\paul{No it's not, you could have $s_j=n$, $p_j-s_j=\ln(n)$ where $s_j<p_j$ for all $n$ but $s_j/(p_j-s_j)\to \infty$. It's the negation of $p_j-s_j=o(s_j)$ also denoted $p_j-s_j = \Omega(s_j)$, it's also $\lim s_j/p_j <1$ } for every $j=1,\ldots,b$. Then 
$P(\hat{S}^{IB\ell_0} = S) \to 1$. \david{as $n$ and $p$ grow? Or as $n$ grows only? Same clarification for Theorem 5.1.} \paul{A1 and A2 are $n\to+\infty$ and $p-s\to\infty$ so it's assuming $n$ and $p$ grow.}
\end{thm}
\paul{ADD under A1, A2, $P(\hat{S}^{IB\ell_0} \subseteq S) \to 1$}

\section{Empirical results}\label{sec:simulation}

\section{Discussion}\label{sec:discussion}

We provided an understanding of how external information can ease variable selection in the Gaussian Sequence model and high-dimensional linear regression. We studied the case where external information helps partition variables into blocks with potentially distinct characteristics and introduced externally-informed $\ell_0$ selectors that let the penalty depend on those blocks. We showed that an oracle externally-informed selector converges faster and has less strict betamin condition than the standard $\ell_0$ oracle. We also provided concrete data analysis methods that incorporate external information to improve variable selection properties without requiring oracle knowledge.

A question that immediately arise for future work is how much the assumption of fixed number of blocks $b$ can be relaxed. Our consistency results are robust to moderate increases in the number of blocks but cannot hold when $b=p$ for example. Another interesting direction of research is understanding the benefit of external information for parameter estimation and estimation error. \paul{Adapting a proof in Rigollet's high-dim statistics notes I could easily give here a bound on estimation error in the Gaussian sequence model. I don't how tight it is though. For linear regresison, there is proof for BIC in Rigollet's notes and I need to check more closely if it can be adapted. Do you know David a proof for the estimation error with $\ell_0$ selection in linear regression? Would it be interesting/relevant to have those bounds here?} \paul{We could also mention externally informed $\ell_1$ in correlated deigns as future direction.}
\david{Castillo, Schmidt-Hieber and van der Vaart have work on estimation error for the sequence model and linear regression, using posterior model probabilities that are analogs to our pseudo-posterior probabilities. I suppose that one could adapt their proof strategy. Miquel Torrens also studied MSE's in the sequence model, his proof strategy was different and more along the lines of what you did here. This said, I'd probably leave this for future work. It might be an easy-ish paper, perhaps.}

We also observe that our results bear similarities with conclusions of works on related problems. In Section \ref{sec:examples}, we discussed how the benefits of externally-informed selectors depend on how discriminative is the external information. \cite{scarlett2012compressed} similarly showed that support recovery in compressed sensing becomes easier as the entropy of the distribution of inclusion probabilities across blocks shrinks. In hypothesis testing, \cite{genovese2006} demonstrated that the power of their informed weighted testing procedure increases when assignment of weights is positively associated with the null hypotheses being false. We expect analogous results to be achievable in other related problems like structural learning in graphical models or variable selection over partially ordered sets as introduced in \cite{taeb}. 

%%%%%%%%%%%%%%%%%%%%%%%%%%%%%%%%%%%%%%%%%%%%%%
%% Single Appendix:                         %%
%%%%%%%%%%%%%%%%%%%%%%%%%%%%%%%%%%%%%%%%%%%%%%
%\begin{appendix}
%\section*{???}%% if no title is needed, leave empty \section*{}.
%\end{appendix}
%%%%%%%%%%%%%%%%%%%%%%%%%%%%%%%%%%%%%%%%%%%%%%
%% Multiple Appendixes:                     %%
%%%%%%%%%%%%%%%%%%%%%%%%%%%%%%%%%%%%%%%%%%%%%%
%\begin{appendix}
%\section{???}
%
%\section{???}
%
%\end{appendix}

%%%%%%%%%%%%%%%%%%%%%%%%%%%%%%%%%%%%%%%%%%%%%%
%% Support information, if any,             %%
%% should be provided in the                %%
%% Acknowledgements section.                %%
%%%%%%%%%%%%%%%%%%%%%%%%%%%%%%%%%%%%%%%%%%%%%%
\begin{acks}[Acknowledgments]
The authors would like to thank ...
\end{acks}
%%%%%%%%%%%%%%%%%%%%%%%%%%%%%%%%%%%%%%%%%%%%%%
%% Funding information, if any,             %%
%% should be provided in the                %%
%% funding section.                         %%
%%%%%%%%%%%%%%%%%%%%%%%%%%%%%%%%%%%%%%%%%%%%%%
\begin{funding}
The first author was supported by ...

The second author was supported in part by ...
\end{funding}

%%%%%%%%%%%%%%%%%%%%%%%%%%%%%%%%%%%%%%%%%%%%%%
%% Supplementary Material, including data   %%
%% sets and code, should be provided in     %%
%% {supplement} environment with title      %%
%% and short description. It cannot be      %%
%% available exclusively as external link.  %%
%% All Supplementary Material must be       %%
%% available to the reader on Project       %%
%% Euclid with the published article.       %%
%%%%%%%%%%%%%%%%%%%%%%%%%%%%%%%%%%%%%%%%%%%%%%
%\begin{supplement}
%\stitle{???}
%\sdescription{???.}
%\end{supplement}

%%%%%%%%%%%%%%%%%%%%%%%%%%%%%%%%%%%%%%%%%%%%%%%%%%%%%%%%%%%%%
%%                  The Bibliography                       %%
%%                                                         %%
%%  imsart-???.bst  will be used to                        %%
%%  create a .BBL file for submission.                     %%
%%                                                         %%
%%  Note that the displayed Bibliography will not          %%
%%  necessarily be rendered by Latex exactly as specified  %%
%%  in the online Instructions for Authors.                %%
%%                                                         %%
%%  MR numbers will be added by VTeX.                      %%
%%                                                         %%
%%  Use \cite{...} to cite references in text.             %%
%%                                                         %%
%%%%%%%%%%%%%%%%%%%%%%%%%%%%%%%%%%%%%%%%%%%%%%%%%%%%%%%%%%%%%

% if your bibliography is in bibtex format, uncomment commands:
\bibliographystyle{imsart-number} % Style BST file (imsart-number.bst or imsart-nameyear.bst)
\bibliography{bibliography.bib}       % Bibliography file (usually '*.bib')


XXXXXXXXXXXXXXXXXXXXXXXXXXXXXXXXXXXXXXXXXXXXXXXXXXXXXXXXXXXXXXXXX\\
EVERYTHING AFTER WILL GO TO SUPPLEMENT\\
XXXXXXXXXXXXXXXXXXXXXXXXXXXXXXXXXXXXXXXXXXXXXXXXXXXXXXXXXXXXXXXXX\\
\section{Auxiliary results}\label{auxiliary}

In this section we collect some technical results. Their proofs are gathered in section~\ref{sec:proofauxi}. 

\begin{lem}\label{lemma:3tailbounds}\textbf{Tail bounds on the maximum and minimum of folded Gaussians} 
\begin{enumerate}[a)]
    \item if $\bs y\sim N_p\left(0,\frac{1}{n}I_n \right)$, $p>1$, and $a \geq \sqrt{\frac{2\ln(p)}{n}}$, $$P\left(\underset{i\in\{1,\ldots,p\}}{\max}|y_i|> a \right) \leq \frac{e^{-\frac{n}{2}\left(a^2-\frac{2\ln (p)}{n}\right)}}{\sqrt{\pi\ln (p)}}$$
    \item if $\bs y\sim N_s\left(\bs\beta,\sigma^2\right)$, $s>1$, and $a \leq \underset{i\in\{1,\ldots,s\}}{\min}|\beta_i|$,
        $$P\left(\underset{i\in\{1,\ldots,s\}}{\min}|y_i| > a\right) \geq P\left(\underset{i\in\{1,\ldots,s\}}{\min}|\beta_i|-\underset{i\in\{1,\ldots,s\}}{\max}|y_i-\beta_i| > a \right)$$
    \item if $\forall i\in\{1,\ldots,s\}$ $y_i\sim N\left(\mu,\frac{1}{n}\right)$, $a<|\mu|$
        $$P\left(\underset{i\in\{1,\ldots,s\}}{\min}|y_i| > a \right) \leq \exp\left\{-\frac{s}{2}\frac{e^{-2n(|\mu|-a)^2/\pi}-e^{-n(a+|\mu|)^2/2}}{\left(1-e^{-2n(|\mu|-a)^2/\pi}\right)^{\frac{1}{2}}+\left(1-e^{-n(a+|\mu|)^2/2}\right)^{\frac{1}{2}}}\right\} $$
\end{enumerate}
\end{lem}


% \section{Proof of Theorem~\ref{theo:suffcondlinearmodel}}

\begin{lem}\label{lem:L0toL1convergencegen} For $\hat S^b$ as in \eqref{eq:Mhat} and any $k$ models $M_1,\ldots,M_k$
$$P(\hat{S}^b\not\in \{M_1,\ldots,M_k\})\;\leq\; (k+1) \sum_{M \in \M\setminus \{M_1,\ldots,M_k\}}\E\left( NC(M)\right).$$
\end{lem}

\begin{lem}\label{lem:boundrhogen}
 For any $M\in\M$ and $T\subseteq S$, let $Q_T=M \cup T$. The non-centrality parameter defined in \eqref{eq:muQM} satisfies: 
\begin{equation}
    \mu_{Q_TM}\; \geq \;n\,\rho(\bs X)\,\sum_{j=1}^b|T_j\setminus M_j|\,\min_{i\in T_j\setminus M_j}{\beta_{i}^{\star}}^2
\end{equation}
\end{lem}

\begin{lem}\label{lemma:s1s3}
Let $W\sim\chi_{\nu}^2(\mu)$ with $\mu \geq 0$, then for any $w>\mu+\nu$
$$P(W > w) \leq e^{-\big(\frac{w+\mu}{2}  - \sqrt{2w( 2 \mu+ \nu)-2 \mu \nu - \nu^2 } \big)} $$
Moreover, if $w$ is increasing such that $\nu=o(w)$ and $\mu=o(w)$, for any $\phi\in(0,1)$ and $n$ large enough
$$P(W > w)  \leq  e^{-\phi\frac{w}{2}}$$
\end{lem}

\begin{lem}\label{lemma:s2} Let $W \sim \chi_\nu^2(\mu)$ with $\mu>0$. For any $w<\mu$, 
$$
P(W<w) \leq \frac{e^{-\frac{1}{2}(\sqrt{\mu}-\sqrt{w})^2}}{(\mu / w)^{\nu / 4}}.$$
Moreover, if $w$ is increasing such that $w=o(\mu)$, for any $\phi\in(0,1)$, for $n$ large enough:
$$
P(W<w) \leq e^{-\phi\frac{\mu}{2}}.$$
\end{lem}

\begin{lem}\label{lemma:s18} 
    Let $W\sim \chi_\nu^2(\mu)$ with $\mu \geq 0$. Assume $g$ positive and increasing such that $\nu=o(\ln(g))$, $\mu=o(\ln(g))$. Let $\bar{u},\underline{u}$ in $(0,1)$ such that $1>\bar{u}>\underline{u}\geq\left(1+g^{\phi}e^{-(\nu+\mu)/2}\right)^{-1}$ where $\phi\in(0,1)$, then for $n$ large enough, we have
$$
\int_{\underline{u}}^{\bar{u}} P\left(W>2\ln \left(\frac{g}{1 / u-1}\right)\right) d u \leq  \frac{1}{g^{\phi}}\Big(\bar{u}-\underline{u}+\phi\ln\Big(\frac{\bar{u}}{\underline{u}}\Big)\Big)
$$

\end{lem}

\begin{lem}\label{lemma:s20}
 Let $W\sim \chi_\nu^2(\mu)$ with $\mu \geq 0$. Assume $g$ positive and increasing such that $\nu=o(\ln(g))$, $\mu=o(\ln(g))$, then for any $\alpha\in(0,1)$ and $n$ large enough, we have
$$
\int_0^1 P\left(W>2 \ln \left(\frac{g}{1 / u-1}\right)\right) d u\;\; \leq\;\;  g^{-\alpha}$$
\end{lem}

\begin{lem}\label{lem:boundexpnc} For $M\in \M\setminus \{T\}$ 
Denote $Q_T= M\cup T$ and $A_T:=\gamma\Delta_{MT}+\tfrac{1-\gamma}{2}\mu_{Q_TM}$ (cf \eqref{eq:DeltaM} and \eqref{eq:muQM}). Suppose that, for some $\gamma \in (1/2,1]$, $|M\setminus T|=o(A_T)$, $\mu_{Q_T T}=o(A_T)$ and $|T\setminus M|=o(\mu_{Q_TM})$ when $T\setminus M\neq \emptyset$. For any $\psi\in(0,1)$ and $n$ large enough:
$$\E\left(NC(M)\right)  \;\;\leq\;\; e^{-\psi A_T}.$$
\end{lem}

\begin{lem}\label{lem:upperboundnoncentralgen}
    For any $T\subseteq S$ and $M\in \M$, let $Q_T=M\cup T$, $$\mu_{Q_TT} \leq n\,\bar{\lambda}\,\sum_{j=1}^b|(S_j\cap M_j )\setminus T_j|\,\max_{i\in (S_j\cap M_j )\setminus T_j}{\beta_{i}^{\star}}^2 \quad \text{where}\quad \bar{\lambda}\,:=\,\lambda_{\max}\Big(\frac1n \bs X_{S}^\top\bs X_{S}\Big)$$
    and
    $$\mu_{Q_TT} \geq n\,\lambda_{\min}\Big(\frac1n \bs X_{S}^\top\bs X_{S}\Big)\,\sum_{j=1}^b|(S_j\cap M_j )\setminus T_j|\,\min_{i\in (S_j\cap M_j )\setminus T_j}{\beta_{i}^{\star}}^2$$
\end{lem}

\section{Proofs of Section~\ref{sec:properties_orth}}


\subsection{Proof of Lemma~\ref{lem:useful-reformulation-GSM}}
First note that by discarding constant term 
 $$\arg\max_{M\in \M} \left\{\max_{\beta\in \mathcal{L}_M}\ell(\bs y; \bs \beta)-\sum_{j=1}^b\kappa_j |M_j| \right\}=\arg\min_{M\in \M} \left\{\min_{\beta\in \mathcal{L}_M}\bigg\{ \tfrac{1}{2}\|\bs y-\sqrt{n}\bs\beta\|^2\right\}+\sum_{j=1}^b\kappa_j |M_j| \bigg\}$$
 
 We also have $\min_{\beta\in \mathcal{L}_M} \tfrac{1}{2}\|\bs y-\sqrt{n}\bs\beta\|^2=\tfrac{1}{2}\|\bs y-\sqrt{n}\tilde{\bs \beta}_M\|^2=\tfrac{1}{2}\|\bs y\|^2-\tfrac{1}{2}\|\sqrt{n}\tilde{\bs \beta}_M\|^2$. The minimization in \eqref{eq:Mhat} can be then replaced with maximization of $\tfrac{1}{2}
 \|\sqrt{n}\tilde{\bs \beta}_M\|^2-\sum_{j=1}^b \kappa_j|M_j|$.

\subsection{Proof of Proposition~\ref{prop:isthresh}}
	Under the assumptions, we can write 
	$$
	\tfrac{1}{2}\|\sqrt{n}\tilde{\bs \beta}_M\|^2-\sum_{j=1}^b\kappa_j |M_j|\;=\;\sum_{j=1}^b\sum_{i\in M_j}(\tfrac{n}{2}{\tilde{ \beta}_i}^2-\kappa_j).
	$$
	Now, Lemma~\ref{lem:useful-reformulation-GSM} gives that \eqref{eq:MhatGSM} can be maximized with respect to each $M_j$ by including $i\in \hat S_j$ whenever $n\tilde\beta_i^2\geq 2\kappa_j$.

\subsection{Proof of Lemma~\ref{lem:l1isthreshold}}
    Denote $\hat{\bs \beta}$, the minimizer of the problem \eqref{eq:genl1}. The MLE under the full model is $\tilde{\bs\beta}=\tfrac{1}{\sqrt{n}} \bs y$ and the optimized function in \eqref{eq:genl1} can be then rewritten as
	$$
	\tfrac{1}{2}\bs y^{\top}\bs y +\tfrac{n}{2}\sum_{j=1}^b\sum_{i\in B_j}\left(\beta_i^2-2\tilde\beta_i\beta_i+2\tfrac{\lambda_j}{n}|\beta_i|\right)
	$$
	which we optimize with respect to each $\beta_i$ separately. For any $i\in B_j$ we have $\hat\beta_i=0$ if and only if $|\tilde \beta_i|\leq \frac{\lambda_j}{n}$. Similarly  for \eqref{eq:genl1ada}, denoting $\hat{\bs \beta}$ the minimizer of the problem, for any $i\in B_j$ we have $\hat\beta_i=0$ if and only if $|\tilde \beta_i|\leq \frac{\lambda_j}{|\beta^{\circ}_i|n}$. If $\bs \beta^{\circ}= \tilde{\bs\beta}$, then for any $i\in B_j$ $\hat{\beta}_i = 0$ if and only if $|\tilde \beta_i|\leq \sqrt{\tfrac{\lambda_j}{n}}$. If $\bs\beta^\circ$ is a LASSO estimate with penalization $\lambda^\circ$, then for any $i\in B_j$ 
 $\hat{\beta}_i = 0$ if and only if $\tilde{\beta}^2_i + \operatorname{sign}(\lambda^{\circ}/n-\tilde{\beta}_i) \frac{\lambda^\circ}{n}\tilde{\beta}_i-\frac{\lambda_j}{n} \leq 0$. Equivalently, for any $i\in B_j$ 
 $\hat{\beta}_i = 0$ if and only if $|\tilde{\beta}_ i| \leq  \frac{1}{2}\left(\lambda^\circ/n+\sqrt{(\lambda^\circ/n)^2+4 \lambda_j/n}\right)$.


\subsection{Proof of Proposition~\ref{prop:suffnecIblockthresh}} By the union bound and by Lemma~\ref{lemma:3tailbounds} a), for $n\geq n_0$, \begin{equation*}
P\left(\hat{S}^{T,b} \not\subseteq S\right)\;\leq\; \sum_{j=1}^b P\Big(\max_{i\in B_j\setminus S_j}|y_i/\sqrt{n}|>\tau_{j}\Big)\;\leq\;\sum_{j=1}^b\frac{e^{-\tfrac{n}{2}\left(\tau_j^2-\tfrac{2\ln(p_j-s_j)}{n}\right)}}{\sqrt{\pi\ln(p_j-s_j)}} 	.
 \end{equation*}
 By A4, the nominator on the right is bounded above by $1$ for all sufficiently large $n$. It follows that 
	$$P\left(\hat{S}^{T,b} \not\subseteq S\right) \;\leq\; \sum_{j=1}^b (\pi\ln (p_j-s_j))^{-1/2} \to 0\qquad \mbox{as } n\to \infty.$$
For the second statement, by independence, we have 
$$P(\hat{S}^{T,b} \subseteq S)\;=\;\prod_{j=1}^b P(\max_{i \in B_j\setminus S_j}|y_i/\sqrt{n}| \leq \tau_j).$$ 
Consider $j$ that satisfies \eqref{cond:necIallblockthresh0}. In particular, there exists $c<1$ such that, for sufficiently large $n$, we have that $\tau_j\leq c\sqrt{2\ln(p_j-s_j)/n}$. Then, for any such $n$,
$$
P\Big(\max_{i \in B_j\setminus S_j}|y_i/\sqrt{n}| \leq \tau_j\Big)\;\leq\; P\Big(\max_{i \in B_j\setminus S_j}|y_i| \leq c \sqrt{2\ln(p_j-s_j)}\Big)\;\leq\;P\Big(\tfrac{1}{\sqrt{2\ln(p_j-s_j)}}\max_{i \in B_j\setminus S_j}y_i \leq c \Big).
$$
On the other hand, results from extreme value theory, Chapter 4 in \cite{galambos}, show that 
 $$\frac{1}{\sqrt{2 \ln(p_j-s_j)}}\max_{i\in B_j\setminus S_j}y_i \;\overset{P}{\longrightarrow} 1$$
 and so 
 $$
 P\Big(\frac{1}{\sqrt{2\ln(p_j-s_j)}}\max_{i \in B_j\setminus S_j}y_i \leq c \Big)\;\to\;0,
 $$ 
 which implies that $P(\hat{S}^{T,b} \subseteq S)\to 0$. 

\subsection{Proof of Proposition~\ref{prop:suffnecIIblockthresh}}
	 By the union bound, 
	$P\left(S \not\subseteq \hat{S}^{T,b}\right)\leq \sum_{j=1}^b  P\left(\min_{i\in S_j}|y_i/\sqrt{n}|\leq\tau_j\right)$. 
	By Lemma~\ref{lemma:3tailbounds} b), for each $j$,	
	$$P\left(\min_{i\in S_j}|y_i/\sqrt{n}|\leq\tau_j\right)\;\leq\; P\left(\max_{i\in S_j}|\tfrac{y_i}{\sqrt{n}}-\beta_i| \geq\beta_{\min,j}^{\star}-\tau_j \right).$$ By Lemma~\ref{lemma:3tailbounds} a)  
	\begin{equation}\label{eq:rate2}
P\left(\max_{i\in S_j}|\tfrac{y_i}{\sqrt{n}}-\beta_i| \geq\beta_{\min,j}^{\star}-\tau_j \right) \;\leq\;\frac{e^{-\tfrac{n}{2}\left((\beta^\star_{\min,j}-\tau_j)^2-\tfrac{2\ln(s_j)}{n}\right)}}{\sqrt{\pi\ln(s_j)}}\;\leq\; (\pi\ln (s_j))^{-1/2}\;\to\;0,		
	\end{equation}
	which concludes the proof. 

\subsection{Proof of Proposition~\ref{prop:necIIblockthresh}} 
a) By independence we have 
$P(S \subseteq \hat{S}^{T,b})=\prod_{j=1}^b P(\min_{i \in S_j}|y_i/\sqrt{n}| > \tau_j)$. Take any $j$, satisfying \eqref{cond:necIIallblockthresh2}. $P(\min_{i \in S_j}|y_i/\sqrt{n}| > \tau_j)=\prod_{i \in B_j} P(|y_i/\sqrt{n}| > \tau_j)$. Denote $i_{\circ}\in B_j$ the entry, such that $\beta_i^*=\beta^*_{\min,\,j}$. From the proof of Lemma \ref{lemma:3tailbounds} b) $$P(|y_{i_{\circ}}/\sqrt{n}|>\tau_j) 
    = P\left(z>\sqrt{n}\left(\tau_j-\beta_{\min,\,j}^{\star}\right)\right) + P\left(z<-\sqrt{n}\left(\tau_j+\beta_{\min,\,j}^{\star}\right)\right)$$ where $z\sim N(0,1)$. Since $\tau_j<\beta^\star_{\min,j}$ satisfies condition \eqref{cond:necIallblockthresh} $P(z<-\sqrt{n}\left(\tau_j+\beta_{\min,\,j}^{\star}\right)) \to 0$. By \eqref{cond:necIIallblockthresh2} $\underset{n \to +\infty}{\lim}P\left(z>\sqrt{n}\left(\tau_j-\beta_{\min,\,j}^{\star}\right)\right)<1$ and $\underset{n \to +\infty}{\lim}P\left(S_j \subseteq \hat{S}^{T,b} \right)<1$.\\
b)	Take any $j=1,\ldots,b$ satisfying \eqref{cond:necIIallblockthresh} By Lemma~\ref{lemma:3tailbounds} c), since $\beta^\star_{\min,j}>\tau_j$:
$$P\left(\min_{i \in S_j}|\frac{y_i}{\sqrt{n}}| > \tau_j \right)\;\leq\;\exp\left\{-\frac{s_j}{2}\frac{e^{-\tfrac{2n}{\pi}(\beta^\star_{\min,j}-\tau_j)^2}-e^{-\tfrac{n}{2}(\tau_j+\beta^\star_{\min,j})^2}}{\left(1-e^{-2n(\beta^\star_{\min,j}-\tau_j)^2/\pi}\right)^{\frac{1}{2}}+\left(1-e^{-n(\tau_j+\beta^\star_{\min,j})^2/2}\right)^{\frac{1}{2}}}\right\}. $$
Let $a_{n}=2\left(1-e^{-2n(\beta^\star_{\min,j}-\tau_j)^2/\pi}\right)^{\frac{1}{2}}+2\left(1-e^{-n(\tau_j+\beta^\star_{\min,j})^2/2}\right)^{\frac{1}{2}}$ and note that $a_n\in (0, 4]$.  Thus, to show that $P\left(\min_{i \in S_j}|\tfrac{y_i}{\sqrt{n}}| > \tau_j \right)$ is bounded away from $1$, it is enough to show that 
$$
\lim_{n\to \infty}{s_j}({e^{-2n(\beta^\star_{\min,j}-\tau_j)^2/\pi}-e^{-n(\tau_j+\beta^\star_{\min,j})^2/2}})\;>\;0.
$$
By \eqref{cond:necIIallblockthresh}, $\lim {s_j} e^{-2n(\beta^\star_{\min,j}-\tau_j)^2/\pi}\geq 1$. To conclude, we show that there exists $c\in(0,1)$ such that
$$
\lim {s_j} e^{-n(\tau_j+\beta^\star_{\min,j})^2/2}\;\leq\;c\;<\;1,
$$ 
where $c$ is as defined in the proposition. We equivalently need
$\sqrt{n}(\tau_j+\beta^\star_{\min,j})\geq\sqrt{{2\ln (s_j/c)}}$ for all $n$ sufficiently large. To show that, note that $\tau_j\geq \beta^\star_{\min,j}$ and $s_j/c\leq p_j$, which gives
$$
\frac{\tau_j+\beta^\star_{\min,j}}{\sqrt{\tfrac{2\ln(s_j/c)}{n}}}\;\geq\;\frac{2\tau_j}{\sqrt{\tfrac{2\ln(p_j)}{n}}}\;=\;\frac{\sqrt{\tfrac{2\ln(p_j-s_j)}{n}}}{\sqrt{\tfrac{2\ln(p_j)}{n}}}\frac{2\tau_j}{\sqrt{\tfrac{2\ln(p_j-s_j)}{n}}}\;=\;\sqrt{\frac{\ln(p_j-s_j)}{\ln(p_j)}}\frac{2\tau_j}{\sqrt{\tfrac{2\ln(p_j-s_j)}{n}}}.
$$
Since the second term on the right converges to something $\geq 2$, it is enough to show that the first term converges to something $>1/2$. Using our assumption that $\tfrac{s_j}{p_j}<1$, we can take $c$ such that $\tfrac{s_j}{p_j}\leq c<1$ and we get
$$
\sqrt{\frac{\ln(p_j-s_j)}{\ln(p_j)}}\;\geq\;\sqrt{\frac{\ln(p_j)+\ln(1-c)}{\ln(p_j)}}\;\to\;1,
$$
which concludes the proof.

\subsection{Proof of Lemma~\ref{lem:necbetaminblockthresh}}
    We have $\frac{\beta_{\min,j}^*}{\sqrt{\tfrac{2\ln(p_j-s_j)}{n}}+\sqrt{\tfrac{\pi}{2}\tfrac{\ln(s_j)}{n}}}= \frac{\beta_{\min,j}^*-\tau_j+\tau_j}{\sqrt{\tfrac{2\ln(p_j-s_j)}{n}}+\sqrt{\tfrac{\pi}{2}\tfrac{\ln(s_j)}{n}}}$ and \eqref{eq:necbetaminblockthresh} implies that $\lim_{n\to \infty}\frac{\tau_j}{\sqrt{\tfrac{2\ln(p_j-s_j)}{n}}} <1$ or $\lim_{n\to \infty}\frac{\beta_{\min,j}^*-\tau_j}{\sqrt{\tfrac{\pi}{2}\tfrac{\ln(s_j)}{n}}} <1$. 
    If $\lim_{n\to \infty}\frac{\tau_j}{\sqrt{\tfrac{2\ln(p_j-s_j)}{n}}} <1$ by Proposition~\ref{prop:suffnecIblockthresh}, $P(\hat{S}^{T,b} = S)\not \to 1$. If $\beta_{\min,j}<\tau_j$ recovery is not possible either by definition. 
    The last case to examine is when $\beta_{\min,j}\geq\tau_j$, $\lim_{n\to \infty}\frac{\tau_j}{\sqrt{\tfrac{2\ln(p_j-s_j)}{n}}} \geq 1$ and $\lim_{n\to \infty}\frac{\beta_{\min,j}^*-\tau_j}{\sqrt{\tfrac{\pi}{2}\tfrac{\ln(s_j)}{n}}} <1$. By Proposition~\ref{prop:necIIblockthresh}, in that case $P(S \subseteq \hat{S}^{T,b} )\not \to 1$ and recovery is not possible asymptotically.

\subsection{Proof of Theorem~\ref{theo:fullrecovGS}}
\paul{XXXXXX ADD proof about bound on convergence rate XXXXXX}
We start by showing that if \eqref{eq:betamin} holds then $\tau^\star_j$ satisfies A4 and A5. Under \eqref{eq:betamin}, there exists $c>1$ such that $\beta^\star_{\min,j}=c(\sqrt{\tfrac{2\ln(p_j-s_j)}{n}} + \sqrt{\tfrac{2\ln(s_j)}{n}})$. Denote $a=\sqrt{\tfrac{2\ln(p_j-s_j)}{n}}$ and $b=\sqrt{\tfrac{2\ln(s_j)}{n}}$. We have
$$
\tau_j^\star\;=\;\frac{c}{2} (a+b)+\frac{(a^2-b^2)}{2c(a+b)}\;=\;\frac{c}{2} (a+b)+\frac{a-b}{2c}\;=\;\frac{c^2+1}{2c}a+\frac{c^2-1}{2c}b.
$$
The fact that $\tau^\star_j>a$ (the lower bound in A4) is now clear because $\tfrac{c^2+1}{2c}> 1$  and $\tfrac{c^2-1}{2c}> 0$. To show that the upper bound in A5 observe that
$$
\beta_{\min,j}^* - \tau^*_j \;=\; c(a+b)- \frac{c^2+1}{2c}a+\frac{c^2-1}{2c}b \;= \; \frac{c^2+1}{2c}a+\frac{c^2-1}{2c}b\;>\;  b 
$$
Then there exists $n_0$ such that for all $n\geq n_0$ \eqref{eq:rateconvgsm} holds. For convenience, denote $d = \frac{1}{n {\beta_{\min}^*}_j} \ln (p_j / s_j-1)$ such that $\tau^*_j= {\beta_{\min}^*}_j/2 + d$. Then
$$e^{-\frac{n}{2}\left({\tau^*_j}^2-\frac{2\ln(p_j-s_j)}{n}\right)}
=e^{-\left[\frac{n}{2}{\beta_{\min}^{\star}}_j^2+\frac{n}{2}d^2 + \frac{1}{2}\left(\ln (p_j- s_j)-\ln (s_j)\right) -\ln (p_j- s_j)\right]}$$
Then, by considering separately the two posssible maximum in $\max\{\ln(p_j-s_j),\ln(s_j)\}$, we get that:
$$\frac{e^{-\frac{n}{2}\left({\tau^*_j}^2-\frac{2\ln(p_j-s_j)}{n}\right)}}{e^{-[\frac{n}{2}{{\beta_{\min}^*}_j}^2-\max\{\ln(p_j-s_j),\ln(s_j)\}]}}= 
e^{-\left[\frac{n}{2}d^2 + \frac{1}{2}\left|\ln (p_j- s_j)-\ln (s_j)\right|\right]} <1$$
The bound in \eqref{eq:optimrateconvgsm} follows from noting that $e^{-\tfrac{n}{2}\big[{\tau_j^*}^2-\tfrac{2\ln(p_j-s_j)}{n}\big]}=e^{-\tfrac{n}{2}\big[(\beta^\star_{\min,j}-{\tau_j^*})^2-\tfrac{2\ln(s_j)}{n}\big]}$ and summing over $j$.


\section{Proofs of Section~\ref{sec:properties_legression}}
\subsection{Proof of Lemma~\ref{lem:useful-reformulation}}
	First note that by discarding constant term 
 $$\arg\max_{M\in \M} \left\{\max_{\beta\in \mathcal{L}_M}\ell(\bs y; \bs \beta)-\sum_{j=1}^b\kappa_j |M_j| \right\}=\arg\min_{M\in \M} \left\{\min_{\beta\in \mathcal{L}_M}\bigg\{ \tfrac{1}{2}\|\bs y-\bs X\bs\beta\|^2\right\}+\sum_{j=1}^b\kappa_j |M_j| \bigg\}$$
 
 We also have $\min_{\beta\in \mathcal{L}_M} \tfrac{1}{2}\|\bs y-\bs X\bs\beta\|^2=\tfrac{1}{2}\|\bs y-\bs X_{\!\!M}\tilde{\bs \beta}^{(M)}\|^2=\tfrac{1}{2}\|\bs y\|^2-\tfrac{1}{2}\|\bs X_{\!\!M}\tilde{\bs\beta}^{(M)}\|^2$. The minimization in \eqref{eq:Mhat} can be then replaced with maximization of $\tfrac{1}{2}
 \|\bs X_{\!\!M}\tilde{\bs\beta}^{(M)}\|^2-\sum_{j=1}^b \kappa_j|M_j|$.

\subsection{Proof of Lemma~\ref{lem:L0toL1convergence}}
The first statement follows directly from Lemma~\ref{lem:L0toL1convergencegen} by taking $\{M_1,\ldots,M_k\}=\{S\}$. For the second statement, note that  $NC(M)=(1+\sum_{N \neq M} e^{C(N)}/ e^{C(M)})^{-1}$. Since for any $M,N\in \M$ and $M\neq N$, $C(N) / C(M)>0$, $NC(M) \leq (1+e^{C(M')-C(M)})^{-1}$.

\subsection{Proof of Lemma~\ref{lemma:s7}}
    This result follows directly from Lemma~S7 in \cite{rossell2022concentration} taking $\phi^*=1$.

\subsection{Proof of Lemma~\ref{lem:boundrho}}
This result follows directly Lemma~\ref{lem:boundrhogen} by taking $T=S$, and observing that $\min_{i\in S_j\setminus M_j}{\beta^*_i}^2\geq {\beta^*_{\min,j}}^2$.

\subsection{Proof of Theorem~\ref{theo:suffcondlinearmodel}} We start by showing that $A_S$ as defined in Lemma~\ref{lem:boundexpnc} for $T=S$ satisfies the conditions of the lemma. Observe that $\Delta_{MS}$, defined in \eqref{eq:DeltaM}, can be rewritten as $\Delta_{MS}=\sum_{j=1}^b(|M_j\setminus S_j|-|S_j\setminus M_j|)\kappa_j$. By  Lemma~\ref{lem:boundrho}, for every $n\in \N$ we have 
\begin{equation}\label{eq:nonoverfittingcond}
   A_S\;=\; \gamma \Delta_{MS}+\frac{1-\gamma}{2}\mu_{Q_SM}\;\geq\; \gamma\sum_{j=1}^b|M_j \setminus S_j|\kappa_j + \sum_{j=1}^b |S_j\setminus M_j |\left(\tfrac{1-\gamma}{2} n \rho(\bs X) {\beta^{\star}_{\min,j}}^{2} - \gamma\kappa_j\right) 
\end{equation} 
Since $M\neq S$, $|M\setminus S|\neq0$ or $|S\setminus M|\neq0$, then by assumptions A6 and A7, for $n$ large enough $A_S>0$. We immediately have $\mu_{Q_S S}=o(A_S)$ since $\mu_{Q_S S}=0$. If $|M\setminus S|=0$, $|M\setminus S|=o(A_S)$ also immediately. Consider now the case $|M\setminus S|\neq 0$. By assumption A7, the last term in \eqref{eq:nonoverfittingcond} is nonnegative and:
$$\frac{|M\setminus S|}{A_S}\leq \bigg[ \gamma\sum_{j=1}^b\frac{|M_j \setminus S_j|}{|M\setminus S|}\kappa_j \bigg]^{-1} \leq \Big[ \gamma\min_{j=1,\ldots,b}\kappa_j \Big]^{-1} $$
where the last inequality follows from $\sum_{j=1}^b\frac{|M_j \setminus S_j|}{|M\setminus S|}=1$.  When $|S\setminus M|\neq0$, by Lemma~\ref{lem:boundrhogen}
$$\frac{|S\setminus M|}{\mu_{Q_SM}} \leq \frac{|S\setminus M|}{n\,\rho(\bs X)\,\sum_{j=1}^b|S_j\setminus M_j|\,{\beta_{\min,j}^{\star}}^2}\leq \Big[n\,\rho(\bs X)\,\max_{j=1,\ldots,b}{\beta_{\min,j}^{\star}}^2\Big]^{-1}  $$
where the last inequality follows from $\frac{|S\setminus M|}{\sum_{j=1}^b|S_j \setminus M_j|}=1$. By assumption A7, we get have $|S\setminus M|=o(\mu_{Q_SM})$ when $|S\setminus M|\neq 0$ too.
We conclude that for every $M\in \mathcal M\setminus \{S\}$, by Lemma~\ref{lem:boundexpnc}, for any $\psi\in(0,1)$ and $n$ large enough:
$$
\E(NC(M))\;\leq\;e^{-\psi A_S}.
$$
Let 
$$A_S^*\;:=\;\gamma\sum_{j=1}^b|M_j \setminus S_j|\kappa_j + \sum_{j=1}^b |S_j\setminus M_j |\left(\tfrac{1-\gamma}{2} n \rho(\bs X) {\beta^{\star}_{\min,j}}^{2} - \gamma\kappa_j\right) $$
By \eqref{eq:nonoverfittingcond}, we have for $n$ large enough:
\begin{equation}\label{eq:boundexp}
    \E(NC(M))\;\leq\;e^{-\psi A_S^*}.
\end{equation}
Assumption~A7 implies there exist $c^{\prime}_j \to +\infty$ such that:
$$\frac{(1-\gamma) n\rho(\bs X)}{2}{{\beta_{\min}^{\star}}_j}^2 \,- \,\kappa_j \;=\; \ln(s_j) + c^{\prime}_{j}.$$
Let $\delta \in (0,1)$ and denote $\bar{m}_j=\max\big\{\frac{2\ln(p_j-s_j)}{b_{j}},\frac{2\ln(s_j)}{c^{\prime}_{j}}\big\}$. Take $\psi=\max_{j=1,\ldots,b}\frac{\xi + \delta + \bar{m}_j}{1+\bar{m}_j}$ for some $\xi \in(0, 1- \delta)$ then $\psi \in (0,1)$ and we have for every $j=1,\ldots,b$:
$$\delta \leq \psi-(1-\psi)\bar{m}_j.$$ 
By definition of $\gamma$ we have: 
$$
\gamma\kappa_j\;\geq\;\frac12 (1+\alpha_j)\kappa_j\;=\;\ln(p_j-s_j)+\frac12(\kappa_j-\ln(p_j-s_j)) \;=\; \ln(p_j-s_j)+ \frac12b_{j}.
$$
From our choice of $\psi,\delta$, we get:
$$\psi\gamma\kappa_j \;\geq\; \psi\big(\ln(p_j-s_j)+ \frac12b_{j}\big) \;\geq\;  \ln(p_j-s_j)+ \delta\frac12b_{j}.$$
By A7 and by our choice of $\psi,\delta$,
$$\psi\left(\tfrac{1-\gamma}{2} n \rho(\bs X) {\beta^{\star}_{\min,j}}^{2} - \gamma\kappa_j\right) \;\geq\; \psi\left(\ln(s_j)+c_j'\right) \;\geq\; \ln(s_j)+\delta \frac{c^{\prime}_j}{2}.
$$
Using the expression for $A_S^*$, the bound in \eqref{eq:boundexp}, and the above inequalities, we get
$$
\E(NC(M))\;\leq\;\exp\{- \sum_{j=1}^b |M_j\setminus S_j|(\ln(p_j-s_j)+\delta \tfrac{b_j}{2})-\sum_{j=1}^b |S_j\setminus M_j|(\ln(s_j)+\delta \tfrac{c^{\prime}_j}{2})\}.
$$

Denote $\mathcal{S}= \sum_{M \in \M\setminus \{S\}} \E\left(NC(M)\right)$ for convenience. We can split $\mathcal{S}$ into sums over the models that have the same number of spurious variables and same number of missing variables in every block, that is the models $M$ such that for all $j$, $|M_j\setminus S_j|=u_j$ and $|S_j\setminus M_j|=w_j$ with $u_j \in \{0,\ldots,p_j-s_j\}$ and $w_j \in \{0,\ldots,s_j\}$, not all simultaneously zero:
$$\mathcal{S} =\sum_{w_1=0}^{s_1}\cdots\sum_{w_b=0}^{s_b}\sum_{u_1=0}^{p_1-s_1}\cdots\sum_{u_b=0}^{p_b-s_b}\sum_{M \neq S:\forall \,j |M_j\setminus S_j|=u_j, |S_j\setminus M_j|=w_j } \E\left(NC(M)\right)$$
The number of models with given collection of $u_j$'s and $w_j$'s is $\prod_{j=1}^b {p_j-s_j \choose u_j}{s_j \choose w_j}$. We thus have:
\begin{equation*}
    \mathcal{S} \leq \prod_{j=1}^b\left(1+\sum_{u_j=1}^{p_j-s_j}{p_j-s_j \choose u_j}e^{- u_j(\ln(p_j-s_j)+\delta \tfrac{b_j}{2})}\right)\left(1+
    \sum_{w_j=1}^{s_j}{s_j \choose w_j}
    e^{-w_j(\ln(s_j)+\delta \tfrac{c^{\prime}_j}{2})}\right) - 1
\end{equation*}
where $-1$ stems from discarding $S$ from the product-sum, which corresponds to the case $u_j=0$ and $w_j=0$ for all $j$. A standard bound on binomial coefficient for $1\leq k\leq n$ is:
\begin{equation*}
    {n \choose k} \leq \left(\frac{n\,e}{k}\right)^k \leq \left(n\,e\right)^k
\end{equation*}
Then
\begin{equation*}
    \mathcal{S} \;\;\leq\;\; \prod_{j=1}^b\left(1+\sum_{u_j=1}^{p_j-s_j}e^{-u_j\left(\delta \tfrac{b_j}{2}-1\right)}\right)\left(1+
    \sum_{w_j=1}^{s_j}    e^{-w_j\left(\delta \tfrac{c^{\prime}_j}{2}-1\right)}\right) - 1
\end{equation*}
Denote
$$
g_j\;=\;e^{1-\delta \tfrac{b_j}{2}},\qquad h_j\;=\;e^{1-\delta \tfrac{c^{\prime}_j}{2}}
$$
where both expressions go to zero as $n$ increases.
For every $j$, we have:
\begin{align*}
    &1+\sum_{u_j=1}^{p_j-s_j}e^{-u_j\left(\delta \tfrac{b_j}{2}-1\right)}
    =\frac{1-g_j^{p_j-s_j+1}}{1-g_j} \\
    &1+\sum_{w_j=1}^{s_j}    e^{-w_j\left(\delta \tfrac{c^{\prime}_j}{2}-1\right)} =\frac{1-h_j^{s_j+1}}{1-h_j}.
\end{align*}
Since both expressions converge to $1$ as $n$ grows, we conclude $\mathcal S\to 0$.

\subsection{Proof of Theorem~\ref{prop:fullrecovreg}}
From the proof of Theorem~\ref{theo:suffcondlinearmodel} we have:
\begin{equation}\label{eq:prodconstreg}
    \mathcal{S} \;\;\leq\;\; \prod_{j=1}^b\left(1+\sum_{u_j=1}^{p_j-s_j}e^{-u_j\left(\delta \tfrac{b_j}{2}-1\right)}\right)\left(1+
    \sum_{w_j=1}^{s_j}    e^{-w_j\left(\delta \tfrac{c^{\prime}_j}{2}-1\right)}\right) - 1
\end{equation}
For every $j$, we have:
\begin{align*}
    &\sum_{u_j=1}^{p_j-s_j}e^{-u_j\left(\delta \tfrac{b_j}{2}-1\right)}
    =g_j\,\frac{1-g_j^{p_j-s_j}}{1-g_j} \\
    &\sum_{w_j=1}^{s_j}    e^{-w_j\left(\delta \tfrac{c^{\prime}_j}{2}-1\right)} =h_j\,\frac{1-h_j^{s_j}}{1-h_j}.
\end{align*}
Each of the sum above is smaller than 1 for $n$ sufficiently large. Then every term of degree larger than 1 in the product in \eqref{eq:prodconstreg} is bounded above by the sum of terms of degree 1. For $n$ large enough, we get: 
 \begin{equation*} \mathcal{S} \leq (2^{2b}-2b)\sum_{j=1}^b \bigg[g_j\,\frac{1-g_j^{p_j-s_j}}{1-g_j}+ h_j\,\frac{1-h_j^{s_j}}{1-h_j}\bigg ]. \end{equation*} 
 Take $$r_b=\max_{j=1,\ldots,b}\bigg\{\frac{1-g_j^{p_j-s_j}}{1-g_j} \,,\,\frac{1-h_j^{s_j}}{1-h_j} \bigg\}.$$
 By Lemma~\ref{lem:L0toL1convergence} and observing that $2e \leq 6$, we get: 
 $$ P(\hat{S}^{b} \neq S)\;\leq  2 \sum_{M \in \M\setminus \{S\}} \E\left(NC(M)\right) \;\leq\ (2^{2b}-2b)6r_b\,\,\sum_{j=1}^b\,\, e^{-\tfrac{\delta}{2}\big[\kappa_j-\ln(p_j-s_j)\big]} + e^{-\tfrac{\delta}{2}\big[\frac{(1-\gamma) n\rho(\bs X)}{2}{{\beta_{\min}^{\star}}_j}^2 - \kappa_j -\ln(s_j) \big]} $$ and 
 $$ P(\hat{S}^{b} \neq S)\;\leq \;(2^{2b}-2b)6r_b\sum_{j=1}^b e^{-\tfrac{\delta}{2}\big[\kappa_j-\ln(p_j-s_j)\big]} + e^{-\frac{\delta}{2}\big[(\sqrt{\frac{(1-\gamma) n\rho(\bs X)}{2}}{{\beta_{\min}^{\star}}_j} - \sqrt{\kappa_j})^2 -\ln(s_j) \big]}.$$
     The bound in \eqref{eq:oraclerateconvreg} follows from proceeding as in the proof of Theorem~\ref{theo:fullrecovGS} with $c=1+\varepsilon$ such that $\lim \varepsilon > 0$. Under \eqref{betamincond:minimalreg}, for all $n\geq n_0$  we have $\sqrt{\kappa^*_j} \geq  (\frac{\varepsilon^2}{2(1+\varepsilon)}+1)\sqrt{\ln(p_j-s_j)}$ and $\frac{\sqrt{(1-\gamma) n\rho(\bs X)}}{4}{{\beta_{\min}^{\star}}_j}-\sqrt{\kappa^*} \geq (\frac{\varepsilon^2}{2(1+\varepsilon)}+1)\sqrt{\ln(s_j)}$
    which implies A6 and A7. Then there exists $r_b\to 1$ and $n_0$ such that for all $n\geq n_0$ \eqref{eq:rateconvreg} holds and proceeding as in the proof of Theorem~\ref{theo:fullrecovGS} shows that for any $j$ 
    $$e^{-\frac{\delta}{2}\left[{\kappa^*_j}-\ln(p_j-s_j)\right]}+e^{-\frac{\delta}{2}\big[(\sqrt{\frac{(1-\gamma) n\rho(\bs X)}{2}}{{\beta_{\min}^{\star}}_j} - \sqrt{\kappa_j^*})^2 -\ln(s_j) \big]} <2 e^{-\frac{\delta}{2}[\frac{(1-\gamma) n\rho(\bs X)}{32}{{\beta_{\min}^*}_j}^2-\max\{\ln(p_j-s_j),\ln(s_j)\}]}. $$

\subsection{Proof of Proposition~\ref{prop:nectypeIreg}}

Recovery of $S$ requires that for every $M\in O_j$, $NC(M)/NC(S)<1$, then it requires $\max_{M\in O_j}\frac{NC(M)}{NC(S)}<1$. By the definition of the normalized criterion $NC$:
$$P(\hat{S}^b=S) \leq P\Big(\max_{M\in O_j}\frac{NC(M)}{NC(S)}<1\Big)\leq  P\big(\max_{M\in O_j}\sqrt{L_{MS}}<\sqrt{2\kappa_j}\Big)$$
By Lemma~\ref{lemma:s7}, for every $M\in O_j$, $L_{MS}\in \chi^2_1$ and there exists $Z_M\sim N(0,1)$ such that $\sqrt{L_{MS}}=|Z_M|$. Since for every $M\in O_j$, $|Z_M|\geq Z_M$, we have:
\begin{equation}\label{eq:upperboundrecov}
    P(\hat{S}^b=S) \leq  1-P\big(\max_{M\in O_j}Z_M\geq\sqrt{2\kappa_j}\Big)
\end{equation}

The set $O_j$ has cardinality $p_j-s_j$, then by Theorem 3.4 in \cite{hartigan} and assumption A2, for any $\epsilon>0$:
$$P\big(\max_{M\in O_j}Z_M \geq \underline{\lambda}_j \sqrt{2 \ln(p_j-s_j)}(1-\epsilon)\big) \to 1$$
If $\lim_{n\to\infty}\frac{\kappa_j}{\underline{\lambda}_j^2\ln(p_j-s_j)} <1$, then $P\big(\max_{M\in O_j}Z_M\geq\sqrt{2\kappa_j}\big) \to 1$ and by \eqref{eq:upperboundrecov} $P(\hat{S}^b=S) \to 0$.

\subsection{Proof of Lemma~\ref{lem:upperboundnoncentral}}
The result follows directly from Lemma~\ref{lem:upperboundnoncentralgen}, with $M=S$ since $Q_T=S$.

\subsection{Proof of Proposition~\ref{prop:nectypeIIreg}}

Let $M=S\setminus S^S$, since $S^S\neq \emptyset$, $M\neq S$. The event $(\hat{S}^b=S)$ implies $\big(\frac{NC(M)}{NC(S)}<1\big)$ then $P(\hat{S}^b=S)\leq P\big(\frac{NC(M)}{NC(S)}<1\big)$. By definition of $NC$, we have
$P\big(\frac{NC(M)}{NC(S)}<1\big) = P(L_{SM}> 2\Delta_{SM})$ where by Lemma~\ref{lemma:s7},  $L_{SM}\sim \chi_{|S\setminus M|}^2(\mu_{SM})$. We  bound next $P(L_{SM}> 2\Delta_{SM})$ with Lemma~\ref{lemma:s1s3}. We first need to check that $2\Delta_{SM} > \mu_{SM} + |S\setminus M|$.
On one side, by Lemma~\ref{lem:upperboundnoncentral} we have:
$\mu_{SM} + |S\setminus M| \leq  \sum_{j=1}^b |S_j\setminus M_j|(n\bar{\lambda}\max_{i\in S_j\setminus M_j}{\beta_i^*}^2+1)$. On the other side, $2\Delta_{SM}=\sum_{j=1}^b |S_j\setminus M_j|2 \kappa_j$. Observe that $S\setminus M \subset S^S$, then by definition of $S^S$ and since $\kappa_j \to +\infty$ for all $j=1,\ldots,n$, for $n$ large enough, $2 \kappa_j >(n\bar{\lambda}\max_{i\in S_j\setminus M_j}|\beta_i^*|+1)$. We then have $2\Delta_{SM} > \mu_{SM} + |S\setminus M|$ and by Lemma~\ref{lemma:s1s3}:
\begin{equation}\label{eq:randomeqeee}
    P(L_{SM}> 2\Delta_{SM}) \leq e^{-\Delta_{SM}\left(1-\sqrt{\frac{\mu_{SM}}{2\Delta_{SM}}}\right)^2}\left(\frac{2\Delta_{SM}}{\mu_{SM}}\right)^{\frac{|S\setminus M|}{4}}
\end{equation}
By Lemma~\ref{lem:upperboundnoncentral} :
$$\frac{\mu_{SM}}{2\Delta_{SM}}\leq \frac{ \sum_{j=1}^b |S_j\setminus M_j|n\bar{\lambda}\max_{i\in S_j\setminus M_j}{\beta_i^*}^2}{\sum_{j=1}^b |S_j\setminus M_j|2 \kappa_j}$$
Observe that:
\begin{eqnarray*}
    \sum_{j=1}^b\frac{n\bar{\lambda}\max_{i\in S_j\setminus M_j}{\beta_i^{\star}}^2}{2\kappa_j}
    &\geq &\sum_{j=1,|S_j\setminus M_j|\neq 0}^b\frac{|S_j\setminus M_j|n\bar{\lambda}\max_{i\in S_j\setminus M_j}{\beta_i^{\star}}^2}{|S_j\setminus M_j|2\kappa_j} \\
    &=& \frac{\sum_{j=1,|S_j\setminus M_j|\neq 0 }^b|S_j\setminus M_j|n\bar{\lambda}\max_{i\in S_j\setminus M_j}{\beta_i^{\star}}^2\Big(1+\sum_{l\neq j}\frac{|S_l\setminus M_l|2\kappa_l}{|S_j\setminus M_j|2\kappa_j}\Big)}{\sum_{j=1}^b|S_j\setminus M_j|2\kappa_j} \\
    &\geq & \frac{\sum_{j=1}^b|S_j\setminus M_j|n\bar{\lambda}\max_{i\in S_j\setminus M_j}{\beta_i^{\star}}^2}{\sum_{j=1}^b \big|S\setminus M_j\big|2\kappa_j} 
\end{eqnarray*}
where last inequality follows from $\Big(1+\sum_{l\neq j}\frac{|S_l^S|\kappa_l}{|S_j^S|\kappa_j}\Big)\geq 1$ for all $j$ such that $|S_j^S|\neq 0$. We have $S_j\setminus M_j \subset S_j^S$ for every $j=1,\ldots,b$, then by definition of the $S_j^S$, $\sum_{j=1}^b\frac{n\bar{\lambda}\max_{i\in S_j\setminus M_j}{\beta_i^{\star}}^2}{2\kappa_j}$ goes to 0 and so does $\frac{\mu_{SM}}{2\Delta_{SM}}$. By \eqref{eq:randomeqeee}, for any $\phi \in (0,1)$, for $n$ large enough:
\begin{equation}
    P(L_{SM}> 2\Delta_{SM}) \leq e^{-\phi\Delta_{SM}}\left(\frac{2\Delta_{SM}}{\mu_{SM}}\right)^{\frac{|S\setminus M|}{4}}
\end{equation} 
By Lemma~\ref{lemma:s7}, for $n$ large enough $\mu_{SM} \geq 2$, then:
\begin{equation}
    P(L_{SM}> 2\Delta_{SM}) \leq e^{-\phi\Delta_{SM}}\left(\Delta_{SM}\right)^{\frac{|S\setminus M|}{4}}
\end{equation} 
and for any $\alpha\in(0,1)$, for $n$ large enough, we have:
\begin{equation}
    P(L_{SM}> 2\Delta_{SM}) \leq e^{-\alpha\phi\Delta_{SM}}
\end{equation} 
where the right hand-side goes to 0 since for all $j=1,\ldots,b$, $\kappa_j \to +\infty$. 

\subsection{Proof of Corollary~\ref{cor:necrecovreg}}
If for some $j\in\{1,\ldots,b\}$ $\lim_{n\to\infty}\frac{\sqrt{n\bar{\lambda}}\beta_{\min,j}^{\star}}{\sqrt{\underline{\lambda}_j\ln(p_j-s_j)}} = 0$ then for that $j$ we have:
$$\lim_{n\to\infty}\frac{\sqrt{n\bar{\lambda}}\beta_{\min,j}^{\star}}{\sqrt{\kappa_j}}\sqrt{\frac{\kappa_j}{\underline{\lambda}_j\ln(p_j-s_j)}} = 0$$
If $\lim_{n\to\infty}\frac{\sqrt{n\bar{\lambda}}\beta_{\min,j}^{\star}}{\sqrt{\kappa_j}}>0$, then $\lim_{n\to\infty}\sqrt{\frac{\kappa_j}{\underline{\lambda}_j\ln(p_j-s_j)}} = 0$ and by Proposition~\ref{prop:nectypeIreg} $P(\hat{S}^b=S)\not\to1$. If $\lim_{n\to\infty}\frac{\sqrt{n\bar{\lambda}}\beta_{\min,j}^{\star}}{\sqrt{\kappa_j}}=0$, then by Proposition~\ref{prop:nectypeIIreg} $P(\hat{S}^b=S)\not\to1$.


\subsection{Proof of Theorem~\ref{thm:convtoT}}
By definition of $\T$, for any $M \not\in\T$, there exists a unique $T^M\in\T$ such that $M\cap S^I=T^M\cap S^I$ and $M\neq T^M$. We next show that $A_{T^M}$ as defined in Lemma~\ref{lem:boundexpnc} satisfies the conditions of the lemma. Observe that $\Delta_{MT^M}$, defined in \eqref{eq:DeltaM}, can be rewritten as $\Delta_{MT^M}=\sum_{j=1}^b(|M_j\setminus T^M_j|-|T^M_j\setminus M_j|)\kappa_j$. By  Lemma~\ref{lem:boundrhogen}, for every $n\in \N$ we have 
\begin{equation}\label{eq:nonoverfittingcond2}
   \gamma \Delta_{MT^M}+\frac{1-\gamma}{2}\mu_{Q_{T^M}M}\;\geq\; \gamma\sum_{j=1}^b|M_j \setminus T^M_j|\kappa_j + \sum_{j=1}^b |T^M_j\setminus M_j |\left(\tfrac{1-\gamma}{2} n \rho(\bs X) \min_{i\in |T^M_j\setminus M_j|}{\beta^{\star}_i}^{2} - \gamma\kappa_j\right) 
\end{equation} 
Since $T^M \subseteq S^L \cup S^I$ and $M\cap S^I=T^M\cap S^I$, $T^M\setminus M \subseteq S^L$ and by definition of $S^L$, the rightmost component in \eqref{eq:nonoverfittingcond2} is nonnegative and positive if $|T^M\setminus M|\neq0$, for $n$ large enough. By assumption A6, component $\gamma\sum_{j=1}^b|M_j \setminus T^M_j|\kappa_j$ is nonnegative and positive if $|M\setminus T^M|\neq0$.  $M\neq T^M$ implies necessarily $|M\setminus T^M|\neq0$ or $|T^M\setminus M|\neq0$ and then, for $n$ large enough, $A_{T^M}>0$. If $|M\setminus T^M|=0$, we immediately have $|M\setminus T^M|=o(A_{T^M})$. If $|M\setminus T^M|\neq 0$, since rightmost component in \eqref{eq:nonoverfittingcond2} is nonnegative, we have:
$$\frac{|M\setminus T^M|}{A_{T^M}}\leq \bigg[ \gamma\sum_{j=1}^b\frac{|M_j \setminus T^M_j|}{|M\setminus T^M|}\kappa_j \bigg]^{-1} \leq \bigg[ \gamma\min_{j=1,\ldots,b}\kappa_j \bigg]^{-1} $$
where the last inequality follows from $\sum_{j=1}^b\frac{|M_j \setminus T^M_j|}{|M\setminus T^M|}=1$. By assumption A6, we also have $|M\setminus T^M|=o(A_{T^M})$ when $|M\setminus T^M|\neq 0$. 
We have $M\setminus T^M\subseteq V\setminus (S^I \cup S^L)$. If $M\cap S^S=\emptyset$, we immediately get $\mu_{Q_{T^M} T^M}=o(A_{T^M})$ because $\mu_{Q_{T^M} T^M}=0$. If $M\cap S^S\neq\emptyset$, by Lemma~\ref{lem:upperboundnoncentralgen}, we then have:
$$\frac{\mu_{Q_{T^M} T^M}}{ A_{T^M}} \leq \frac{n\,\bar{\lambda}\,\sum_{j=1}^b|(S_j\cap M_j )\setminus T^M_j|\,\max_{i\in (S_j\cap M_j )\setminus T^M_j}{\beta_{i}^{\star}}^2}{ \gamma\sum_{j=1}^b|M_j \setminus T^M_j|\kappa_j}$$
For all $j=1,\ldots,b$, $|M_j \setminus T^M_j| \geq |(S_j\cap M_j )\setminus T^M_j|$. Moreover, since  $M\setminus T^M\subseteq V\setminus (S^I \cup S^L)$ and $T^M\cap S^I=M\cap S^I$, $(S_j\cap M_j )\setminus T^M_j \subseteq S^S_j$ for all $j=1,\ldots,b$. We then get:
$$\frac{\mu_{Q_{T^M} T^M}}{ A_{T^M}} \leq \frac{n\,\bar{\lambda}\,\sum_{j=1}^b|(S_j\cap M_j )\setminus T^M_j|\,\max_{i\in S^S_j}{\beta_{i}^{\star}}^2}{ \gamma\sum_{j=1}^b|(S_j\cap M_j ) \setminus T^M_j|\kappa_j}$$
An upper bound on the right hand-side above follows from:
\begin{eqnarray*}
    \sum_{j=1}^b\frac{n\bar{\lambda}\max_{i\in S^S_j}{\beta_i^{\star}}^2}{\gamma\kappa_j}
    &\geq &\sum_{j=1,|(S_j\cap M_j )\setminus T^M_j|\neq 0}^b\frac{|(S_j\cap M_j )\setminus T^M_j|n\bar{\lambda}\max_{i\in S^S_j}{\beta_i^{\star}}^2}{|(S_j\cap M_j )\setminus T^M_j|\gamma\kappa_j} \\
    &=& \frac{\sum_{j=1,|(S_j\cap M_j )\setminus T^M_j|\neq 0 }^b|(S_j\cap M_j )\setminus T^M_j|n\bar{\lambda}\max_{i\in S^S_j}{\beta_i^{\star}}^2\Big(1+\sum_{l\neq j}\frac{|(S_l\cap M_l)\setminus T^M_l|\gamma\kappa_l}{|(S_j\cap M_j )\setminus T^M_j|\gamma\kappa_j}\Big)}{\sum_{j=1}^b|(S_j\cap M_j )\setminus T^M_j|\gamma\kappa_j} \\
    &\geq & \frac{\sum_{j=1}^b|(S_j\cap M_j )\setminus T^M_j|n\bar{\lambda}\max_{i\in S^S_j}{\beta_i^{\star}}^2}{\sum_{j=1}^b \big|(S\cap M_j)\setminus T^M_j\big|\gamma\kappa_j} 
\end{eqnarray*}
where last inequality follows from $\Big(1+\sum_{l\neq j}\frac{|(S_l\cap M_l)\setminus T^M_l|\gamma\kappa_l}{|(S_j\cap M_j )\setminus T^M_j|\gamma\kappa_j}\Big)\geq 1$ for all $j$. By definition of the $S_j^S$, $\sum_{j=1}^b n\bar{\lambda}\max_{i\in S^S_j}{\beta_i^{\star}}^2\,/\,\gamma\kappa_j$ goes to 0 and $\mu_{Q_{T^M} T^M}=o(A_{T^M})$ when $M\cap S^S\neq\emptyset$ too. When $|T^M\setminus M|\neq0$, by Lemma~\ref{lem:boundrhogen}
$$\frac{|T^M\setminus M|}{\mu_{Q_TM}} \leq \frac{|T\setminus M|}{n\,\rho(\bs X)\,\sum_{j=1}^b|T_j\setminus M_j|\,\min_{i\in T_j\setminus M_j}{\beta_{i}^{\star}}^2} $$
And since by definition of $T^M$, $|T^M\setminus M| = |S^L\setminus M|$:
$$\frac{|T^M\setminus M|}{\mu_{Q_TM}} \leq \frac{|S^L\setminus M|}{n\,\rho(\bs X)\,\sum_{j=1}^b|S^L_j\setminus M_j|\,\min_{i\in S^L_j\setminus M_j}{\beta_{i}^{\star}}^2}\leq \Big[n\,\rho(\bs X)\,\max_{i \in S^L}{\beta_{i}^{\star}}^2\Big]^{-1}  $$
where the last inequality follows from $\frac{|S^L\setminus M|}{\sum_{j=1}^b|S^L_j \setminus M_j|}=1$. By definition of $S^L$, we get have $|T^M\setminus M|=o(\mu_{Q_{T^M}M})$ when $|T^M\setminus M|\neq 0$.
We conclude that for every $M\in \mathcal M\setminus \T$, by Lemma~\ref{lem:boundexpnc}, for any $\psi\in(0,1)$ and $n$ large enough:
$$
\E(NC(M))\;\leq\;e^{-\psi A_{T^M}}.
$$
Let 
$$A_{T^M}^*\;:=\;\gamma\sum_{j=1}^b|M_j \setminus T^M_j|\kappa_j + \sum_{j=1}^b |T^M_j\setminus M_j |\left(\tfrac{1-\gamma}{2} n \rho(\bs X) \min_{i\in S^L_j}{\beta^{\star}_i}^{2} - \gamma\kappa_j\right) $$
By \eqref{eq:nonoverfittingcond} and because $\min_{i\in |T^M_j\setminus M_j|}{\beta^{\star}_i}^{2}\geq \min_{i\in S^L_j}{\beta^{\star}_i}^{2}$, we have for $n$ large enough:
\begin{equation}\label{eq:boundexp2}
    \E(NC(M))\;\leq\;e^{-\psi A_{T^M}^*}.
\end{equation}
The defintion of $S^L_j$ implies there exist $c^{\prime}_j \to +\infty$ such that:
$$\frac{(1-\gamma) n\rho(\bs X)}{2}\min_{i\in S^L_j}{\beta^{\star}_i}^{2}\,- \,\kappa_j \;=\; \ln(s_j) + c^{\prime}_{j}.$$
Let $\delta \in (0,1)$ and denote $\bar{m}_j=\max\big\{\frac{2\ln(p_j-s_j)}{b_{j}},\frac{2\ln(s_j)}{c^{\prime}_{j}}\big\}$. Take $\psi=\max_{j=1,\ldots,b}\frac{\xi + \delta + \bar{m}_j}{1+\bar{m}_j}$ for some $\xi \in(0, 1- \delta)$ then $\psi \in (0,1)$ and we have for every $j=1,\ldots,b$:
$$\delta \leq \psi-(1-\psi)\bar{m}_j.$$ 
By definition of $\gamma$ we have: 
$$
\gamma\kappa_j\;\geq\;\frac12 (1+\alpha_j)\kappa_j\;=\;\ln(p_j-s_j)+\frac12(\kappa_j-\ln(p_j-s_j)) \;=\; \ln(p_j-s_j)+ \frac12b_{j}.
$$
From our choice of $\psi,\delta$, we get:
$$\psi\gamma\kappa_j \;\geq\; \psi\big(\ln(p_j-s_j)+ \frac12b_{j}\big) \;\geq\;  \ln(p_j-s_j)+ \delta\frac12b_{j}.$$
By A7 and by our choice of $\psi,\delta$,
$$\psi\left(\tfrac{1-\gamma}{2} n \rho(\bs X) \min_{i\in S^L_j}{\beta^{\star}_i}^{2} - \gamma\kappa_j\right) \;\geq\; \psi\left(\ln(s_j)+c_j'\right) \;\geq\; \ln(s_j)+\delta \frac{c^{\prime}_j}{2}.
$$
Using the expression for $A_{T^M}^*$, the bound in \eqref{eq:boundexp2}, and the above inequalities, we get
$$
\E(NC(M))\;\leq\;\exp\{- \sum_{j=1}^b |M_j\setminus T^M_j|(\ln(p_j-s_j)+\delta \tfrac{b_j}{2})-\sum_{j=1}^b |T^M_j\setminus M_j|(\ln(s_j)+\delta \tfrac{c^{\prime}_j}{2})\}.
$$

Denote $\mathcal{S}= \sum_{M \in \M\setminus \T} \E\left(NC(M)\right)$ for convenience. Recall that for $M \in \M\setminus \T$ there exists a unique $T^M\in\T$ such that $M\cap S^I=T^M\cap S^I$.  We can split $\mathcal{S}$ into sums over the models $M$ such that, for all $j$, $|M_j\setminus T^M_j|=u_j$ and $|T^M_j\setminus M_j|=w_j$ with $u_j \in \{0,\ldots,p_j-s_j,\ldots,p_j-s_j+|S_j^S|\}$ and $w_j \in \{0,\ldots,|S^L_j|\}$, not all simultaneously zero:
$$\mathcal{S} =\sum_{w_1=0}^{|S^L_1|}\cdots\sum_{w_b=0}^{|S^L_b|}\sum_{u_1=0}^{p_1-s_1+|S_1^S|}\cdots\sum_{u_b=0}^{p_b-s_b+|S_b^S|}\sum_{M \neq T^M:\forall \,j |M_j\setminus T^M_j|=u_j, |T^M_j\setminus M_j|=w_j } \E\left(NC(M)\right)$$
The number of models with given collection of $u_j$'s and $w_j$'s is $\prod_{j=1}^b {p_j-s_j+|S_j^S| \choose u_j}{|S_j^L|\choose w_j}$. We thus have:
\begin{equation*}
    \mathcal{S} \leq \prod_{j=1}^b\left(1+\sum_{u_j=1}^{p_j-s_j+|S_j^S|}{p_j-s_j +|S_j^S|\choose u_j}e^{- u_j(\ln(p_j-s_j)+\delta \tfrac{b_j}{2})}\right)\left(1+
    \sum_{w_j=1}^{|S_j^L|}{|S_j^L| \choose w_j}
    e^{-w_j(\ln(s_j)+\delta \tfrac{c^{\prime}_j}{2})}\right) - 1
\end{equation*}
where $-1$ stems from discarding the cases $M=T^M$ ($u_j=0$ and $w_j=0$ for all $j$). A standard bound on binomial coefficient for $1\leq k\leq n$ is:
\begin{equation*}
    {n \choose k} \leq \left(\frac{n\,e}{k}\right)^k \leq \left(n\,e\right)^k
\end{equation*}
Then
\begin{equation*}
    \mathcal{S} \;\;\leq\;\; \prod_{j=1}^b\left(1+\sum_{u_j=1}^{p_j-s_j+|S_j^S|}e^{-u_j\left(\delta \tfrac{b_j}{2}-\ln\big(1+\frac{|S^S_j|}{p_j-s_j}\big)-1\right)}\right)\left(1+
    \sum_{w_j=1}^{|S_j^L|}    e^{-w_j\left(\ln\big(\frac{s_j}{|S_j^L|}\big)+\delta \tfrac{c^{\prime}_j}{2}-1\right)}\right) - 1
\end{equation*}
Denote
$$
g_j\;=\;e^{1+\ln\big(1+\frac{|S^S_j|}{p_j-s_j}\big)-\delta \tfrac{b_j}{2}},\qquad h_j\;=\;e^{1-\ln\big(\frac{s_j}{|S_j^L|}\big)-\delta \tfrac{c^{\prime}_j}{2}}
$$
Since $|S^S_j|=O(p_j-s_j)$ and $s_j\geq|S_j^L|$, $g_j$ and $h_j$ go to zero as $n$ increases.
For every $j$, we have:
\begin{align*}
    &1+\sum_{u_j=1}^{p_j-s_j+|S_j^S|}e^{-u_j\left(\delta \tfrac{b_j}{2}-1\right)}
    =\frac{1-g_j^{p_j-s_j+|S_j^S|+1}}{1-g_j} \\
    &1+\sum_{w_j=1}^{|S_j^L|}    e^{-w_j\left(\delta \tfrac{c^{\prime}_j}{2}-1\right)} =\frac{1-h_j^{|S_j^L|+1}}{1-h_j}.
\end{align*}
Since both expressions converge to $1$ as $n$ grows, we conclude $\mathcal S\to 0$ which proves (i). Finally, (ii) follows from (i) and observing that by Lemma~\ref{lem:L0toL1convergencegen}, 
$P(\hat{S}^{b} \not\in \T ) \,\,\leq \,\, (2^{|S^I|}+1) \sum_{M \in \M\setminus \{\T\}} \E\left(NC(M)\right) = (2^{|S^I|}+1) \mathcal{S}$.


\section{Proofs of Section~\ref{sec:informed_l0}}

% Then we have what I believed are tight bounds:
% \begin{equation}\label{eq:boundsshat}
%     \frac{s}{p}\,\,-\,\,B^{\circ} \leq \frac{\widehat{s}}{p}  \leq \frac{s}{p}\,\,+\,\, A^{\circ}
% \end{equation}
% You intuition was correct and you can ensure $\mathrm{E}\big[A^{\circ}]\to 0$ and $\mathrm{E}\big[\frac{\widehat{s}}{p}\big]\to\frac{s}{p}$ under milder condition than selection consistency. Setting $\kappa^{\circ}= \Theta(\ln(\frac{p-s}{\ln(p)})) + c$ where $c\to\infty$ is sufficient for $\mathrm{E}\big[A^{\circ}]\to 0$ while for consistency $\ln(p-s) + c$ where $c\to\infty$ is sufficient.

% The factor $[\ln(p)]^{-1}$ comes from the averaging by $p$. To have selection selection consistency we required (among other things) $\mathrm{E}\big[\sum_{M \in \M | |M| > S}NC^{\circ}(M)\big]\to 0$, for $\mathrm{E}[A^{\circ}]$ we have the same sum but with weights $\frac{|M|-s}{p}<1$. Proving something similar for $\frac{\widehat{s_j}}{p_j}$ does not seem straightforward though.\\

% Now, this does not solve the issue you had with a strict hypothesis in empirical Bayes. 
% Consider the empirical Bayes method , in step 2 we set penalty:
% $$\kappa_j = \ln(\frac{p_j}{\hat{s_j}}-1)+\frac{1}{2}\ln(n)$$
% For \textbf{consistency of the empirical Bayes} method we need for all $j=1,\ldots,b$:
% \begin{align}
%     \label{eq:typeI}&\kappa_j = \ln(p_j-s_j)+c_j\\ 
%     \label{eq:typeII}&g(n)\beta_{\min,j}^*-\sqrt{\kappa_j} = \sqrt{\ln(s_j)} + d_j 
% \end{align}
% where $c_j,d_j\to\infty$ and $g(n)=\sqrt\frac{{(1-\gamma)n\rho(\bs X)}}{8}$
% We have:
% $$\ln(\frac{p_j}{\hat{s_j}}-1) = \ln(\frac{p_j}{s_j}-1)+\ln(\frac{p_j-\hat{s_j}}{p_j-s_j}) - \ln(\frac{\hat{s_j}}{s_j})$$

\subsection{Proof of Lemma~\ref{lem:empiricalbayesinterpret}}
Denote $H(\bs \theta)=p(\bs y \mid \bs \theta)$. For any $j=1,\ldots,b$, the partial derivative with respect to $\theta^{(j)}$ is:
    \begin{equation}\label{eq:logderivative}
        \frac{\partial \ln H(\bs \theta)}{\partial \theta^{(j)}}=\frac{\partial H(\bs \theta)}{\partial \theta^{(j)}}\,H(\boldsymbol{\theta})^{-1}
    \end{equation}
    Observe that:
    \begin{equation}\label{eq:partialmarglikelihood}
    H(\bs \theta)=\sum_{M \in \M}p(\bs y \mid M)p(M \mid \bs \theta)
    \quad \text{and}\quad 
    \frac{\partial H(\bs \theta)}{\partial \theta^{(j)}}=\sum_{M\in \M}p(\bs y \mid M)\frac{\partial p(M \mid \bs \theta)}{\partial \theta^{(j)}}
    \end{equation}
    By our choice of prior~\eqref{eq:bernoulliprior}:
    \begin{equation}\label{eq:partialprobmodel}
    %p(M \mid \bs \theta)=\prod_{j=1}^b \theta^{(j)}^{\sum_{i\in B_j}m_i} \, (1-\theta^{(j)})^{p_j - \sum_{i\in B_j}m_i}\quad\text{and}\quad 
    \frac{\partial p(M \mid \bs \theta)}{\partial \theta^{(j)}}= p(M \mid \bs \theta) \bigg(\frac{\sum_{i\in B_j}m_i}{\theta^{(j)}}-\frac{p_j-\sum_{i\in B_j}m_i}{1-\theta^{(j)}} \bigg)        
    \end{equation}
    Replacing \eqref{eq:partialprobmodel} into \eqref{eq:partialmarglikelihood}, we get:
    \begin{equation*}\label{eq:partialmarglikelihood2}
        \frac{\partial H(\bs \theta)}{\partial \theta^{(j)}}=\frac{1}{\theta^{(j)}}\sum_{i \in B_j}\sum_{M\in \M: m_i=1}p(\bs y \mid M)p(M \mid \bs \theta) - \frac{1}{1-\theta^{(j)}}\bigg[p_j H(\bs \theta)-\sum_{i \in B_j}\sum_{M\in \M: m_i=1}p(\bs y \mid M)p(M \mid \bs \theta)
        \bigg]
    \end{equation*}
    Note that:
    $$\frac{\sum_{M\in \M: m_i=1}p(\bs y \mid M)p(M \mid \bs \theta)}{\sum_{M\in \M}p(\bs y \mid M)p(M \mid \bs \theta)} = P(m_i = 1\mid \bs y)$$
    By \eqref{eq:logderivative}, we then get the following expression of the partial derivative for any $j=1,\ldots,b$:
    $$\frac{\partial \ln H(\bs \theta)}{\partial \theta^{(j)}}=\frac{1}{\theta^{(j)}}\sum_{i \in B_j}P(m_i = 1\mid \bs y) - \frac{1}{1-\theta^{(j)}}\bigg[p_j-\sum_{i \in B_j}P(m_i = 1\mid \bs y)
        \bigg]$$
    Setting the partial derivatives to 0 and solving for $\theta^{(j)}$ gives the result.

\subsection{Proof of Proposition~\ref{lem:shatconsistence}}
% Observe that
%     \begin{equation}\label{eq1}
% \begin{cases}
%     \text{if } \, \beta_i^{\star} \neq 0, \,\, P(\hat{M}_j = S_j \mid \bs y) \leq P(m_i = 1\mid \bs y)\leq 1 \\
%     \text{if } \, \beta_i^{\star} = 0, \,\, 0\leq P(m_i = 1\mid\bs y) \leq P(\hat{M}_j \neq S_j \mid \bs y) \\
% \end{cases}
% \end{equation}
% We also have that:
% $$\widehat{s}_j=\sum_{i  \in B_j:\beta_i^{\star}\neq 0}P(\beta_i \neq 0\mid\bs y)+\sum_{i  \in B_j:\beta_i^{\star}=0}P(\beta_i \neq 0\mid\bs y)$$
% From the inequalities \eqref{eq1}, we get:
% \begin{align*}
%     &s_j\big(1-P(\hat{M}_j \neq S_j\mid\bs y)\big)\leq \sum_{\beta_i^{\star} \neq 0}P(\beta_i \neq 0\mid\bs y) \leq s_j
% \end{align*}
% and,
% \begin{align*}
%     &0 \leq \sum_{\beta_i^{\star} = 0}P(\beta_i \neq 0\mid\bs y) \leq \big(p_j-s_j\big)P(\hat{M}_j \neq S_j\mid\bs y)
% \end{align*}
% Then:
% $$s_j\big(1-P(\hat{M}_j \neq S_j\mid\bs y)\big)\leq \widehat{s}_j\leq s_j+\big(p_j-s_j\big)P(\hat{M}_j \neq S_j\mid\bs y)$$
% Dividing by $p_j$ and resting $s_j/p_j$, we have:
% $$-\frac{s_j}{p_j}P(\hat{M}_j \neq S_j\mid\bs y)\leq \frac{\widehat{s}_j}{p_j}-\frac{s_j}{p_j}\leq\big(1-\frac{s_j}{p_j}\big)P(\hat{M}_j \neq S_j\mid\bs y)$$
% And:
% $$\Big|\frac{\widehat{s}_j}{p_j}-\frac{s_j}{p_j}\Big|\leq\max\Big\{1-\frac{s_j}{p_j},\frac{s_j}{p_j}\Big\}P(\hat{M}_j \neq S_j\mid\bs y) \leq P(\hat{M}_j \neq S_j\mid\bs y)$$
% %$\hat{M}$ is assumed to be a consistent estimator of $S$ then $P(\hat{M} \neq S\mid\bs y) \to 0$ and $\frac{\widehat{s}_j}{p_j} \to s_j/p_j$.
% \\
% \\
% -------------------------------------\\
% Observe that for any $j=1,\ldots,b$ and any $i\in B_j$
%     \begin{equation}\label{eq1}
% \begin{cases}
%     \text{if } \, \beta_i^{\star} \neq 0, \,\,  \sum_{M \in \M| S_j \subseteq M_j} NC(M) \leq \sum_{M \in \M| i\in M}NC(M) \leq 1 \\
%     \text{if } \, \beta_i^{\star} = 0, \,\, 0\leq \sum_{M \in \M| i\in M}NC(M)\leq \sum_{M \in \M| M_j \not\subseteq S_j} NC(M) \\
% \end{cases}
% \end{equation}
% We also have that:
% $$\frac{\widehat{s}_j}{p_j}=\frac{1}{p_j}\sum_{i  \in B_j:\beta_i^{\star}\neq 0}\sum_{M \in \M| i\in M}NC(M)
% +\frac{1}{p_j}\sum_{i  \in B_j:\beta_i^{\star}=0}\sum_{M \in \M| i\in M}NC(M)$$
% From the inequalities \eqref{eq1}, we get:
% \begin{align*}
%     &\frac{s_j}{p_j}\big(1-\sum_{M \in \M| S_j \not\subseteq M_j} NC(M) \big)\leq \frac{1}{p_j}\sum_{i  \in B_j:\beta_i^{\star}\neq 0}\sum_{M \in \M| i\in M}NC(M) \leq  \frac{s_j}{p_j}
% \end{align*}
% and,
% \begin{align*}
%     &0 \leq \frac{1}{p_j} \sum_{i  \in B_j:\beta_i^{\star}= 0}\sum_{M \in \M| i\in M}NC(M) \leq \big(1-\frac{s_j}{p_j}\big)\sum_{M \in \M| M_j \not\subseteq S_j} NC(M)
% \end{align*}
% Then:
% $$\frac{s_j}{p_j}\big(1-\sum_{M \in \M| S_j \not\subseteq M_j} NC(M)\big) \leq \frac{\widehat{s}_j}{p_j}\leq \frac{s_j}{p_j}+\big(1-\frac{s_j}{p_j}\big)\sum_{M \in \M| M_j \not\subseteq S_j} NC(M)$$
% And:
% $$\Big|\frac{\widehat{s}_j}{p_j}-\frac{s_j}{p_j}\Big|\leq\max\Bigg\{\big(1-\frac{s_j}{p_j}\big) \sum_{M \in \M| M_j \not\subseteq S_j} NC(M),\frac{s_j}{p_j}\sum_{M \in \M| S_j \not\subseteq M_j} NC(M)\Bigg\} \leq \sum_{M \in \M| M_j \neq S_j} NC(M)$$

% ---------------------------------------\\
% Observe that for any $j=1,\ldots,b$ and any $i\in B_j$:
%     \begin{equation}\label{eq1}
% \begin{cases}
%     \text{if } \, \beta_i^{\star} \neq 0, \,\,  P( S_j\subseteq \widehat{M}_j) \leq P(\beta_i \neq 0) \leq 1 \\
%     \text{if } \, \beta_i^{\star} = 0, \,\, 0\leq P(\beta_i \neq 0)\leq P(\widehat{M}_j \not\subseteq S_j) \\
% \end{cases}
% \end{equation}
% We also have that:
% $$\frac{\widehat{s}_j}{p_j}=\frac{1}{p_j}\sum_{i  \in B_j:\beta_i^{\star}\neq 0}P(\beta_i \neq 0)
% +\frac{1}{p_j}\sum_{i  \in B_j:\beta_i^{\star}=0}P(\beta_i \neq 0)$$
% From the inequalities \eqref{eq1}, we get:
% \begin{align*}
%     &\frac{s_j}{p_j}\big(1-P( S_j\not\subseteq \widehat{M}_j) \big)\leq \frac{1}{p_j}\sum_{i  \in B_j:\beta_i^{\star}\neq 0}P(\beta_i \neq 0) \leq  \frac{s_j}{p_j}
% \end{align*}
% and,
% \begin{align*}
%     &0 \leq \frac{1}{p_j}\sum_{i  \in B_j:\beta_i^{\star}= 0}P(\beta_i \neq 0) \leq \big(1-\frac{s_j}{p_j}\big)P(\widehat{M}_j \not\subseteq S_j)
% \end{align*}
% Then:
% $$\frac{s_j}{p_j}\big(1-P( S_j\not\subseteq \widehat{M}_j)\big) \leq \frac{\widehat{s}_j}{p_j}\leq \frac{s_j}{p_j}+\big(1-\frac{s_j}{p_j}\big)P(\widehat{M}_j \not\subseteq S_j)$$
% And:
% $$\Big|\frac{\widehat{s}_j}{p_j}-\frac{s_j}{p_j}\Big|\leq\max\Bigg\{\big(1-\frac{s_j}{p_j}\big) P(\widehat{M}_j \not\subseteq S_j),\frac{s_j}{p_j}P( S_j\not\subseteq \widehat{M}_j)\Bigg\} \leq P(\widehat{M}_j \neq S_j)$$


For every $j=1,\ldots,b$, we have the decomposition:
\begin{gather}\label{eq:decompshat}
    \E\bigg(\frac{\widehat{s_j}}{p_j}\bigg)
   =A_j\,+\, B_j \\
       \text{where} \quad A_j=\frac{1}{p_j}\sum_{i \in B_j}\sum_{M \in \T|i\in M}\E(NC(M)) \quad\text{and}\quad  B_j=\frac{1}{p_j}\sum_{i \in B_j}\sum_{M \in \M\setminus \T|i\in M}\E(NC(M)) \nonumber
\end{gather}
A bound on $B_j$ for every $j=1,\ldots,b$ is:
$$B_j = \sum_{M \in \M\setminus \T}\E(NC(M))\sum_{i \in B_j}\frac{I(i\in M_j)}{p_j}\leq \sum_{M \in \M\setminus \T}\E(NC(M)) $$
where the inequality follows from $\sum_{i \in B_j}\frac{I(i\in M_j)}{p_j}\leq 1$ for all $M$. By Theorem~\ref{thm:convtoT}, for every $j=1,\ldots,b$, $B_j$ vanishes.
For every $j=1,\ldots,b$, $A_j$ can be furthered decomposed:
\begin{align*}
  A_j
   &= \sum_{M \in \T}\E(NC(M))\sum_{i \in B_j}\frac{I(i\in M_j)}{p_j}\\
   &= \sum_{M \in \T}\E(NC(M))\sum_{i \in S^L_j}\frac{I(i\in M_j)}{p_j}+\sum_{M \in \T}\E(NC(M))\sum_{i \in B_j\setminus S^L_j}\frac{I(i\in M_j)}{p_j}\\
   &= \frac{|S^L_j|}{p_j}\sum_{M \in \T}\E(NC(M))+\sum_{M \in \T}\E(NC(M))\sum_{i \in B_j\setminus S^L_j}\frac{I(i\in M_j)}{p_j}
\end{align*}
where the last equality follows from $I(i\in M_j)=1$ for all $i\in S^L$ and $M\in\T$. By  Theorem~\ref{thm:convtoT}, for every $j=1,\ldots,b$, $\lim_{n\to\infty}A_j \geq \frac{|S^L_j|}{p_j}$. By \eqref{eq:decompshat}, since $B_j$ vanishes, $\lim_{n\to\infty}\E\Big(\frac{\widehat{s_j}}{p_j}\Big)\geq \frac{|S^L_j|}{p_j}$ for every $j=1,\ldots,b$.

Since $I(i\in M_j)=0$ for all $i\in B_j\setminus(S^L_j\cup S^I_j)$ and $M\in\T$, $A_j$ can be further written:
\begin{align*}
  A_j
   &= \frac{|S^L_j|}{p_j}\sum_{M \in \T}\E(NC(M))+\sum_{M \in \T}\E(NC(M))\sum_{i \in S^I_j}\frac{I(i\in M_j)}{p_j}
\end{align*}
Using that $\sum_{i \in S^I_j}\frac{I(i\in M_j)}{p_j}\leq \frac{|S^I_j|}{p_j}$ for all $M$, we have:
\begin{align*}
    \frac{s_j}{p_j}-\E\bigg(\frac{\widehat{s_j}}{p_j}\bigg) = \frac{s_j}{p_j}-A_j - B_j
    &\geq \frac{s_j}{p_j}- \frac{|S^L_j|+|S^I_j|}{p_j}\sum_{M \in \T}\E(NC(M)) - B_j
\end{align*}
By Theorem~\ref{thm:convtoT}, $\sum_{M \in \T}\E(NC(M))\to 1$ and $B_j\to 0$, and then $ \frac{s_j}{p_j}-\lim_{n\to+\infty}\E\big(\frac{\widehat{s_j}}{p_j}\big) \geq  \frac{s_j-|S^L_j|+|S^I_j|}{p_j}= \frac{|S^S_j|}{p_j}$ for every $j=1,\ldots,b$, which proves the upper bound.


\subsection{Proof of Theorem \ref{theo:empbayesselconsist}}
Denote for any $M\in\M$, $NC^{\circ}(M)$, the normalized criterion value for model $M$ under step 1 penalty $\kappa^{\circ}$. For this choice of penalty and every $j=1,\ldots,b$, we have:
\begin{equation*}
    \frac{\widehat{s_j}}{p_j}
   = \frac{1}{p_j}\sum_{M \in \M}NC^{\circ}(M)\sum_{i \in B_j}I(i\in M) =\sum_{M \in \M | M \neq S}\frac{|M_j|}{p_j}NC^{\circ}(M) +\frac{s_j}{p_j}NC^{\circ}(S) 
\end{equation*}
Using that $NC^{\circ}(S)=1-\sum_{M \in \M | M \neq S}NC(M)$, we get
\begin{equation*}
   \frac{\widehat{s}_j}{p_j}
   =\sum_{M \in \M | M \neq S}\frac{|M_j|-s_j}{p_j}NC^{\circ}(M) \,\,+\,\, \frac{s_j}{p_j}
\end{equation*}
Splitting the sum in the right-hand term above between models $M$ that are overfitting in $B_j$, that is $|M_j|\geq s_j$ and those that are underfitting in $B_j$, $|M_j|<s_j$, we can write:
\begin{equation}\label{eq:rewriteshat}
  \frac{\widehat{s}_j}{p_j}  = \frac{s_j}{p_j}\,\,+\,\, O_j^{\circ}-\,\,U_j^{\circ}  
\end{equation}
where:
$$
    O_j^{\circ} = \sum_{M \in \M | |M_j| > s_j}\frac{|M_j|-s_j}{p_j}NC^{\circ}(M)\quad\text{and}\quad
    U_j^{\circ} = \sum_{M \in \M | |M_j| < s_j}\frac{s_j-|M_j|}{p_j}NC^{\circ}(M)
$$
By \eqref{eq:rewriteshat}, we have the following decomposition of step 2 penalties :
\begin{align}\label{eq:decompstep2pen}
    \kappa_j
    &= \ln(p_j-s_j) + \ln\Big(\frac{\sqrt{n}}{s_j}\Big) + \ln\Big(\frac{p_j-\widehat{s}_j}{p_j-s_j}\Big) - \ln\Big(\frac{\widehat{s}_j}{s_j}\Big) \nonumber\\
    &= \ln(p_j-s_j)+\ln(\frac{\sqrt{n}}{s_j})+\ln\Big(1-\frac{p_j(O^{\circ}_j-U^{\circ}_j)}{p_j-s_j}\Big)-\ln\Big(\frac{\widehat{s}_j}{s_j}\Big)
\end{align}
We now show the $\kappa_j$ satisfy assumption~A6. Since $U_j^{\circ}$ is non negative, a lower bound on $\kappa_j$ for every $j=1,\ldots,b$ is:
\begin{equation*}
    \kappa_j
    \geq  \ln(p_j-s_j)+\ln(\frac{\sqrt{n}}{s_j})+\ln\Big(1-\frac{p_j O^{\circ}_j}{p_j-s_j}\Big)-\ln\Big(\frac{\widehat{s}_j}{s_j}\Big)
\end{equation*}
We have :
$$\frac{p_j O^{\circ}_j}{p_j-s_j} =  \sum_{M \in \M | |M_j| > s_j}\frac{|M_j|-s_j}{p_j-s_j}NC^{\circ}(M)\leq \sum_{M \in \M | |M_j| > s_j}NC^{\circ}(M)$$ 
If $M$ is such that $|M_j|>s_j$, then $M\not\in\T$ and by Theorem~\ref{thm:convtoT}, $\ln\Big(1-\frac{p_j O^{\circ}_j}{p_j-s_j}\Big)$ vanishes in probability. 
Let $\varepsilon_n$ be a positive increasing sequence such that $\varepsilon_n=o(\sqrt{n}s^{-1})$. By the Markov inequality:
$$P\Big(-\ln\Big(\frac{\widehat{s}_j}{s_j}\Big)>-\ln(\varepsilon_n)\Big)\geq 1-\frac{\E(\widehat{s}_j/p_j)}{s_j/p_j} \varepsilon_n^{-1} $$
By Proposition~\ref{lem:shatconsistence}, we have
$$\lim_{n\to\infty}P\Big(-\ln\Big(\frac{\widehat{s}_j}{s_j}\Big)>-\ln(\varepsilon_n)\Big)\geq 1-\lim_{n\to\infty}\frac{s_j-|S^S_j(\kappa^{\circ})|}{s_j} \varepsilon_n^{-1} \geq 1-\lim_{n\to\infty}\varepsilon_n^{-1}=1  $$
We then get that with probability going to 1:
\begin{equation}\label{eq:finallowerboundkappaj}
  \kappa_j 
    \geq  \ln(p_j-s_j)+\ln\Big(\frac{\sqrt{n}}{s_j\varepsilon_n}\Big) 
\end{equation}
Since by assumption $s_j=o(\sqrt{n})$ and $\varepsilon_n=o(\sqrt{n}s_j^{-1})$, the $\kappa_j$ satisfy assumption~A6 with probability going to 1.

We go on with showing that assumption~A9 implies A7 holds. Observe:
$$\ln\Big(\frac{\widehat{s}_j}{s_j}\Big)=\ln\Big(1+\frac{p_j(O^{\circ}_j-U^{\circ}_j)}{s_j}\Big) $$
Since $O_j^{\circ}$ is nonnegative, from \eqref{eq:decompstep2pen}, for every $j=1,\ldots,b$, an upper bound on $\kappa_j$ is:
\begin{equation}\label{eq:gggggggg}
    \kappa_j \leq \ln\big(p_j-s_j\big)+\ln\Big(\frac{\sqrt{n}}{s_j}\Big)+\ln\bigg(\frac{1+\frac{p_j}{p_j-s_j}U^{\circ}_j}{1-\frac{p_j}{s_j}U^{\circ}_j}\bigg)
\end{equation}
We have:
\begin{align*}
    \frac{p_j}{p_j-s_j}U^{\circ}_j = \sum_{M \in \M | |M_j| < s_j}\frac{s_j-|M_j|}{p_j-s_j}NC^{\circ}(M) \\
    \frac{p_j}{s_j}U^{\circ}_j = \frac{p_j-s_j}{s_j} \sum_{M \in \M | |M_j| < s_j}\frac{s_j-|M_j|}{p_j-s_j}NC^{\circ}(M)
\end{align*}
We split the sum in the right-hand terms above between models $M\in \T$ and $M\not\in \T$. If $M\in \T$, then by definition $s_j-|M_j|\leq s_j-|S^L_j(\kappa^{\circ})|$. Using that $ \sum_{M \in \T}NC^{\circ}(M)\leq 1$ we have: 
$$ \sum_{M \in \T | |M_j| < s_j}\frac{s_j-|M_j|}{p_j-s_j}NC^{\circ}(M) \leq \frac{s_j-|S^L_j(\kappa^{\circ})|}{p_j-s_j}$$
For the $M\not\in \T$, using that $s_j-|M_j|\leq\max\{p_j-s_j,s_j\}$:
$$\sum_{M \in \M\setminus \T | |M_j| < s_j}\frac{s_j-|M_j|}{p_j-s_j}NC^{\circ}(M) \leq \frac{s_j-|S^L_j(\kappa^{\circ})|}{p_j-s_j}\sum_{M \in \M\setminus \T }NC^{\circ}(M)$$
By assumption, $\frac{s_j-|S^L_j(\kappa^{\circ})|}{p_j-s_j}$ is bounded so by Theorem~\ref{thm:convtoT}, with probability going to 1:
\begin{align*}
    \frac{p_j}{p_j-s_j}U^{\circ}_j \leq \frac{s_j-|S^L_j(\kappa^{\circ})|}{p_j-s_j}\\
    \frac{p_j}{s_j}U^{\circ}_j \leq  \frac{s_j-|S^L_j(\kappa^{\circ})|}{s_j}
\end{align*}

By replacing $\frac{p_j}{p_j-s_j}U^{\circ}_j$ and $\frac{p_j}{s_j}U^{\circ}_j$ in \eqref{eq:gggggggg} by the bounds above, we get that  with probability going to 1:
\begin{equation}\label{eq:finalupperboundkappaj}
  \kappa_j \leq \ln\big(p_j-|S^L_j(\kappa^{\circ})|\big)+\ln\Big(\frac{\sqrt{n}}{s_j}\Big)
\end{equation}
Moreover the lower bound \eqref{eq:finallowerboundkappaj}, holds for any $\varepsilon_n=o(\sqrt{n}s^{-1})$. Taking $\varepsilon_n=\ln(\sqrt{n}s^{-1})$, we have with probability going to 1:
\begin{equation}\label{eq:rrrrrrrrr}
    1-\frac12\Big(1+\max_{j=1, \ldots,b}\frac{\ln(p_j-s_j)}{\kappa_j}\Big) \geq 1-\psi
\end{equation}
By \eqref{eq:finalupperboundkappaj} and \eqref{eq:rrrrrrrrr} assumption~A9 is sufficient for assumption~A7 to hold. Since A6 and A7 hold, by Theorem \ref{theo:suffcondlinearmodel}, $P(\hat{S}^{EB}\neq S) \to 0$.

\subsection{Proof of Theorem \ref{theo:informedl0consist}}
Proceeding as in the proof of Theorem ~\ref{theo:empbayesselconsist}, we have the following decomposition of step 2 penalties :
\begin{align*}
    \kappa_j
    &= \ln(p_j-s_j) + \ln(\sqrt{n}) + \ln\Big(\frac{p_j-\widehat{s}_j}{p_j-s_j}\Big)  \nonumber\\
    &= \ln(p_j-s_j)+\ln(\sqrt{n})+\ln\Big(1-\frac{p_j(O^{\circ}_j-U^{\circ}_j)}{p_j-s_j}\Big)
\end{align*}
Since $U_j^{\circ}$ is non negative, lower bound on those is:
\begin{equation*}
    \kappa_j
    \geq  \ln(p_j-s_j)+\ln(\sqrt{n})+\ln\Big(1-\frac{p_j O^{\circ}_j}{p_j-s_j}\Big)
\end{equation*}
From the proof of Theorem ~\ref{theo:empbayesselconsist}, by Theorem~\ref{thm:convtoT}, $\ln\Big(1-\frac{p_j O^{\circ}_j}{p_j-s_j}\Big)$ vanishes in probability. With probability going to 1, we then have:
\begin{equation}\label{eq:finallowerboundkappaj2}\kappa_j 
    \geq  \ln(p_j-s_j)+\ln(\sqrt{n})
\end{equation}
and the $\kappa_j$ satisfy assumption~A6.

Since $O_j^{\circ}$ is nonnegative, from \eqref{eq:decompstep2pen}, for every $j=1,\ldots,b$, an upper bound on $\kappa_j$ is:
\begin{equation*}
    \kappa_j \leq \ln\big(p_j-s_j\big)+\ln(\sqrt{n})+\ln\bigg(1+\frac{p_j}{p_j-s_j}U^{\circ}_j\bigg)
\end{equation*}
As in the proof of Theorem ~\ref{theo:empbayesselconsist}, by Theorem~\ref{thm:convtoT}, with probability going to 1:
\begin{equation*}
    \frac{p_j}{p_j-s_j}U^{\circ}_j \leq \frac{s_j-|S^L_j(\kappa^{\circ})|}{p_j-s_j}
\end{equation*}
and, for any $\delta >0$:
\begin{equation}\label{eq:finalupperboundkappaj2}
\kappa_j \leq \ln\big(p_j-|S^L_j(\kappa^{\circ})|\big)+\ln(\sqrt{n})
\end{equation}
Moreover by \eqref{eq:finallowerboundkappaj2}, we have with probability going to 1:
\begin{equation}\label{eq:rrrrrrrrr2}
    1-\frac12\Big(1+\max_{j=1, \ldots,b}\frac{\ln(p_j-s_j)}{\kappa_j}\Big) \geq 1-\psi
\end{equation}
By \eqref{eq:finalupperboundkappaj2} and \eqref{eq:rrrrrrrrr2} assumption~A10 is sufficient for assumption~A7 to hold. Since A6 and A7 hold, by Theorem \ref{theo:suffcondlinearmodel}, $P(\hat{S}^{I\ell_0}\neq S) \to 0$.

\section{Proofs of auxiliary results}\label{sec:proofauxi}

\subsection{Proof of Lemma~\ref{lemma:3tailbounds}}

	a) \,\, By union bound and the identical distribution of the $y_i$'s,
\begin{center}
    $P\left(\underset{i\in\{1,\ldots,p\}}{\max}|y_i|> a \right)     \leq \sum_{i=1}^{p}{ P\left(|\sqrt{n}y_i|>\sqrt{n}a \right)}
    = p\; P\left(|z|U\sqrt{n}a\right)$
\end{center}
    where $z\sim N(0,1)$. 
    By symmetry and the standard tail bound for standard normal, $P(z \geq \delta) \leq \frac{1}{\sqrt{2\pi}} \frac{1}{\delta} e^{-\delta^{2} / 2}$:
\begin{center}
  $P\left(\underset{i\in\{1,\ldots,p\}}{\max}|y_i|> a \right) 
     \leq  p\; 2 \, P\left(z>\sqrt{n}a\right) \leq \frac{1}{\sqrt{\pi}}\frac{\sqrt{2}}{\sqrt{n}a}p e^{-\frac{n}{2} a^{2}}$
\end{center}
Since $a \geq \sqrt{\frac{2\ln(p)}{n}}$, $\frac{\sqrt{2}}{\sqrt{n}a}\leq \frac{1}{\ln(p)}$. Taking $a^2=a^2-\frac{2\ln(p)}{n}+\frac{2\ln(p)}{n}$, we get \begin{center}
$ P\left(\underset{i\in\{1,\ldots,p\}}{\max}|y_i|> a \right) \leq \frac{1}{\sqrt{\pi\ln(p)}} e^{-\frac{n}{2} \left(a^{2}-\frac{2\ln(p)}{n}\right)}$ \\\end{center}
b) $ \underset{i \in \{1,\ldots,s\}}{\min}|\beta_i|- \underset{i \in \{1,\ldots,s\}}{\max}|y_i-\beta_i| > a$ implies $\forall i \; \underset{i \in \{1,\ldots,s\}}{\min}|\beta_i| - |y_i-\beta_i| > a $. Since $\forall i \; |\beta_i|< |\beta_i-y_i|+|y_i|$ we have that $ \forall i \; \min_{i \in \{1,\ldots,s\}}|\beta_i| - |y_i-\beta_i| > a$ implies $\forall i \; |y_i| > a $ and $ \underset{i \in \{1,\ldots,s\}}{\min}|y_i| > a$. By the properties of the probability $P$: 
\begin{center}$ P\left(\underset{i \in \{1,\ldots,s\}}{\min}|\beta_i|-\underset{i \in \{1,\ldots,s\}}{\max}|y_i-\beta_i| > a\right) \leq P\left(\underset{i \in \{1,\ldots,s\}}{\min}|y_i| > a\right) $ \end{center}
c)\,\, We start by bounding the probability $P(|y_i|)>a$ for any $i$. We have that:
$$P(|y_i|>a) = P(y_i>a) + P(y_i<-a) = P(z>\frac{a-\mu}{\sigma}) + P(z<-\frac{a+\mu}{\sigma}) \quad \text{where }z\sim N(0,1)$$
If $\mu>0$, then $P(|y_i|>a) =  P(z>\frac{a-|\mu|}{\sigma}) + P(z<-\frac{a+|\mu|}{\sigma})$. If $\mu<0$, by symmetry, $P(z>\frac{a-\mu}{\sigma})=P(z>\frac{a+|\mu|}{\sigma})=P(z<-\frac{a+|\mu|}{\sigma})$ and $P(z<-\frac{a+\mu}{\sigma})=P(z<\frac{|\mu|-a}{\sigma})=P(z>\frac{a-|\mu|}{\sigma})$. Then for any $\mu$, $P(|y_i|>a)  = P(z>\frac{a-|\mu|}{\sigma}) + P(z<-\frac{a+|\mu|}{\sigma}) $ and using the symmetry of the standard Gaussian for any $a<|\mu|$: 
\begin{equation}\label{eq:random}
    P(|y_i|>a) = 1+P(0<z<\frac{|\mu|-a}{\sigma})-P(0<z<\frac{a+|\mu|}{\sigma})
\end{equation}
From \cite{chu}, for any $\delta>0$, $\frac{1}{2}\left(1-e^{-\delta^2/2}\right)^{\frac{1}{2}}\leq P(z \leq \delta)-\frac{1}{2}\leq \frac{1}{2}\left(1-e^{-2\delta^2/\pi}\right)^{\frac{1}{2}}$. Thus, for any $\gamma,\delta>0$:
$$ P(0<z \leq \gamma)-P(0<z \leq \delta)\leq
\frac{1}{2}\left(1-e^{-2\gamma^2/\pi}\right)^{\frac{1}{2}}-\frac{1}{2}\left(1-e^{-\delta^2/2}\right)^{\frac{1}{2}}$$
Applying the above to \eqref{eq:random}, we get:
$$P(|y_i|>a) \leq 
1+ \frac{1}{2}\left(\left(1-e^{-2(|\mu|-a)^2/\pi\sigma^2}\right)^{\frac{1}{2}}-\left(1-e^{-(a+|\mu|)^2/2\sigma^2}\right)^{\frac{1}{2}}\right)\\$$

Using that $x^{\frac{1}{2}}-y^{\frac{1}{2}}=\frac{x-y}{x^{\frac{1}{2}}+y^{\frac{1}{2}}}$ for $x,y>0$, we get:
$$P(|y_i|>a) \leq 1- \frac{e^{-2(|\mu|-a)^2/\pi\sigma^2}-e^{-(a+|\mu|)^2/2\sigma^2}}{2\left(\left(1-e^{-2(|\mu|-a)^2/\pi\sigma^2}\right)^{\frac{1}{2}}+\left(1-e^{-(a+|\mu|)^2/2\sigma^2}\right)^{\frac{1}{2}}\right)}$$

Finally, $P(\min_{i}|y_i|>a)=\exp\left(s\ln{P(|y_i|>a)}\right)$, using that $\ln\left(1+ x \right)< x$ for $x>-1$ we get the desired inequality.

\subsection{Proof of Lemma~\ref{lem:L0toL1convergencegen}}
Suppose $NC(M_1)+\ldots+NC(M_k)>\frac{k}{k+1}$ then for any $M\neq \{M_1,\ldots,M_k\}$, $NC(M)<\frac{1}{k+1}$. In addition, if $NC(M_1)+\ldots+NC(M_k)>\frac{k}{k+1}$ then necessarily $\max_{l=1,\ldots,k}NC(M_l)> \frac{1}{k+1}$ and then $\hat{S}^{b}\in \{M_1,\ldots,M_k\}$. Consequently, 
$$
P(\hat{S}^b\not\in \{M_1,\ldots,M_k\})\; \leq\; P\Big(NC(M_1)+\ldots+NC(M_k)\leq\frac{k}{k+1}\Big)$$
We have:
$$P\Big(NC(M_1)+\ldots+NC(M_k)\leq\frac{k}{k+1}\Big)\;=\;P\Big(\sum_{M \in \M\setminus \{M_1,\ldots,M_k\}} NC(M)\geq \frac{1}{k+1}\Big)$$
The result follows from the the Markov's inequality applied to the right-hand side above.

\subsection{Proof of Lemma~\ref{lem:boundrhogen}} 
Clearly $M \subseteq Q_T$ then $\mu_{Q_TM}$, as defined in \eqref{eq:muQM}, satisfies
\begin{eqnarray*}
    \mu_{Q_TM}&=&\|\big(I_n-P_M\big) \bs X_{Q_T\setminus M} \bs \beta^{\star}_{Q_T\setminus M}\|^2\\
    &=&\|\big(I_n-P_M\big) \bs X_{T\setminus M} \bs \beta^{\star}_{T\setminus M}\|^2\\
    &= & n {\beta^{\star}_{T\setminus M}}^\top \left(\tfrac1n \bs X_{T\setminus M}^\top(I_n-P_M)\bs X_{T\setminus M}\right)\beta^{\star}_{T\setminus M}\\
    &\geq & n\lambda_{\min}\left(\tfrac1n \bs X_{T\setminus M}^\top(I_n-P_M)\bs X_{T\setminus M}\right)\|\beta^{\star}_{T\setminus M}\|^2,
\end{eqnarray*}
where the first equality follows from observing that $Q_T\setminus M=T\setminus M$. Since $T$ is a susbet of $S$, by reordering columns, $X_{S\setminus M}=[X_{T\setminus M},X_{S\setminus (M\cup T)}]$ and $\tfrac1n \bs X_{T\setminus M}^\top(I_n-P_M)\bs X_{T\setminus M}$ is a diagonal block of $\tfrac1n \bs X_{S\setminus M}^\top(I_n-P_M)\bs X_{S\setminus M}$. By interlacing, $\lambda_{\min}\left(\tfrac1n \bs X_{T\setminus M}^\top(I_n-P_M)\bs X_{T\setminus M}\right) \geq \lambda_{\min}\left(\tfrac1n \bs X_{S\setminus M}^\top(I_n-P_M)\bs X_{S\setminus M}\right)$. Finally observing that 
$\lambda_{\min}\left(\tfrac1n \bs X_{T\setminus M}^\top(I_n-P_M)\bs X_{T\setminus M}\right) \geq \rho(\bs X)$ and $\|\beta^{\star}_{T\setminus M}\|^2\geq \sum_{j=1}^b|T_j\setminus M_j|\,\min_{i\in T_j\setminus M_j}{\beta_{i}^{\star}}^2$, gives the desired result.

\subsection{Proof of Lemma~\ref{lemma:s1s3}}
By Lemma 8.1 in \cite{birge}, we have for any $x>0$
$$P(W> (\nu+\mu)+2 \sqrt{(\nu+2 \mu) x}+2 x) \leq e^{-x}$$
The function $f:x\mapsto (\nu+\mu)+2 \sqrt{(\nu+2 \mu) x}+2 x$ is  one-to-one between $\R^+$ and $( \nu+
\mu,+\infty)$. We then have  for any $w>\mu+\nu$:
$$P(W > w) \leq e^{-f^{-1}(w)} \:=\; e^{-\big(\frac{w+\mu}{2}  - \sqrt{2w(\nu+2 \mu)-2 \mu \nu - \nu^2 }\big) } $$
Observe that:
$$\frac{w+\mu}{2}  - \sqrt{2w( 2 \mu+ \nu)-2 \mu \nu - \nu^2 }= \frac{w}{2}\bigg(1+\frac{\mu}{w} -\sqrt{\frac{8(2 \mu+\nu)}{w}}\bigg(1-\frac{\nu\mu}{2w\mu+w\nu}-\frac{\nu^2}{2(2w\mu+w\nu)} \bigg)\bigg)$$
 If $\nu=o(w)$ and $\mu=o(w)$, we then have $\frac{w+\mu}{2}  - \sqrt{2w( 2 \mu+ \nu)-2 \mu \nu - \nu^2 }= \frac{w}{2}(1+o(1))$ and for any $\phi\in(0,1)$ and $n$ large enough:
$$P(W > w) \leq  e^{-\phi\frac{w}{2}}$$

\subsection{Proofs of Lemma~\ref{lemma:s2}}
The firs part follows directly from Lemma S2 in \cite{rossell2022concentration}.

\subsection{Proof of Lemma~\ref{lemma:s18}}
For any $u\in (\underline{u},\bar{u})$, $2\ln \left(\frac{g}{1 / u-1}\right) \geq 2\ln \left(\frac{g}{1 / \underline{u}-1}\right)$. Since $\underline{u}\geq (1+g^{\phi}e^{-(\nu+\mu)/2})^{-1}$,  we have for any $u\in (\underline{u},\bar{u})$:
$$2\ln \left(\frac{g}{1 / u-1}\right) \geq 2(1-\phi)\ln(g)+\nu+\mu $$
and for any $u\in (\underline{u},\bar{u})$, $\nu=o\left(2\ln \left(\frac{g}{1 / u-1}\right)\right)$ and $\mu=o\left(2\ln \left(\frac{g}{1 / u-1}\right)\right)$. By Lemma~\ref{lemma:s1s3}, for any $\phi\in(0,1)$ and $n$ large enough:
\begin{equation}\label{eq:randomm}
    \int_{\underline{u}}^{\bar{u}} P\left(W>2\ln \left(\frac{g}{1 / u-1}\right)\right) d u< \frac{1}{g^{\phi}} \int_{\underline{u}}^{\bar{u}}(1 / u-1)^{\phi} d u .
\end{equation}

Applying the change of variables $v=1 / u-1$ to the integral above gives
$$
\int_{\underline{u}}^{\bar{u}}(1 / u-1)^{\phi} d u=\int_{1 / \bar{u}-1}^{1 / \underline{u}-1} \frac{v^{\phi}}{(v+1)^2} d v
$$
Note that $(1+x)^r \leq 1+r x$  for $x \geq-1$ and $r\in[0,1]$ and that $v^{\phi}=(v-1 + 1)^{\phi}$ with $v-1>-1$ since $v<1 / \bar{u}-1$ and $\bar{u}<1$. Then $v^{\phi}=(v-1 + 1)^{\phi}\leq 1 + \phi (v-1)\leq 1 + \phi(v+1)$. As a result:
\begin{equation}\label{eq:random2}
\int_{\underline{u}}^{\bar{u}}(1 / u-1)^{\phi} d u<\int_{1 / \bar{u}-1}^{1 / \underline{u}-1} \frac{1}{(v+1)^2} +\frac{\phi}{v+1} d v = \bar{u}-\underline{u}+\phi\ln\Big(\frac{\bar{u}}{\underline{u}}\Big)
\end{equation}
The results follows inputing the bound from \eqref{eq:random2} in \eqref{eq:randomm}.


\subsection{Proof of Lemma~\ref{lemma:s20}}
Since a probability is bounded by $1$, for any $a\in(0,1)$.
\begin{equation}\label{eq:basebound}
    \int_0^1 P\left(W>2 \ln \left(\frac{g}{1 / u-1}\right)\right) d u \;\leq\; 2 a+\int_{a}^{1-a} P\left(W>2\ln\left(\frac{g}{1 / u-1}\right)\right) d u.
\end{equation}

Take $a=\left(1+g^{\phi}e^{-(\nu+\mu)/ 2}\right)^{-1}$ for some $\phi\in(\alpha,1)$. By Lemma~\ref{lemma:s18} with $\underline{u}=a$ and $\bar{u}=1-a$:
$$
\int_0^1 P\left(W>2 \ln \left(\frac{g}{1 / u-1}\right)\right) d u \;\leq\; \frac{2}{1+g^{\phi}e^{-(\nu+\mu)/2}}+ 
\frac{1-2a+\ln(g^{\phi}e^{-(\nu+\mu)/2})}{g^{\phi}} $$
Since $1+g^{\phi}e^{-(\nu+\mu)/ 2}>g^{\phi}e^{-(\nu+\mu)/ 2}$ and $1-2a-\frac{\nu+\mu}{2}\leq \ln(g^\phi)$ for $n$ large enough, we have:
$$
\int_0^1 P\left(W>2 \ln \left(\frac{g}{1 / u-1}\right)\right) d u \;\leq\; g^{-\phi}\big(2e^{(\nu+\mu)/ 2}+2\ln(g^{\phi})\big) $$
We have:
\begin{equation}\label{eq:random24}
    \frac{g^{-\phi}2e^{(\nu+\mu)/ 2}}{g^{-\alpha}} \;=\;e^{-(\phi-\alpha)\ln(g)\big(1-\frac{\nu+\mu}{2(\phi-\alpha)\ln(g)}- \frac{\ln(2)}{(\phi-\alpha)\ln(g)}\big)}
\end{equation}
and 
\begin{equation}\label{eq:random25}
    \frac{g^{-\phi}2\ln(g^{\phi})}{g^{-\alpha}} \;=\;e^{-(\phi-\alpha)\ln(g)\big(1-\frac{\ln(\phi\ln(g))}{(\phi-\alpha)\ln(g)}- \frac{\ln(2)}{(\phi-\alpha)\ln(g)}\big)}
\end{equation}
Since $\alpha<\phi$, $g$ is increasing and $\nu=o(\ln(g))$ and $\mu=o(\ln(g))$, both expressions in \eqref{eq:random24} and \eqref{eq:random25} vanish as $n$ grows and for $n$ large enough:
$$
\int_0^1 P\left(W>2 \ln \left(\frac{g}{1 / u-1}\right)\right) d u \;\leq\; g^{-\alpha}$$

\subsection{Proof of Lemma~\ref{lem:boundexpnc}}
Since $M\in \mathcal M\setminus \{T\}$, by Lemma~\ref{lem:L0toL1convergence} (ii),  $NC(M)\leq (1+e^{C(T)-C(M)})^{-1}$.  By the standard formula for expectation of a nonnegative random variable in terms of its CDF, we get:
\begin{eqnarray*}
  \E\left(NC(M)\right)  &\leq& \int_0^1 P\left((1+e^{C(T)-C(M)})^{-1}\geq u\right){\rm d} u\\
  &=&\int_0^1 P\left(C(T)-C(M) \leq \ln{(\tfrac1u-1)}\right){\rm d}u\\
  &=&\int_0^1 P\Big(-\tfrac12 L_{TM} \geq \Delta_{MT}-\ln{(\tfrac1u-1)}\Big) {\rm d}u.
\end{eqnarray*}
Let $Q_T=M \cup T$ and recall that $L_{TM}=L_{Q_TM}-L_{Q_TT}$. For any $\gamma\in (\frac{1}{2},1]$
$$
-\tfrac12 L_{TM} -(\Delta_{MT}-\ln{(\tfrac1u-1)})\;=\;\left(\tfrac12 L_{Q_TT}-\ln{\Bigl(\frac{e^{\gamma\Delta_{MT}}}{\tfrac1u-1}\Bigl)}\right)-\Bigl(\tfrac12 L_{Q_TM}+(1-\gamma)\Delta_{MT}\Bigl).
$$
% \begin{equation}
%     L_{MT}-2\sum_{j=1}^b\Delta_{M,j}-2\ln{(1 / u-1)}= L_{Q_TM}-2\ln{\Bigl(\frac{e^{\gamma\sum_{j=1}^b\Delta_{M,j}}}{1 / u-1}\Bigl)}-\Bigl(L_{Q_TM}+2(1-\gamma)\sum_{j=1}^b\Delta_{M,j}\Bigl)
% \end{equation}
Observe that for any random variables $U$, $V$, and any $\epsilon,\gamma' \geq 0$, the event $
\{U-V \geq 0\}$ implies $\{U \geq \gamma' \epsilon\}\cup
\{V  < \gamma' \epsilon\}$. Let $U=\tfrac12 L_{Q_TT}-\ln{\Bigl(\frac{e^{\gamma\Delta_{MT}}}{\tfrac1u-1}\Bigl)}$ and $V=\tfrac12 L_{Q_TM}+(1-\gamma)\Delta_{MT}$. Take $\epsilon=\mu_{Q_TM}$ and $\gamma'=\tfrac12({1-\gamma})$. Recall that $A_T=\gamma\Delta_{MT}+\gamma'\mu_{Q_TM}$ and
$$
\{U\geq \gamma'\epsilon\}\;=\;\{\tfrac12 L_{Q_TT}\geq \ln{\Bigl(\frac{e^{\gamma\Delta_{MT}}}{\tfrac1u-1}}\Bigl)+\gamma'\mu_{Q_TM}\}\;=\;\{\tfrac12 L_{Q_TT}\geq \ln{\Bigl(\frac{e^{A_T}}{\tfrac1u-1}}\Bigl)\},
$$
$$
\{V<\gamma'\epsilon\}\;=\;\{\tfrac12 L_{Q_TM}<-(1-\gamma)\Delta_{MT}+\gamma'\mu_{Q_TM}\}\;=\;\{\tfrac12 L_{Q_TM}<\gamma'(\mu_{Q_TM}-2\Delta_{MT})\}.
$$
By the union bound we have:
\begin{equation}\label{eq:prooflemmas20nonspu}
    \E\left(NC(M)\right)  \leq \int_0^1P\Bigl(\tfrac12 L_{Q_TT}\geq \ln{\Bigl(\frac{e^{A_T}}{\tfrac1u-1}\Bigl)}\Bigl){\rm d}u+ P\Bigl(\tfrac12 L_{Q_TM}<\gamma'(\mu_{Q_TM}-2\Delta_{MT})\Bigl).
\end{equation}
We first bound the first term in \eqref{eq:prooflemmas20nonspu}. If $M \subset T$ then $Q_T=T$, $L_{Q_TT}=0$, and this term is zero. Suppose now that $T$ is a strict subset of $Q_T$. By Lemma~\ref{lemma:s7}, $L_{Q_TT}$ has a $\chi^2$-distribution. More precisely, $L_{Q_TT} \sim \mathcal{\chi}_{|Q_T\setminus T|}^2(\mu_{Q_TT})$ {but $\mu_{Q_TT}=0$ because $\beta^\star_{Q_T\setminus T}=0$}. Let $\nu=|M\setminus T|$. By definition $\mu_{Q_TM}$ is nonnegative and $A_T \geq \gamma \Delta_{MT}$. Since by assumption $|M\setminus T|=o(\ln(e^{A_T}))$ and $\mu_{Q_TT}=o(\ln(e^{A_T}))$. Take $\alpha\in(\psi,1)$, by Lemma~\ref{lemma:s20}, t and $n$ large enough:
\begin{equation}\label{eq:boundLQS}
\int_0^1P\Bigl(L_{Q_TT}>2\ln{\Bigl(\frac{e^{A_T}}{1 / u-1}\Bigl)}\Bigl){\rm d}u \;<\; 
e^{-\alpha A_T}.
\end{equation}
% or, equivalently, 
% \begin{equation}\label{eq:boundLQT}
% \int_0^1P\Bigl(L_{Q_TT}>2\ln{\Bigl(\frac{e^{\gamma\Delta_{M}+\gamma'\mu_{Q_TM}}}{1 / u-1}\Bigl)}\Bigl){\rm d}u \;<\; 16
% e^{-\gamma\Delta_{M}+\gamma'\mu_{Q_TM}}\left(\gamma\Delta_{M}+\gamma'\mu_{Q_TM}\right)^{\frac{|M\setminus T|}{2}+1}.
% \end{equation}
We now bound the second term in \eqref{eq:prooflemmas20nonspu}. If $T \subset M$, then $Q_T=M$, $L_{Q_TM}=0$, $\mu_{Q_TM}=0$, and this term is zero. If $M$ is a strict subset of $Q_T$ then, by Lemma~\ref{lemma:s7}, $L_{Q_TM} \sim \chi_{|Q_T\setminus M|}^2(\mu_{Q_TM})$ %(note that $Q_T\setminus M=T\setminus M$ and $\mu_{Q_TM}=\mu_{Q_TM}$). 
Clearly, when $\mu_{Q_TM}\leq 2\Delta_{MT}$, this term is also zero so suppose  $\mu_{Q_TM}> 2\Delta_{MT}$.  
We have  
$$P\left(L_{Q_TM}<2\gamma'(\mu_{Q_TM}-2\Delta_{MT})\right)\; \leq \; P\left(L_{Q_TM}<2\gamma'\mu_{Q_TM}\right)$$ 
If $M$ is a strict subset of $Q_T$,   $Q_T\setminus M=T\setminus M\neq \empty$ and by assumption $|T\setminus M|=o(\mu_{Q_TM})$, then by Lemma~\ref{lemma:s2}, for any $\phi\in(0,1)$, for $n$ large enough:
\begin{equation}\label{eq:aaaaa}
    P\left(L_{Q_TM}<2\gamma'\mu_{Q_TM}\right) \;\;\leq\;\; e^{-\phi\frac{\mu_{Q_TM}}{2}}
\end{equation}
Moreover we have:
$$
A_T\;=\;\frac{1-\gamma}{2}\mu_{Q_TM}+\gamma\Delta_{MT}\;=\;\Big(1-\gamma+\gamma\frac{2\Delta_{MT}}{\mu_{Q_TM}}\Big)\frac{\mu_{Q_TM}}{2}
$$
Since $\mu_{Q_TM}>2\Delta_{MT}$, $\big(1-\gamma+\gamma\frac{2\Delta_{MT}}{\mu_{Q_TM}}\big)<1$, by \eqref{eq:aaaaa}, for $n$ large enough:
\begin{equation}\label{eq:boundLQM}
    P\left(L_{Q_TM}<2\gamma'(\mu_{Q_TM}-2\Delta_{MT})\right) \;\leq\; e^{-A_T} \;\leq\; e^{-\alpha A_T}
\end{equation}
Summing the bounds in \eqref{eq:boundLQS} and \eqref{eq:boundLQM} gives that for $n$ large enough $\E\left(NC(M)\right) < 2e^{-\alpha A_T}$. Since $\psi<\alpha$, for $n$ large enough $\E\left(NC(M)\right) < e^{-\psi A_T}$


\subsection{Proof of Lemma~\ref{lem:upperboundnoncentralgen}}
By Lemma~\ref{lemma:s7}, we have
\begin{align*}
    \mu_{Q_TT}&={\bs \beta^{\star}_{Q_T\setminus T}}^{\top}\bs X_{Q_T\setminus T}^{\top}  \big(I_n-P_T\big) \bs X_{Q_T\setminus T} \bs \beta^{\star}_{Q_T\setminus T}\\
    &={\bs \beta^{\star}_{M\setminus T}}^{\top}\bs X_{M\setminus T}^{\top}  \big(I_n-P_T\big) \bs X_{M\setminus T} \bs \beta^{\star}_{M\setminus T}\\
    &={\bs \beta^{\star}_{(S\cap M)\setminus T}}^{\top}\bs X_{(S\cap M)\setminus T}^{\top}  \big(I_n-P_T\big) \bs X_{(S\cap M)\setminus T} \bs \beta^{\star}_{(S\cap M)\setminus T}
\end{align*}
where the third equality follows from $\bs \beta^*_{M\setminus S}=0$.
We start with the upper bound. Denote for any square matrix $A$, its largest eigenvalue 
$\lambda_{\max}(A)$. We have:
$$\mu_{Q_TT}\leq n\lambda_{\max}\Big(\frac{1}{n}\bs X_{(S\cap M)\setminus T}^{\top}  \big(I_n-P_T\big) \bs X_{(S\cap M)\setminus T}\Big) \|\bs \beta^{\star}_{(S\cap M)\setminus T}\|_2^2$$
Let $B=\tfrac{1}{n}\bs X_{(S\cap M)\setminus T}^{\top}  \big(I_n-P_T\big) \bs X_{(S\cap M)\setminus T}$ and $C=\tfrac{1}{n}\bs X_{(S\cap M)\cup T}^{\top}\bs X_{(S\cap M)\cup T}$. $B$ is the Schur complement $\big(C \,  / \, \tfrac{1}{n}\bs X_T^{\top} \bs X_T\big)$. The inverse $B^{-1}$ is then a diagonal block of $C^{-1}$. By interlacing. $\lambda_{min}(B^{-1})\geq \lambda_{min}(C^{-1})$ and then $\lambda_{\max}(B)\leq \lambda_{\max}(C)$. We also have $(S\cap M)\cup T \subseteq S$, then by interlacing again $\lambda_{\max}(C)\leq \lambda_{\max}\Big(\tfrac{1}{n}\bs X_S^{\top}\bs X_S\Big)=\bar{\lambda}$. The upper bound follows from also observing that $\|\bs \beta^{\star}_{(S\cap M)\setminus T}\|_2^2\leq \sum_{j=1}^b |(S_j\cap M_j)\setminus T_j|\max_{i\in (S_j\cap M_j)\setminus T_j}{\beta_i^*}^2$.

A lower bound on $\mu_{Q_TT}$ is: 
$$\mu_{Q_TT}\geq n\lambda_{\min}\Big(\frac{1}{n}\bs X_{(S\cap M)\setminus T}^{\top}  \big(I_n-P_T\big) \bs X_{(S\cap M)\setminus T}\Big) \|\bs \beta^{\star}_{(S\cap M)\setminus T}\|_2^2$$
The final bound is obtained by similar interlacing argument as for the upper bound and the bound $\|\bs \beta^{\star}_{(S\cap M)\setminus T}\|_2^2\geq \sum_{j=1}^b |(S_j\cap M_j)\setminus T_j|\min_{i\in (S_j\cap M_j)\setminus T_j}{\beta_i^*}^2$

\section{Properties in the Gaussian sequence model with $s$ finite}\label{sec:appendixGSMsfinite}

We derive here selection properties of the block $\ell_0$ penalties in the Gaussian Sequence model dropping assumption~A3 and assuming instead:
\begin{enumerate}[start=8,label={(A\arabic*)}]
    \item $\text{for all } j, \;\; s_j \,< \, \infty$
\end{enumerate}

Changing A3 for A8 implies redeveloping results relative to the probability of false negatives and consequently betamin conditions. Observe that the Lemmas and Properties of Section~\ref{sec:GSMgen} as well as Property \ref{prop:suffnecIblockthresh} in Section~\ref{sec:properties_legression} do not depend on A3, the probability of false negatives or a betamin condition. They hold equally in the regime where $s_j$ is finite for all $j=1,\ldots,b$.

Consider the condition:
\begin{enumerate}[start=9,label={(A\arabic*)}]
    \item $\text{for all } j, \;\; \sqrt{n}(\beta_{\min,j}^{\star}-\tau_j) \to \infty$
\end{enumerate}
Under A8, A9 can be shown to be sufficient for $\hat{S}^{T,b}$ to hold the screening property. Moreover in Proposition~\ref{prop:necIIblockthresh} a) we showed A9 was a general necessary condition for the screening property without assuming A3. A9 is then necessary and sufficient for $\hat{S}^{T,b}$ to hold the property under A8.
\begin{prop}\label{prop:suffnecIIblockthreshbis}
 In the Gaussian sequence model \eqref{eq:GSM}, under assumptions A1, A2 and A9, $\hat{S}^{T,b}$ satisfies:
$$P(S\subseteq \hat{S}^{T,b} ) \to 1$$
\end{prop}
\begin{proof}
    By the proof of Proposition~\ref{prop:suffnecIblockthresh}, we have: \begin{equation*}
P\left(\hat{S}^{T,b} \not\subseteq S\right)\;\leq\; \sum_{j=1}^b\frac{e^{-\tfrac{n}{2}\left((\beta^\star_{\min,j}-\tau_j)^2-\tfrac{2\ln(s_j)}{n}\right)^2}}{\sqrt{\pi\ln(s_j)}}
 \end{equation*}
 where the right-hand side goes to 0 by assumption A9.
\end{proof}
By Proposition~\ref{prop:suffnecIblockthresh}, to guarantee model recovery with high probability the block thresholds $\tau_j$ must be greater or equal to $\sqrt{2\ln(p_j-s_j)/n}$. Under A8, we then get the minimal betamin condition, for all $j=1,\ldots,b$:
\begin{equation}\label{eq:nectypeIIGSMfinite}
    \sqrt{n}\beta_{\min,j}^{\star}-\sqrt{2\ln(p_j-s_j)} \to \infty
\end{equation}
Condition \eqref{eq:nectypeIIGSMfinite} is necessary for asymptotic support recovery with $\hat{S}^{T,b}$:
\begin{lem}\label{lem:necbetaminblockthresh2}
    Suppose there exists $j\in\{1,\ldots,b\}$ such that
    \begin{equation}\label{eq:necbetaminblockthreshbis}
      \lim_{n\to \infty}\sqrt{n}\beta_{\min,j}^*-\sqrt{2\ln(p_j-s_j)} < +\infty  
    \end{equation}
then under assumptions A1, A2 and A8,\;  $P(\hat{S}^{T,b} = S)\not \to 1$
\end{lem}

\begin{proof}

  We have $\sqrt{n}\beta_{\min,j}^*-\sqrt{2\ln(p_j-s_j)} \,=\, \sqrt{n}\beta_{\min,j}^*-\sqrt{n}\tau_j+\sqrt{2\ln(p_j-s_j)}\Big(\tfrac{\sqrt{n}\tau_j}{\sqrt{2\ln(p_j-s_j)}}-1\Big)$ and \eqref{eq:necbetaminblockthreshbis} implies that there exists $c\in\R^+$ such that:
  \begin{equation}\label{eq:proofnecbetaminbis}
      \lim_{n\to \infty}\sqrt{n}\beta_{\min,j}^*-\sqrt{n}\tau_j+\sqrt{2\ln(p_j-s_j)}\Big(\tfrac{\sqrt{n}\tau_j}{\sqrt{2\ln(p_j-s_j)}}-1\Big)\leq c
  \end{equation}

Consider the case $\lim_{n\to +\infty}\tfrac{\sqrt{n}\tau_j}{\sqrt{2\ln(p_j-s_j)}}-1<0$. Then by Proposition~\ref{prop:suffnecIblockthresh},
$P(\hat{S}^{T,b} = S)\not \to 1$. 
Now consider the case $\lim_{n\to +\infty}\tfrac{\sqrt{n}\tau_j}{\sqrt{2\ln(p_j-s_j)}}-1 \geq0$. Then by \eqref{eq:proofnecbetaminbis},  $\lim_{n\to \infty}\sqrt{n}\beta_{\min,j}^*-\sqrt{n}\tau_j \leq c$ and by Proposition~\ref{prop:necIIblockthresh} a), $P(\hat{S}^{T,b} = S)\not \to 1$.
\end{proof}

Theorem~\ref{theo:fullrecovGS} does not assume A3 but requires the earlier betamin assumption~A5. Now under A8, A9 implies A5. As a result, when $s_j$ is finite for all $j=1,\ldots,b$, Theorem~\ref{theo:fullrecovGS} holds equally assuming A8 and A9 instead of A5.

It follows that the gains in terms of convergence rate with  informed block thresholding under A8 are the same as under A3. To discuss gains in terms of betamin conditions, we assume without loss of generality that $\beta_{\min}^{\star}$ is located in block $b$. By \eqref{eq:nectypeIIGSMfinite}, the minimal betamin conditions with block informed thresholding $\hat{S}^{T,b}$ and standard thresholding $\hat{S}^1$ are respectively:
\begin{align}
    \sqrt{n}\beta_{\min,b}^{\star}-\sqrt{2\ln(p_b-s_b)} \to \infty \quad\text{and}\quad
    \sqrt{n}\beta_{\min}^{\star}-\sqrt{2\ln(p-s)} \to \infty
\end{align}
The relevant quantities to compare the minimal requirement on the growth rate of betamin are thus $\sqrt{\tfrac{\ln(p_b-s_b)}{n}}$ and $\sqrt{\tfrac{\ln(p-s)}{n}}$. Since $p_b-s_b \leq p - s$, the minimal betamin condition with $\hat{S}^{T,b}$ is always milder than with $\hat{S}^{T,1}$ under A8 again. Moreover, observe that under Scenarios 1 to 3 in Section~\ref{sec:examples}, the relevant ratio to compare minimal betamin condition was also $\Theta\Big(\sqrt{\tfrac{\ln(p_b-s_b)}{\ln(p-s)}}\Big)$. As a result, those scenarios also illustrate the magnitude of the gains in terms of minimal betamin conditions under A8.


\section{Non linear block $\ell_0$ penalty in high-dimensional linear regression}\label{sec:appendixnonlinear}
\sdescription{Short description of Supplement C.}

In this section we derive selection properties of the block $\ell_0$ penalties in the linear regression setting without assuming linearity. 

Let $\eta_j$ for $j=1,\ldots,b$ be any non-negative and increasing penalty functions on the natural numbers. The selector based on those $\ell_0$ block penalties is:
\begin{equation}
  \widehat S^b\;\in\;\arg\max_{M\in \M} \left\{\max_{\beta\in \mathcal{L}_M} \ell(\bs y;\bs \beta)-\sum_{j=1}^b\eta_j (|M_j|)\right\},\qquad M=M_1\cup \ldots\cup M_b.	
\end{equation}
where $\ell$ is the log-likelihood function.

Rewrite the difference in penalty between any model $M$ and the true model $S$:
\begin{equation}\label{eq:DeltaMbis}
\Delta_{MS} \; := \; \sum_{j=1}^b \eta_j(|M_j|)-\eta_j(|S_j|)
\end{equation}
%The proof of Lemma~\ref{lem:boundexpnc} does not rely on the linearity of block penalties and the Lemma holds identically with the redefined $\Delta_{M}$. 
We define the normalized difference in block penalty between any model $M$ and the true support $S$ as:
 $$\tilde{\kappa}_{M, j}\; := \; \begin{cases}
     \frac{\eta_j(|M_j|)-\eta_j(|S_j|)}{|M_j|\,-\,|S_j|} \;\;\; \text{if}\;\;\; |M_j|\,\neq\,|S_j| \\
     0 \;\;\;\text{if}\;\;\; |M_j|\,=\,|S_j|
 \end{cases}  $$
The quantity $\tilde{\kappa}_{M, j}$ plays a similar role as $\kappa_j$, representing the average penalty incurred for adding a variable from block $B_j$ to the model $M$. However, unlike in the case of linear penalties, $\tilde{\kappa}_{M, j}$ is not constant but a function of $|M_j|$ and $|S_j|$. For any $M\in\M$, $\tilde{\kappa}_{M, j}$ is exactly $\kappa_j$ if the penalty is linear as defined in Section \ref{sec:intro_blockpen}.

As in the linear case, consistency is achieved by setting conditions on the $\tilde{\kappa}_{M, j}$ and $\beta^*_{\min, j}$:
\begin{enumerate}[start=9,label={(A\arabic*)}]
    \item for each $j$, there exist $b_{j}\to \infty$ (as $n\to \infty$) such that for all $M \in \M$ such that $|M_j| \neq |S_j|$ and $|M_j \setminus S_j| >0$, for $n$ sufficiently large, and  all $j=1,\ldots,b$:
\begin{equation}\label{cond:suffIregnonlinpen}
    \tilde{\kappa}_{M, j} =\;\ln\bigg(\frac{p_j-s_j}{|M_j \setminus S_j| }\bigg) + b_{j}
\end{equation}

\item  for each $j$, there exist $c_{j}\to +\infty$ such that for $n$ sufficiently large, for all $j=1,\ldots,b$
    \begin{equation}\label{cond:suffIIregnonlinpen}
        \sqrt{\frac{(1-\gamma) n\rho(\bs X)}{2}}{{\beta_{\min}^{\star}}_j} \,- \,\sqrt{\ln(p_j-s_j) + b_{j}} \;=\; \sqrt{\ln(s_j)} + c_{j} .
    \end{equation}
where $\alpha_j:=\frac{\ln(p_j-s_j)}{\ln(p_j-s_j) + b_{j}}$ and $\gamma:=\frac12(1+\max_j \alpha_j)\in (\tfrac12,1)$ for $n$ large enough.
\end{enumerate}
\begin{thm}\label{theo:suffcondlinearmodelnonlinpen}
Suppose that A1, A2, A7 and A10 hold, then 
    % \begin{align}
    %  &\qquad \tilde{\kappa}_{M, j} > \ln(p_j-s_j)  \label{cond:suffIregnonlinpen}\\
    %    &\qquad \sqrt{\frac{2(1-\gamma) n\rho(\bs X)}{8}}{{\beta_{\min}^{\star}}_j}  \geq \sqrt{\tilde{\kappa}_{M, j}}+{\sqrt{\ln(s_j)}}  .\label{cond:suffIIregnonlinpen}
    % \end{align}
Then $$\sum_{M \in \M\setminus \{S\}}\E\left(NC(M)\right)\to 0 \qquad \mbox{and} \qquad P(\widehat{S}^b = S) \to 1$$
\end{thm}

\begin{proof}
The proof is essentially the same as proof Theorem~\ref{theo:suffcondlinearmodel} replacing $\kappa_j$ by $\tilde{\kappa}_{M, j}$ for every $j=1,\ldots,b$. We rewrite here the main intermediate results of the proof Theorem~\ref{theo:suffcondlinearmodel} and redevelop only the steps altered by the non-linearity of the penalty. We refer the reader to the proof of Theorem~\ref{theo:suffcondlinearmodel} for more details.

Observe that $\Delta_{MS}$, defined in \eqref{eq:DeltaMbis}, can be rewritten as $\Delta_{MS}=\sum_{j=1}^b(|M_j\setminus S_j|-|S_j\setminus M_j|)\tilde{\kappa}_{M}$. By  Lemma~\ref{lem:boundrho}, for every $n\in \N$ we have 
\begin{equation}\label{eq:nonoverfittingcondgen}
   A_S\;=\; \gamma \Delta_{MS}+\frac{1-\gamma}{2}\mu_{Q_SM}\;\geq\; \gamma\sum_{j=1}^b|M_j \setminus S_j|\tilde{\kappa}_{M} + \sum_{j=1}^b |S_j\setminus M_j |\left(\tfrac{1-\gamma}{2} n \rho(\bs X) {\beta^{\star}_{\min,j}}^{2} - \gamma\tilde{\kappa}_{M, j}\right) 
\end{equation} 
As for Theorem~\ref{theo:suffcondlinearmodel}, by assumption A10 we have $\gamma\sum_{j}|M_j\setminus S_j|\tilde{\kappa}_{M, j}\geq \tfrac{3}{4}|M\setminus S|$ as long as $\min_j \ln(p_j-s_j)\geq \tfrac32$, which always holds in our asymptotic regime for $n$ sufficiently large. Observe that the proof of Lemma~\ref{lem:boundexpnc} does not rely on the linearity of block penalties and the Lemma holds identically with the redefined $\Delta_{MS}$.  By Lemma~\ref{lem:boundexpnc} we have:
$$
\E(NC(M))\;\leq\;16 e^{-A_S}A_S^{|M\setminus S|/2+1}+e^{-A_S},
$$
where $A_S=\gamma\Delta_{MS}+\tfrac{1-\gamma}{2}\mu_{Q_SM}$. Similarly to Theorem~\ref{theo:suffcondlinearmodel}, defining:
$$A_S^*\;:=\;\gamma\sum_{j=1}^b|M_j \setminus S_j|\tilde{\kappa}_{M, j} + \sum_{j=1}^b |S_j\setminus M_j |\left(\tfrac{1-\gamma}{2} n \rho(\bs X) {\beta^{\star}_{\min,j}}^{2} - \gamma\tilde{\kappa}_{M, j}\right) $$
we get 
$$
\E(NC(M))\;\leq\;16 e^{-A_S^*}{A_S^*}^{|M\setminus S|/2+1}+e^{-A_S^*}.
$$
Note that 
$$A_S^*\;=\;\gamma\sum_{j \in \{1,\ldots, b\} \text{s.t.} |M_j \setminus S_j| \neq 0}^b|M_j \setminus S_j|\tilde{\kappa}_{M, j} + \sum_{j=1}^b |S_j\setminus M_j |\left(\tfrac{1-\gamma}{2} n \rho(\bs X) {\beta^{\star}_{\min,j}}^{2} - \gamma\tilde{\kappa}_{M, j}\right) $$
Furthermore, condition \eqref{cond:suffIIregnonlinpen} implies there exist $c^{\prime}_j \to +\infty$ such that:
$$\frac{(1-\gamma) n\rho(\bs X)}{2}{{\beta_{\min}^{\star}}_j}^2 \,- \,\tilde{\kappa}_{M, j} \;=\; \ln(s_j) + c^{\prime}_{j}.$$
We also have that for $x$ sufficiently large and any $\psi \in (0,1)$
\begin{equation}\label{eq:boundexp3}
16e^{-x}x^{\tfrac12|M\setminus S|+1}+e^{-x}\leq e^{-\psi x}.    
\end{equation}
To choose $\psi$ we proceed as in the proof of Theorem~\ref{theo:suffcondlinearmodel}. Let $\delta \in (0,1)$ and denote:
$$\bar{m}_j=\begin{cases}
    \max\bigg\{\tfrac{2\ln\big(\tfrac{p_j-s_j}{|M_j\setminus S_j|}\big)}{b_{j}},\tfrac{2\ln(s_j)}{c^{\prime}_{j}}\bigg\},\;\;\; \text{if}\;\;|M_j \setminus S_j| \neq 0 \\
    \frac{2\ln(s_j)}{c^{\prime}_{j}},\;\;\; \text{if}\;\;|M_j \setminus S_j| \neq 0
\end{cases}$$
Take $\psi=\max_{j=1, \ldots,b}\frac{a + \delta + \bar{m}}{1+\bar{m}}$ for some $a \in(0, 1- \delta)$ then $\psi \in (0,1)$ and for every $j=1,\ldots,b$:
$$\delta \leq \psi-(1-\psi)\bar{m}_j.$$ 
By definition of $\gamma$ and $\alpha_j$, we have for every $j=1,\ldots,b$ such that $|M_j \setminus S_j|\neq0$:
$$
\gamma\tilde{\kappa}_{M, j}\;\geq\;\tfrac12 (1+\alpha_j)\tilde{\kappa}_{M, j}\geq \tfrac12 \bigg(1+\tfrac{\ln\big(\tfrac{p_j-s_j}{|M_j\setminus S_j|}\big)}{\tilde{\kappa}_{M, j}}\bigg)\tilde{\kappa}_{M, j} \;=\; \ln\big(\tfrac{p_j-s_j}{|M_j\setminus S_j|}\big)+ \tfrac12b_{j}.
$$
From our choice of $\psi,\delta$, we further get:
$$\psi\gamma\tilde{\kappa}_{M, j} \;\geq\; \psi\bigg(\ln\big(\tfrac{p_j-s_j}{|M_j\setminus S_j|}\big)+ \frac12b_{j}\bigg) \;\geq\; \ln\big(\tfrac{p_j-s_j}{|M_j\setminus S_j|}\big)+ \delta\frac12b_{j}.$$
Observe that, for any $M$, $\tilde{\kappa}_{M, j} \leq \ln(p_j-s_j)+b_j$. Then, by \eqref{cond:suffIIregnonlinpen} and our choice of $\psi,\delta$:
$$\psi\left(\tfrac{1-\gamma}{2} n \rho(\bs X) {\beta^{\star}_{\min,j}}^{2} - \gamma\tilde{\kappa}_{M, j}\right) \;\geq\; \psi\left(\ln(s_j)+c_j'\right) \;\geq\; \ln(s_j)+\delta \frac{c^{\prime}_j}{2}.
$$
Using the expression for $A_S^*$, the bound in \eqref{eq:boundexp3}, and the above inequalities, we get
$$
\E(NC(M))\;\leq\;\exp\{- \sum_{j \in \{1,\ldots, b\} \text{s.t.} |M_j \setminus S_j| > 0} |M_j\setminus S_j|\bigg(\ln\bigg(\frac{p_j-s_j}{|M_j\setminus S_j|}\bigg)+\delta \tfrac{b_j}{2}\bigg)-\sum_{j=1}^b |S_j\setminus M_j|(\ln(s_j)+\delta \tfrac{c^{\prime}_j}{2})\}.
$$
Following the same subsequent steps of the proof of Theorem~\ref{theo:suffcondlinearmodel} shows selection consistency.
\end{proof}

The EBIC penalty is equivalent to considering that all the variables belong to a single block, that is $b=1$, and setting $\eta_1(|M|)=\psi\ln{p\choose |M|} + \tfrac{1}{2}\ln(n)$. The EBIC penalty can be shown to satisfy assumption~A10 whenever $s=o(\ln(n))$ and the following corollary holds:
\begin{cor}
    Suppose that $b=1$, $\eta_1(|M|)=\psi \ln {p\choose |M|} + \tfrac{1}{2}\ln(n)$ with $\psi\geq1$, $s=o(\sqrt{n})$, $\max|M|_{M\in\M} \leq \tfrac{p}{2}$ and that there exists $c\to +\infty$ such that for $n$ sufficiently large
    \begin{equation}\label{eq:betaminebic}
        \sqrt{\frac{\ln(n)-2(s-2)}{2 \ln(p)+\ln(n)-2(s-2)}\frac{n\rho(\bs X)}{16}}{{\beta_{\min}^{\star}}} \,- \,\sqrt{\psi\ln(p) + \frac{1}{2}\ln(n)} \;=\; \sqrt{\ln(s)} + c .
    \end{equation}
    Then A10 and A11 hold and under A1 and A2: $$\sum_{M \in \M\setminus \{S\}}\E\left(NC(M)\right)\to 0 \qquad \mbox{and} \qquad P(\widehat{S}^b = S) \to 1$$    
\end{cor}

\begin{proof}
    We start by showing $\eta_1$ satisfies assumption~A10. Let $M\in\M$ and denote $|M\setminus S|=u$ and $|S\setminus M|=w$. We have $|M|=u+s-w$ and $|M|-|S|=u-w$. As in A10 consider the case $u-w>0$ and $u>0$, then:
    \begin{equation*}
        \tilde{\kappa}_{M,1}=\frac{\psi}{u-w}\ln\bigg[\frac{{p \choose |M|}}{{p \choose |S|}}\bigg]+ \frac{1}{2}\ln n
    \end{equation*}
    A well-known property of binomial coefficients is that for any positive integers $n$, $h$, $k$ we have:
    \begin{equation}\label{eq:idbinom}
        \binom{n}{h}\binom{n-h}{k}=\binom{n}{h+k}\binom{h+k}{h}
    \end{equation}
    Consider first the case in which $u-w>0$, take $n=p$, $k=u-w$ and $h=s$ in \eqref{eq:idbinom}, we have:
    \begin{equation}\label{eq:identitybinom}
        \frac{{p \choose |M|}}{{p \choose s}}=
        %\frac{{n \choose h+k}}{{n \choose h}}= \frac{{n -h\choose k}}{{h+k \choose h}} =
        \frac{{p -s \choose u-w}}{{|M| \choose s}}
    \end{equation}
    % \begin{equation}\label{eq:identitybinom}
    %     \frac{{p \choose |M|}}{{p \choose s}}=\frac{{p \choose p-s+w-u}}{{p \choose s}} =
    %     \frac{{n \choose h+k}}{{n \choose n-h}}=\frac{{n \choose h+k}}{{n \choose h}}= \frac{{n -h\choose k}}{{h+k \choose h}} =
    %     \frac{{p -s \choose u-w}}{{|M| \choose s}}
    % \end{equation}
A standard bounds on binomial coefficient for $1\leq k\leq n$ are:
\begin{equation}\label{eq:boundbinom}
  \left(\frac{n}{k}\right)^k \leq  {n \choose k} \leq \frac{n^n}{k^k(n-k)^{n-k}}
\end{equation}
Combining \eqref{eq:identitybinom} and \eqref{eq:boundbinom}, we get:
$$ \frac{\psi}{u-w}\ln\bigg[\frac{{p \choose |M|}}{{p \choose s}}\bigg] \geq \psi\ln\bigg(\frac{p-s}{u-w}\bigg)-\psi\bigg[\frac{s}{u-w}\ln\bigg(\frac{u-w}{s}+1\bigg) +\ln\bigg(1+\frac{s}{u-w}\bigg) \bigg]$$
We have:
$$\psi\ln\bigg(\frac{p-s}{u-w}\bigg) \geq \psi\ln\bigg(\frac{p-s}{u}\bigg) \geq \ln\bigg(\frac{p-s}{u}\bigg)$$
Observe also that $f:x \mapsto \frac{\ln(x+1)}{x}+\ln(1+{x}^{-1})$ for $x>0$ is decreasing. We have $\frac{u-w}{s} \geq s^{-1}$ and then: 
\begin{equation}\label{eq:lowerboundb1ebic2}
        \tilde{\kappa}_{M,1} \geq \ln\big(\frac{p-s}{u}\big)-s\ln(s^{-1}+1)+ \ln \big(\tfrac{\sqrt{n}}{1+s}\big)
\end{equation}
 We have that $\lim_{n\to+\infty}s\ln(s^{-1}+1)< +\infty$ and since $s=o(\sqrt{n})$, $\ln(\tfrac{\sqrt{n}}{s+1}) \to +\infty$, then $\eta_1$ satisfies assumption~A10 in the case $u-w>0$.\\
 Consider now the case where $u-w<0$. We have:
     \begin{equation*}
        \tilde{\kappa}_{M,1}=\frac{\psi}{w-u}\ln\bigg[\frac{{p \choose s}}{{p \choose |M|}}\bigg]+ \frac{1}{2}\ln n
    \end{equation*}
 Taking $n=p$, $k=w-u$ and $h=|M|$ in \eqref{eq:idbinom}, we obtain:
     \begin{equation}
        \frac{{p \choose s}}{{p \choose |M|}}=
        %\frac{{n \choose h+k}}{{n \choose h}}= \frac{{n -h\choose k}}{{h+k \choose h}} =
        \frac{{p - |M| \choose w-u}}{{s \choose |M|}}
    \end{equation}
By bounds in \eqref{eq:boundbinom}, we get:
$$ \frac{\psi}{w-u}\ln\bigg[\frac{{p \choose s}}{{p \choose |M|}}\bigg] \geq \psi\ln\bigg(\frac{p-|M|}{w-u}\bigg)+\psi\bigg[\frac{s}{w-u}\ln\bigg(1-\frac{w-u}{s}\bigg) -\ln\bigg(\frac{s}{w-u}-1\bigg) \bigg]$$
Since $|M|<s\leq p/2$ and $\psi\geq 1$ we have:
$$\psi\ln\bigg(\frac{p-|M|}{w-u}\bigg) \geq \psi\ln\bigg(\frac{p-s}{w-u}\bigg) \geq \psi\ln\bigg(\frac{p-s}{w}\bigg)\geq \ln\bigg(\frac{p-s}{u}\bigg)$$
Observe that $g:x \mapsto \frac{\ln(1-x)}{x}-\ln({x}^{-1}-1)$ for $x\in(0,1)$ is increasing. We have $\frac{w-u}{s} \geq s^{-1}$ and then: 
\begin{equation}
        \tilde{\kappa}_{M,1} \geq \ln\big(\frac{p-s}{u}\big)-s\ln(1-s^{-1})+ \ln \big(\tfrac{\sqrt{n}}{s-1}\big)
\end{equation}
Note that the righ-handside above is well defined. Since $u>0$ and $u-w<0$, we have $w\geq2$ and $|M|<s$. The later is not possible unless $s-1>0$. We have that $\lim_{n\to+\infty}s\ln(1-s^{-1})< +\infty$ and since $s=o(\sqrt{n})$, $\ln(\tfrac{\sqrt{n}}{s-1}) \to +\infty$, then $\eta_1$ satisfies assumption~A10 in the case $u-w<0$ too.

 We now show that if \eqref{eq:betaminebic} holds then A11 holds. Since A10 holds there exists $b_1$ such that for all $M\in\M$ such that $|M|\neq s$ and $|M\setminus S|>0$, $\tilde{\kappa}_{M, 1} =\;\ln\bigg(\frac{p-s}{|M \setminus S| }\bigg) + b_{1}$. By inequalities \eqref{eq:lowerboundb1ebic1} and \eqref{eq:lowerboundb1ebic2}, for any of those $M$, $b_1 \geq \tfrac{1}{2}\ln(n)-(s+2)$. By assumption~A11 we have: $1-\gamma=1-\frac{1}{2}\frac{\ln(p-s)}{})$ XXX -MISSING XXX
\end{proof}


\end{document}
